\chapt{श्रावणमासमाहात्म्यम्}

\sect{प्रथमोऽध्यायः}

\uvacha{शौनक उवाच}

\twolineshloka
{सूत सूत महाभाग व्यासशिष्य ह्यकल्मष}
{त्वदीयवदनाम्भोजान्नानाख्यानानि शृण्वताम्} %॥१॥

\twolineshloka
{तृप्तिर्न जायते भूयः श्रवणेच्छा प्रवर्धते}
{कार्तिकस्य च माहात्म्यं तुलासंस्थे दिवाकरे} %॥२॥

\twolineshloka
{माघमासस्य माहात्म्यं मकरस्थे विभावसौ}
{वैशाखमासमाहात्म्यं तथा मेषगते रवौ} %॥३॥

\twolineshloka
{तत्र तत्र च ये धर्माः कथिताः सर्वशस्त्वया}
{एतेभ्योऽप्यधिकः कश्चिन्मासश्चेत्तव सम्मतः} %॥४॥

\twolineshloka
{धर्म ईशप्रियो नित्यं तं त्वं कथय साम्प्रतम्}
{यच्छ्रुत्वा पुनरन्यत्र श्रोतुमिच्छा न नो भवेत्} %॥५॥

\onelineshloka
{श्रद्धालोः श्रोतुरग्रे तु वक्ता गोप्यं न कारयेत्} %॥६॥


\uvacha{सूत उवाच}
\twolineshloka
{शृणुध्वं मुनयः सर्वे भवतो वाक्यगौरवात्}
{तुष्टोऽहं न च गोप्यं मे भवदग्रे तु किञ्चन} %॥७॥

\twolineshloka
{अदाम्भिक्यं तथास्तिक्यमशठत्वं सुभक्तिता}
{शुश्रूषत्वं विनीतत्वं ब्रह्मण्यत्वं सुशीलता} %॥८॥

\twolineshloka
{ध्रुवत्वं च शुचित्वं च तपस्वित्वानसूयते}
{एते द्वादशसङ्ख्याका गुणाः श्रोतुः प्रकीर्तिताः} %॥९॥

\twolineshloka
{ते सर्वेऽपि भवत्स्वेव तुष्टस्तत्त्वं ब्रवीम्यतः}
{सनत्कुमारो मेधावी धर्मजिज्ञासुरानतः} %॥१०॥

\onelineshloka
{ईश्वरं परिपप्रच्छ भक्त्या परमया युतः} %॥११॥


\uvacha{सनत्कुमार उवाच}
\twolineshloka
{देवदेव महाभाग तव योगिध्येयपदाम्बुज}
{व्रतानि बहुशस्त्वत्तः श्रुता धर्माश्च सर्वशः} %॥१२॥

\twolineshloka
{तथापि श्रोतुमिच्छेका वर्तते हृदि साम्प्रतम्}
{द्वादशस्वपि मासेषु मासः श्रेष्ठतमोऽस्ति यः} %॥१३॥

\twolineshloka
{प्रीतिकरोऽत्यन्तं सिद्धिदः सर्वकर्मणाम्}
{अन्यमासे कृतं कर्म तदेवास्मिन्कृतं यदि} %॥१४॥

\twolineshloka
{स्यादनन्तफलं देव तं मासं वक्तुमर्हसि}
{तत्रत्यान्सर्वधर्माश्च लोकानुग्रहकाम्यया} %॥१५॥


\uvacha{ईश्वर उवाच}
\twolineshloka
{सनत्कुमार वक्ष्यामि सुगोप्यमपि सुव्रत}
{शुश्रूषुत्वेन भक्त्या च प्रीतोऽस्मि विधिनन्दन} %॥१६॥

\twolineshloka
{द्वादशस्वपि मासेषु श्रावणो मेऽतिवल्लभः}
{श्रवणार्हं यन्माहात्म्यं तेनासौ श्रावणो मतः} %॥१७॥

\twolineshloka
{श्रवण पौर्णमास्यां ततोऽपि श्रावणः स्मृतः}
{यस्य श्रवणमात्रेण सिद्धिदः श्रावणोऽप्यतः} %॥१८॥

\twolineshloka
{स्वच्छत्वाच्च नभस्तुल्यो नभा इति ततः स्मृतः}
{तत्रत्यधर्मगणनां कर्तुं कः शक्नुयाद्भुवि} %॥१९॥

\twolineshloka
{सर्वतो यत्फलं वक्तुं चतुरास्योऽभवद्विधिः}
{द्रष्टुं यत्फलमाहात्म्यं सहस्त्राक्षोऽभवद् वृषा} %॥२०॥

\twolineshloka
{अनन्तो यत्फलं वक्तुं सहस्त्रद्वयजिह्वकः}
{किं बहूक्तेन कोऽप्येतद् द्रष्टुं वक्तुं च न क्षमः} %॥२१॥

\twolineshloka
{एतत्कलामपि मुने लभन्ते नान्यमासकाः}
{सर्वव्रतमयश्चैषः सर्वधर्ममयस्तथा} %॥२२॥

\twolineshloka
{नैकोऽपि वासरो यत्र व्रतशून्यः प्रदृश्यते}
{प्रायेण तिथयश्चापि व्रतवत्योऽत्र मासि वै} %॥२३॥

\twolineshloka
{अत्रोच्यते मया यद्यदर्थवादो न सोऽत्र हि}
{आर्तैर्जिज्ञासुभिर्भक्तैस्तथार्थार्थिमुमुक्षुभिः} %॥२४॥

\onelineshloka*
{चतुर्विधैरपि जनैः सेव्यः स्वस्वेष्टकाङ्क्षिभिः}

\uvacha{सनत्कुमार उवाच}
\twolineshloka
{भगवन् यत्त्वया प्रोक्तं व्रतशून्यो न वासरः}
{प्रायेण तिथिरप्यत्र तन्ममाचक्ष्व सत्तम} %॥२५॥

\twolineshloka
{कस्यां तिथौ किं व्रतं स्यात्कस्मिन्वारे च किं व्रतम्}
{तत्र तत्राधिकारी कः किं फलं कीदृशो विधिः} %॥२६॥

\twolineshloka
{केन केनापि चाचीर्णमुद्यापनविधिश्च कः}
{प्रधानं पूजनं कुत्र जागणश्चापि तद्विधिः} %॥२७॥

\twolineshloka
{को देवः कुत्र पूज्यः स्यात्सामग्री पूजनस्य का}
{कस्य व्रतस्य कः कालस्तत्सर्वं कथय प्रभो} %॥२८॥

\twolineshloka
{त्वत्प्रियश्च कथं मासः पवित्रः केन हेतुना}
{मासेऽस्मिन्नवतारः कः श्रेष्ठश्चायं कुतोऽभवत्} %॥२९॥

\twolineshloka
{अस्मिन्मासे च के धर्मा अनुष्ठेया वद प्रभो}
{प्रश्नेऽपि च कियज्ज्ञानं ममाज्ञस्य तवाग्रतः} %॥३०॥

\twolineshloka
{अशेषेण समाचक्ष्व पृष्टादन्यच्च यद्भवेत्}
{जनानां तारणार्थाय कृपालो कृपया वद} %॥३१॥

\twolineshloka
{रवौ सोमे भौमवारे बुधे सुरगुरौ कवौ}
{शनैश्चरदिने वापि तत्सर्वं वद मे विभो} %॥३२॥

\twolineshloka
{सर्वेषामादिभूतस्त्वमादिदेवस्ततः स्मृतः}
{एकस्य विधिबाधाभ्यामन्यबाधाविधी यथा} %॥३३॥

\twolineshloka
{अन्येषामल्पदेवत्वान्महादेवस्ततः स्मृतः}
{देवत्रयाश्रयेऽश्वत्थे उपर्यास्ते स्थितिस्तव} %॥३४॥

\twolineshloka
{शिवस्त्वं शुभरूपत्वादघौघहरणाद्धरः}
{तव चैवादिदेवत्वे प्रमाणं शुक्लवर्णकः} %॥३५॥

\twolineshloka
{प्रकृतौ शुक्लवर्णेऽन्ये वर्णाः स्युर्विकृतिं गताः}
{यतः कर्पूरगौरस्त्वमादिदेवस्ततो ह्यसि} %॥३६॥

\twolineshloka
{गणपत्याधारभूतान्मूलाधाराच्चतुर्दलात्}
{स्वाधिष्ठानाभिधात्पद्मात्षद्दलाद् ब्रह्मदैवतात्} %॥३७॥

\twolineshloka
{मणिपूराद्दशदलान्मण्डलाद्विष्ण्वधिष्ठितात्}
{ब्रह्मविष्णूपरिस्थस्त्वं वदतीदं च मुख्यताम्} %॥३८॥

\twolineshloka
{एकस्य तेऽर्चनाद्देव पञ्चायतनपूजनम्}
{जायतेऽन्यसुरे चैव सम्भवेन्न हि सर्वथा} %॥३९॥

\twolineshloka
{स्वयं शिवस्त्वं वामोरौ शक्तिर्गणपतिस्तथा}
{दक्षिणोरावक्ष्णि सूर्यो हृदये भक्तराड्ढरिः} %॥४०॥

\twolineshloka
{अन्नस्य ब्रह्मरूपत्वाद्रसात्मत्वाद्धरेरपि}
{भोक्तृत्वाच्च तवेशान श्रेष्ठत्वे कस्य संशयः} %॥४१॥


\threelineshloka
{विरक्तत्वं शिक्षयिष्यन् श्मशाने पर्वते स्थितिः}
{उतामृतत्वस्येशानो मन्त्रलिङ्गेन सूक्तके}
{पौरुषे प्रतिपाद्योऽसि इति प्राहुर्महर्षयः} %४२

\twolineshloka
{जगत्संहारकं हालाहलं केन धृतं गले}
{महाप्रलयकालाग्निं भाले धर्तुं च कः क्षमः} %॥४३॥

\twolineshloka
{भवान्धकूपपतने हेतुः केन हतः स्मरः}
{किं वर्ण्यं भागधेयं ते यद्वक्तुं हीदृशो भवान्} %॥४४॥

\twolineshloka
{त्वां स्तोतुं जन्मकोट्यापि वराकोऽहं न च क्षमः}
{कृत्वा मयि कृपामेव मत्प्रश्नान्वक्तुमर्हसि} %॥४५॥


{॥इति श्रीस्कन्दपुराणे ईश्वरसनत्कुमारसंवादे श्रावणमासमाहात्म्ये प्रथमोऽध्यायः॥१॥}



\sect{द्वितीयोऽध्यायः}

\uvacha{ईश्वर उवाच}
\twolineshloka
{साधु साधु महाभाग विनीतोऽसि विरिञ्चिज}
{श्रोता गुणयुतो यस्माच्छ्रद्धालुर्भक्तिभूषितः} %॥१॥

\twolineshloka
{श्रावणे मासि विनयाद्यत्पृष्टं भवतानघ}
{अपृष्टमपि ते वक्ष्ये प्रेम्णा परमया मुदा} %॥२॥

\twolineshloka
{प्रियो भवति चाद्वेष्टा नम्रस्त्वं च तथाविधः}
{पञ्चमो मस्तकश्छिन्नः प्रोद्धतस्य पितुस्तव} %॥३॥

\twolineshloka
{त्वं च तं मत्सरं त्यक्त्वा मां यतः शरणं गतः}
{अतो वक्ष्यामि ते तात भूत्वा चैकमनाः शृणु} %॥४॥

\twolineshloka
{कुर्यान्नक्तव्रतं योगिन् श्रावणे नियतो नरः}
{रुद्राभिषेकं कुर्वीत मासमात्रं दिने दिने} %॥५॥

\twolineshloka
{स्वप्रीतिविषयस्थापि कस्यचित्त्यागमाचरेत्}
{पुष्पैः फलैश्च धान्यैश्च तुलसीमञ्जरीदलैः} %॥६॥

\twolineshloka
{बिल्वपत्रैर्लक्षपूजां शङ्करस्य समाचरेत्}
{कोटिलिङ्गादि कर्तव्यं ब्राह्मणांश्चैव भोजयेत्} %॥७॥

\twolineshloka
{धारणापारणे कुर्यादुपोषणमथापि च}
{पञ्चामृताभिषेकं च मम प्रीतिकरं परम्} %॥८॥

\twolineshloka
{अस्मिन्मासे कृतं यद्यत्तदानन्त्याय कल्पते}
{भूमिशायी ब्रह्मचारी सत्यवाक्यो भवेन्मुने} %॥९॥

\twolineshloka
{न नयेन्मासमेनं तु व्रतवन्ध्यं कदाचन}
{अनोदनं समश्नीयाद्धविष्यान्नमथापि वा} %॥१०॥

\twolineshloka
{पत्रे चैव समश्नीयाच्छाकमात्रं त्यजेद् व्रती}
{किञ्चिद्व्रती सर्वथा स्याद्भक्तिमान्मुनिसत्तम} %॥११॥

\twolineshloka
{सदाचारो भूमिशायी प्रातः स्नायी जितेन्द्रियः}
{मत्पूजां प्रत्यहं कुर्यादेकाग्रकृतमानसः} %॥१२॥

\twolineshloka
{पुरश्चरणमप्यत्र मन्त्राणां सिद्धिदं परम्}
{शिवषड्वर्णमन्त्रस्य गायत्र्याश्च जपं चरेत्} %॥१३॥

\twolineshloka
{प्रदक्षिणा नमस्कारान् वेदपारायणं तथा}
{जपः पुरुषसूक्तस्य अधिकं फलदो भवेत्} %॥१४॥

\twolineshloka
{ग्रहयज्ञः कोटिहोमो लक्षहोमोऽयुतस्तथा}
{कृतः फलति सद्योऽत्र वाञ्छितार्थफलप्रदः} %॥१५॥

\twolineshloka
{अस्मिन्मासे चैकदिनं यो वन्ध्यं व्रततो नयेत्}
{स याति नरकं घोरं यावदाभूतसम्प्लवम्} %॥१६॥

\twolineshloka
{यथायं मे प्रियो मासस्तथा किञ्चिन्न मे प्रियम्}
{कामिनः फलदश्चायं निष्कामस्य तु मोक्षदः} %॥१७॥

\twolineshloka
{तत्र व्रतानि ये धर्मास्तान्मत्तः शृणु सत्तम}
{रवौ रविव्रतं सोमे मत्पूजा नक्तभोजनम्} %॥१८॥

\twolineshloka
{प्रथमं सोममारभ्य व्रतं स्याद्रोटकाभिधम्}
{सार्धमासत्रयं तत्स्यात्सर्वकामार्थसिद्धिदम्} %॥१९॥

\twolineshloka
{भौमे मङ्गलगौर्याश्च तदह्नोर्बुधजीवयोः}
{शुक्रे जीवन्तिकायाश्च आञ्जनेयनृसिंहयोः} %॥२०॥

\twolineshloka
{शनौ व्रतं समादिष्टं तिथिष्वथ मुने शृणु}
{नभः शुक्लद्वितीयायां व्रतमौदुम्बराभिधम्} %॥२१॥

\twolineshloka
{गौरीवतं तृतीयायां श्रावणे शुक्लपक्षके}
{तथा शुक्लचतुर्थ्यां तु दूर्वागणपतिव्रतम्} %॥२२॥

\twolineshloka
{विनायकीति तस्याश्च संज्ञा स्यादपरा मुने}
{नागानां पूजने शस्ता पञ्चमी शुक्लपक्षके} %॥२३॥

\twolineshloka
{इमां रौरवकल्यादि जानीहि मुनिसत्तम}
{सूपौदनव्रतं षष्ठ्यां सप्तम्यां शीतलाव्रतम्} %॥२४॥

\twolineshloka
{पवित्रारोपणं देव्या वसौ भूते यथा भवेत्}
{शुक्लकृष्णनवम्योस्तु नक्तव्रतविधिः स्मृतः} %॥२५॥

\twolineshloka
{दशम्यां शुक्लपक्षे तु आशासंज्ञं व्रतं भवेत्}
{पक्षद्वये विशेषोऽस्मिन्नेकादश्योस्तु कश्चन} %॥२६॥

\twolineshloka
{पवित्रारोपणं शुक्लद्वादश्यां तु हरेः स्मृतम्}
{द्वादश्यां श्रीधरं पूज्य परां गतिमवाप्नुयात्} %॥२७॥

\twolineshloka
{उत्सर्जनमुपाकर्म सभादीपस्तथैव च}
{उपाकर्म सभायां तु रक्षाबन्धस्ततः परम्} %॥२८॥

\twolineshloka
{श्रवणाकर्म तत्रैव तथा सर्पबलिः तथा सर्पबलिः स्मृतः}
{हयग्रीवस्यावतारः पूर्णिमायां तु सप्तकम्} %॥२९॥

\twolineshloka
{नभःकृष्णे तु सङ्कष्टचतुर्थीव्रतमुच्यते}
{ज्ञेया मानवकल्पादिः श्रावणे कृष्णपञ्चमी} %॥३०॥

\twolineshloka
{पूर्णावतारः कृष्णस्य कृष्णाष्टम्यां द्विजोत्तम}
{अवतारः समभवद् व्रतं तत्र महोत्सवम्} %॥३१॥

\twolineshloka
{इयं मन्वादिरपि च ज्ञातव्या मुनिपुङ्गव}
{अमायां श्रावणे मासि पिठोराव्रतमुच्यते} %॥३२॥

\twolineshloka
{कुशानां ग्रहणं चैव वृषभाणां च पूजनम्}
{शुक्लाद्यतिथिमारभ्य तत्तत्तिथिषु देवताः} %॥३३॥

\twolineshloka
{वह्निर्देवः प्रतिपदि द्वितीया ब्रह्मदेवता}
{तृतीयायास्तथा गौरी चतुर्थ्यां गणनायकः} %॥३४॥

\twolineshloka
{सर्पाधिपा पञ्चमी स्यात्षष्ठी स्यात्स्कन्ददेवता}
{सप्तम्यां भास्करो देवः शिवदेवाष्टमी तिथिः} %॥३५॥

\twolineshloka
{दुर्गाधिपा तु नवमी दशम्यन्तकदेवता}
{एकादश्यधिपाश्चैव विश्वेदेवाः प्रकीर्तिताः} %॥३६॥

\twolineshloka
{द्वादश्याश्च तु हरिः कामस्त्रयोदश्यधिपो मतः}
{चतुर्दश्यां शिवश्चैव पौर्णमास्याः शशी पतिः} %॥३७॥

\twolineshloka
{अमायाः पितरो देवा एते तिथ्यधिपाः स्मृताः}
{स देवस्तत्र पूज्यः स्याद्यस्य देवस्य या तिथिः} %॥३८॥

\twolineshloka
{अगस्त्यस्योदयो मासे प्रायेणात्रैव जायते}
{कथयामि च तं कालं शृणुष्वैकमना मुने} %॥३९॥

\twolineshloka
{सिंहसङ्क्रान्तिदिवसाद्यदा द्वादश यान्ति च}
{चत्वारिंशच्च घटिकास्तदागस्योदयो भवेत्} %॥४०॥

\twolineshloka
{सप्ताहानि ततः पूर्वमगस्त्यार्घ्यं समाचरेत्}
{द्वादशेष्वपि मासेषु आदित्यो भिन्नसंज्ञया} %॥४१॥

\twolineshloka
{तपते श्रावणे तत्र गभस्तिरितिसंज्ञितः}
{तत्पूजनं च कर्तव्यं मासेऽस्मिन्भक्तितत्परैः} %॥४२॥

\twolineshloka
{चतुर्षु यानि मासेषु वर्ज्यानि शृणु सत्तम}
{श्रावणे च त्यजेच्छाकं दधि भाद्रपदे तथा} %॥४३॥

\twolineshloka
{दुग्धमाश्वयुजे मासि कार्तिके द्विदलं त्यजेत्}
{इत्यादीनि समस्तानि तानि कर्तुमशक्नुवन्} %॥४४॥

\twolineshloka
{एकस्मिन् श्रावणे मासि कुर्वंस्तत्फलभाग्भवेत्}
{उद्देशोऽयं मया प्रोक्तः सङ्क्षेपात्तव मानद} %॥४५॥

\twolineshloka
{अत्रत्यानां व्रतानां तु धर्माणां मुनिसत्तम}
{केनापि विस्तरो वक्तुं नालं वर्षशतैरपि} %॥४६॥

\twolineshloka
{मम प्रीत्यै हरेर्वापि कुर्याद् व्रतमशेषतः}
{आवयोर्नहि भेदोऽस्ति परमार्थविचारतः} %॥४७॥

\twolineshloka
{कल्पयन्त्यत्र ये भेदं ते वै निरयगामिनः}
{सनत्कुमार तस्मात्त्वं श्रावणे धर्ममाचर} %॥४८॥


{॥इति श्रीस्कन्दपुराणे ईश्वरसनत्कुमारसंवादे श्रावणमासमाहात्म्ये श्रावणव्रतोद्देशकथनं नाम द्वितीयोऽध्यायः॥२॥}




\sect{तृतीयोऽध्यायः}

\uvacha{सनत्कुमार उवाच}
\twolineshloka
{भगवन् व्रतसङ्घस्य उद्देशः कथितस्त्वया}
{तृप्तिर्न जायते स्वामिन् विस्तराद्वक्तुमर्हसि} %॥१॥

\onelineshloka
{यच्छ्रुत्वा कृतकृत्योऽहं भविष्यामि सुरेश्वर} %॥२॥


\uvacha{ईश्वर उवाच}
\twolineshloka
{नक्तव्रतेन योगीश श्रावणं यो नयेत्सुधीः}
{द्वादशस्वपि मासेषु स नक्तफलभाग्भवेत्} %॥३॥

\twolineshloka
{दिनावसानपूर्व तु नक्तं स्याद्रात्रिभोजनम्}
{तत्राद्यास्त्रिघटीस्त्यक्त्वा कालः स्यान्नक्तभोजने} %॥४॥

\twolineshloka
{ततः सन्ध्या त्रिघटिका अस्तादुपरि भास्वतः}
{चत्वारीमानि कर्माणि सन्ध्यायां परिवर्जयेत्} %॥५॥

\twolineshloka
{आहारं मैथुनं निद्रां स्वाध्यायं च चतुर्थकम्}
{गृहस्थयतिभेदेन तद्व्यवस्थां च मे शृणु} %॥६॥

\twolineshloka
{आत्मनो द्विगुणी छाया मन्दीभवति भास्करे}
{यतेर्नक्तं तु तत्प्रोक्तं न नक्तं निशि भोजनम्} %॥७॥

\twolineshloka
{नक्षत्रदर्शनान्नक्तं गृहस्थस्य बुधैः स्मृतम्}
{यतेर्दिनाष्टमे भागे रात्रौ तस्य निषिध्यते} %॥८॥

\twolineshloka
{नक्तं निशायां कुर्वीत गृहस्थो विधिसंयुतः}
{यतिश्च विधवा चैव विधुरश्च ससूर्यकम्} %॥९॥

\twolineshloka
{विधुरश्चेत्पुत्रवान् स्यात्स तु रात्रौ समाचरेत्}
{अनाश्रमोऽप्याश्रमी स्यादपत्नीकोऽपि पुत्रवान्} %॥१०॥

\twolineshloka
{एवं यथाधिकारं तु कुर्यान्नक्तव्रतं सुधीः}
{अत्र नक्तव्रती मासे परां गतिमवाप्नुयात्} %॥११॥

\twolineshloka
{प्रातःस्नानं करिष्यामि ब्रह्मचर्यव्रते स्थितः}
{भोक्ष्यामि नक्तं भूशय्यां करिष्ये प्राणिनां दयाम्} %॥१२॥

\twolineshloka
{प्रारब्धेऽस्मिन्वते देव यद्यपूर्णे म्रिये ह्यहम्}
{तदा सम्पूर्णतां यातु प्रसादात्ते जगत्पते} %॥१३॥

\twolineshloka
{इति सङ्कल्प्य मेधावी मासे नक्तव्रतं चरेत्}
{एवं नक्तव्रतं कुर्वन्मम प्रियतमो भवेत्} %॥१४॥

\twolineshloka
{विप्रद्वारातिरुद्रेण महारुद्रेण वा स्वयम्}
{अभिषेकं मासमात्रं रुद्रेण प्रत्यहं चरेत्} %॥१५॥

\twolineshloka
{तस्य प्रीणाम्यहं वत्स जलधाराप्रियो यतः}
{कुर्याद्रुद्रेण वा होमं मम प्रीतिकरं परम्} %॥१६॥

\twolineshloka
{स्वस्य यद्रोचतेऽत्यन्तं भोज्यं वा भोग्यमेव वा}
{सङ्कल्प्य द्विजवर्याय दत्त्वा मासे स्वयं त्यजेत्} %॥१७॥

\twolineshloka
{अतः परं शृणु मुने लक्षपूजाविधिं परम्}
{श्रीकामो बिल्वपत्रैश्च दूर्वाभिः शान्तिकामुकः} %॥१८॥

\twolineshloka
{आयुष्कामेन कर्तव्यं चम्पकैः पूजनं हरेः}
{विद्याकामेन कर्तव्यं मल्लिकाजातिभिस्तथा} %॥१९॥

\twolineshloka
{शिवविष्ण्वोः प्रसन्नत्वं तुलसीभिः प्रसिध्यति}
{पुत्रकामेन कर्तव्यं बार्हतैः पूजनं शुभम्} %॥२०॥

\twolineshloka
{दुःस्वप्नप्रशमार्थाय शस्तं धान्यप्रपूजनम्}
{रङ्गवल्ल्यादिभिर्देवं देवस्याग्रे विनिर्मितैः} %॥२१॥

\twolineshloka
{पद्मादिभिः स्वस्तिकाद्यैश्चक्राद्यैः पूजयेद्विभुम्}
{एवं हि सर्वपुष्यैश्च सर्वकामार्थसिद्धये} %॥२२॥

\twolineshloka
{लक्षपूजां प्रकुर्याच्चेत्सुप्रसन्नो हरो भवेत्}
{उद्यापनं ततः कार्यं मण्डपं चैव साधयेत्} %॥२३॥

\twolineshloka
{वेदिका च प्रकर्तव्या मण्डपस्य त्रिभागतः}
{पुण्याहवाचनं कृत्वा आचार्यं वरयेत्ततः} %॥२४॥

\twolineshloka
{गीतवादित्रनिर्घोषैर्ब्रह्मघोषेण भूयसा}
{प्रविश्य मण्डपे तस्मिन् रात्रौ जागरणं चरेत्} %॥२५॥

\twolineshloka
{वेदिकायां प्रकर्तव्यं लिङ्गतोभद्रमुत्तमम्}
{तन्मध्ये तण्डुलैः कुर्यात्कैलासं च सुशोभनम्} %॥२६॥

\twolineshloka
{कलशं स्थापयेत्तत्र ताम्र चैव महाप्रभम्}
{पञ्चपल्लवसंयुक्तं न्यसेद् वस्त्रं सुसूक्ष्मकम्} %॥२७॥

\twolineshloka
{सौवर्णी प्रतिमां तत्र स्थापयेत्पार्वतीपतेः}
{पूजां तत्र प्रकुर्वीत पञ्चामृतपुरःसरैः} %॥२८॥

\twolineshloka
{सनैवेद्यैर्गीतवादित्रनृत्यकैः}
{वेदशास्त्रपुराणैश्च रात्रौ जागरणं चरेत्} %॥२९॥

\twolineshloka
{प्रभातसमये सुस्नातः सुशुचिर्भवेत्}
{स्थण्डिलं कारयेत्तत्र स्वशाखोक्तविधानतः} %॥३०॥

\twolineshloka
{होमं च सतिलाज्येन पायसेन च कारयेत्}
{मूलमन्त्रेण गायत्र्या शिवनाम्नां सहस्रकैः} %॥३१॥

\twolineshloka
{येन मन्त्रेण पूजा तु कृता तेनैव होमयेत्}
{शर्कराघृतमिश्रेण चरुणा जुहुयात्ततः} %॥३२॥

\twolineshloka
{ततः स्विष्टकृतं हुत्वा पूर्णाहुतिमनन्तरम्}
{आचार्यं पूजयेत्सम्यग्वस्त्रालङ्कारभूषणैः} %॥३३॥

\twolineshloka
{ब्राह्मणान्पूजयेत्पश्चात्तेभ्यो दद्याच्च दक्षिणाम्}
{येन येन प्रकुर्याच्च लक्षपूजामुमापतेः} %॥३४॥

\twolineshloka
{तत्तद्दद्यात्सुवर्णेन कृत्वा शम्भुं प्रपूजयेत्}
{यदि दीपः कृतस्तेन तद्दानं चैव कारयेत्} %॥३५॥

\twolineshloka
{सुवर्णवर्तिकां कृत्वा दीपमात्रं च रौप्यकम्}
{गोघृतेन समायुक्तं सर्वकामार्थसिद्धये} %॥३६॥

\twolineshloka
{क्षमापयेत्ततो देवं ब्राह्मणान्भोजयेच्छतम्}
{एवं यः कुरुते पूजां तस्य प्रीणाम्यहं मुने} %॥३७॥

\twolineshloka
{तत्रापि श्रावणे कुर्यात्तदानन्त्याय कल्पते}
{स्वप्रीतिविषयः कश्चित्पदार्थस्त्यज्यते यदि} %॥३८॥

\twolineshloka
{मदर्पणधिया चात्र मासि तस्य फलं शृणु}
{इहामुत्र च तत्प्राप्तिर्भवेल्लक्षगुणाधिका} %॥३९॥

\twolineshloka
{सकामत्वे तु चैवं स्यान्निष्कामत्वे परा गतिः}
{रुद्राभिषेकं कुर्वाणस्तत्रत्याक्षरसङ्ख्यया} %॥४०॥

\twolineshloka
{प्रत्यक्षरं कोटिवर्ष रुद्रलोके महीयते}
{पञ्चामृतस्याभिषेकादमृतत्वं समश्नुते} %॥४१॥

\twolineshloka
{तस्मिन्मासे भूमिशायी फलं तस्यापि मे शृणु}
{प्रवालनिर्मितां श्रेष्ठां गजदन्तभवामपि} %॥४२॥

\twolineshloka
{पाटीरनिर्मितां चापि खचितां नवरत्नकैः}
{निःसीममृदुपक्षीन्द्रविशेषां द्विजसत्तम} %॥४३॥

\twolineshloka
{वरवस्त्रेण सञ्छन्नां तूलिकां चात्र शोभनाम्}
{दशोपबर्हणैर्युक्तां शय्यां स लभते शुभाम्} %॥४४॥

\twolineshloka
{रम्याङ्गनासमायुक्तां रत्नदीपविभूषिताम्}
{ब्रह्मचर्येण चाप्यत्र वीर्यपुष्टिर्भवेद् दृढा} %॥४५॥

\twolineshloka
{ओजो बलं देहदार्ढ्यं यद्धर्मस्योपकारकम्}
{प्रत्यक्षैव भवेत्तस्य ब्रह्मप्राप्तिर्न संशयः} %॥४६॥

\twolineshloka
{निष्कामस्य सकामस्य स्वर्गं देवाङ्गना शुभा}
{अत्र मौनव्रतधरो महान्वक्ता प्रजायते} %॥४७॥

\twolineshloka
{अहोरात्रं दिने वापि भुक्तिकालेऽथवा पुनः}
{घण्टायाः पुस्तकस्यापि व्रतान्ते दानमाचरेत्} %॥४८॥

\twolineshloka
{सर्वशास्त्रप्रवीणः स्याद्वेदवेदाङ्गपारगः}
{वाचस्पतिसमो बुद्धौ मौनमाहात्म्यतो भवेत्} %॥४९॥

\onelineshloka
{मौनिनः कलहो नास्ति तस्मान्मौनव्रतं परम्} %॥५०॥


{॥इति श्रीस्कन्दपुराणे ईश्वरसनत्कुमारसंवादे श्रावणमासमाहात्म्ये नक्तव्रतलक्षपूजा-भूमिशयनमौनादिव्रतकथनं नाम तृतीयोऽध्यायः॥३॥}



\sect{चतुर्थोऽध्यायः}

\uvacha{ईश्वर उवाच}
\twolineshloka
{सनत्कुमार वक्ष्यामि धारणापारणाव्रतम्}
{पुण्याहं वाचयेत्पूर्वमारभ्य प्रतिपद्दिनम्} %॥१॥

\twolineshloka
{सङ्कल्पयेन्मम प्रीत्यै धारणापारणाव्रतम्}
{एकस्मिन्धारणं कुर्यात्पारणं च तथापरे} %॥२॥

\twolineshloka
{उपवासो धारणे स्यात्पारणे भोजनं भवेत्}
{समाप्ते मासि चैवात्र कुर्यादुद्यापनं व्रती} %॥३॥

\twolineshloka
{समाप्ते श्रावणे मासि पुण्याहं कारयेत्पुरा}
{आचार्यं वरयेत्पश्चाद् ब्राह्मणांश्चैव मानद} %॥४॥

\twolineshloka
{पार्वतीशङ्करस्यापि प्रतिमां स्वर्णनिर्मिताम्}
{पूर्णकुम्भे तु संस्थाप्य पूजयेन्निशि भक्तितः} %॥५॥

\twolineshloka
{रात्रौ जागरणं कुर्यात्पुराणश्रवणादिभिः}
{प्रातरग्निं समाधाय होमं कुर्याद्यथाविधि} %॥६॥

\twolineshloka
{त्र्यम्बकेणेति मन्त्रेण जुहुयाच्च तिलौदनम्}
{तथैव शिवगायत्र्या जुहुयाच्च घृतौदनम्} %॥७॥

\twolineshloka
{षडक्षरेण मन्त्रेण पायसं जुहुयात्ततः}
{पूर्णाहुतिं ततो हुत्वा होमशेषं समापयेत्} %॥८॥

\twolineshloka
{ब्राह्मणान्भोजयेत्पश्चादाचार्यं चैव पूजयेत्}
{एवं कृत्वा महाभाग ब्रह्महत्यादिपातकैः} %॥९॥

\twolineshloka
{मुच्यते नात्र सन्देहस्तस्मात्कुर्यान्महाव्रतम्}
{शृणु मासोपवासस्य श्रावणे विधिमादरात्} %॥१०॥

\twolineshloka
{सङ्कल्पयेत्तु प्रतिपद्दिने प्रातर्वतं मुने}
{नारी वा पुरुषो वापि संयतात्मा जितेन्द्रियः} %॥११॥

\twolineshloka
{ततोऽर्चयेदमायां तु शङ्करं लोकशङ्करम्}
{सम्पूजयेत् षोडशभिरुपचारैर्वृषध्वजम्} %॥१२॥

\twolineshloka
{ब्राह्मणान्पूजयेच्चैव वस्त्रालङ्करणादिभिः}
{भोजयेच्च यथाशक्त्या प्रणिपत्य विसर्जयेत्} %॥१३॥

\twolineshloka
{एवं मासोपवासस्तु मम प्रीतिकरो भवेत्}
{रुद्रवर्तिविधानं च सम्मितं लक्षसङ्ख्यया} %॥१४॥

\twolineshloka
{शृणुष्वावहितो भूत्वा सर्वसिद्धिकरं नृणाम्}
{कार्पासतन्तुभिः कार्या एकादशभिरादरात्} %॥१५॥

\twolineshloka
{वर्तयस्ता रुद्रवर्तिसंज्ञाः प्रीतिकरा मम}
{श्रावणस्याद्यदिवसे सङ्कल्प्य विधिपूर्वकम्} %॥१६॥

\twolineshloka
{देवदेवं महादेवं लक्षवर्तिभिरादरात्}
{नीराजयामि गौरीशं श्रावणे मासि भक्तितः} %॥१७॥

\twolineshloka
{पूजयित्वा प्रतिदिनं वर्तीनां तु सहस्रशः}
{नीराजयेदन्त्यदिने सहस्राण्येकसप्ततिः} %॥१८॥

\twolineshloka
{अथवा प्रत्यहं दद्यात्सहस्त्रत्रयमादृतः}
{चरमे तु दिने दद्यात्सहस्त्राणि त्रयोदश} %॥१९॥

\twolineshloka
{एकस्मिन्वा दिने रुद्रवर्तिलक्षं प्रदीपयेत्}
{सुघृतेनापि बहुना स्निग्धास्ता मम वल्लभाः} %॥२०॥

\onelineshloka
{पूजयित्वा तु विश्वेशं शृणुयाच्च कथां ततः} %॥२१॥


\uvacha{सनत्कुमार उवाच}

\threelineshloka
{देवदेव जगन्नाथ जगदानन्दकारक}
{व्रतस्यास्य प्रभावं मे कृपां कृत्वा वद प्रभो}
{केन चीर्णं व्रतमिदं विधिरुद्यापने कथम्} %२२


\uvacha{ईश्वर उवाच}
\twolineshloka
{शृणु वैधात्र यत्नेन व्रतानामुत्तमं व्रतम्}
{रुद्रवर्त्या महापुण्यं सर्वोपद्रवनाशनम्} %॥२३॥

\twolineshloka
{प्रीतिसौभाग्यजननं पुत्रपौत्रसमृद्धिदम्}
{शङ्करप्रीतिजननं शिवलोकं परं पदम्} %॥२४॥

\twolineshloka
{रुद्रवर्तिसमं नास्ति त्रिषु लोकेषु सुव्रतम्}
{अत्रैवोदाहरन्तीममितिहासं पुरातनम्} %॥२५॥

\twolineshloka
{क्षिप्रानद्यास्तटे रम्ये पुरी उज्जयिनी शुभा}
{तस्यामासीत्सुगन्धाख्या वारस्त्री ह्यतिसुन्दरा} %॥२६॥

\twolineshloka
{तया शुल्कं कृतं स्वीये सुरते तु सुदुःसहम्}
{सुवर्णानां शतं लोके प्रतिज्ञां कृतवत्यथ} %॥२७॥

\twolineshloka
{युवानश्च तया विप्रा भ्रंशिताश्च सुगन्धया}
{राजानो राजपुत्राश्च नग्नीकृत्य पुनः पुनः} %॥२८॥

\twolineshloka
{तेषां भूषां गृहीत्वा तु धिक्कृतास्तु सुगन्धया}
{एवं हि बहवो लोका लुण्ठितास्ते सुगन्धया} %॥२९॥

\twolineshloka
{तस्यास्तु देहगन्धेन कोशमात्रं सुगन्धितम्}
{रूपलावण्यकान्तीनां सा प्रतिष्ठा धरातले} %॥३०॥

\twolineshloka
{षट्त्रिंशद्रागभार्याणां षड्रागाणां गायने}
{तत्सन्तत्या अनन्तायाः कुशला गानकर्मणि} %॥३१॥

\twolineshloka
{नृत्ये रम्भादिकाः सर्वास्तर्जयन्ती सुराङ्गनाः}
{गत्या गजांश्च हंसांश्च विहसन्ती पदे पदे} %॥३२॥

\twolineshloka
{कामबाणान्प्रेरयन्ती कटाक्षैर्भूकृतैश्च तैः}
{कदाचित्सा गता क्षिप्रां कौतुकाविष्टमानसा} %॥३३॥

\twolineshloka
{ददर्श सा मनोरम्यामृषिभिः परिसेविताम्}
{केचिद् ध्यानपरा विप्राः केचिज्जपपरायणाः} %॥३४॥

\twolineshloka
{रताः शिवार्चने केचिद्विष्णोः केचित्प्रपूजने}
{तेषां मध्ये वसिष्ठो हि तया दृष्टो महामुने} %॥३५॥

\twolineshloka
{तस्या धर्मेऽभवद् बुद्धिस्तद्दर्शनमहत्त्वतः}
{विगताशा जीवनेऽपि विषयेषु विशेषतः} %॥३६॥

\twolineshloka
{विनम्रकन्धरा भूत्वा प्रणिपत्य पुनः पुनः}
{स्वकर्मपरिहाराय पप्रच्छ मुनिपुङ्गवम्} %॥३७॥


\uvacha{सुगन्धोवाच}

\threelineshloka
{अनाथनाथ विप्रेन्द्र सर्वविद्याविशारद}
{मया कृतानि योगीश पापानि सुबहून्यपि}
{तत्सर्वपापनाशाय उपायं ब्रूहि मे प्रभो} %३८


\uvacha{ईश्वर उवाच}
\twolineshloka
{एवमुक्तस्तया विप्रो वसिष्ठो मुनिरादरात्}
{तस्याः कर्म प्रविज्ञाय सोऽब्रवीद्दीनवत्सलः} %॥३९॥


\uvacha{वसिष्ठ उवाच}
\twolineshloka
{शृणुष्वावहिता भूत्वा तव पापस्य सङ्क्षयः}
{येन जायेत पुण्येन तत्सर्वं कथयामि ते} %॥४०॥

\twolineshloka
{गच्छ वाराणसीं भद्रे त्रिषु लोकेषु विश्रुताम्}
{गत्वा कुरुष्व तत्क्षेत्रे व्रतं त्रैलोक्यदुर्लभम्} %॥४१॥

\twolineshloka
{रुद्रवर्त्य महापुण्यं शिवप्रीतिकरं परम्}
{कृत्वा व्रतमिदं भद्रे प्राप्स्यसि त्वं परां गतिम्} %॥४२॥


\uvacha{ईश्वर उवाच}
\twolineshloka
{ततः सा कोशमादाय भृत्यं चैव समित्रकम्}
{गत्वा काशीं व्रतं चक्रे वसिष्ठोक्तविधानतः} %॥४३॥

\onelineshloka
{ततः सा सशरीरेण तस्मिन् लिङ्गे लयं ययौ} %॥४४॥

\twolineshloka
{एवं या कुरुते नारी व्रतमेतत्सुदुर्लभम्}
{यं यं चिन्तयते कामं तं तं प्राप्नोत्यसंशयः} %॥४५॥

\twolineshloka
{माहात्म्यं शृणु माणिक्यवर्तीनामपि सुव्रत}
{व्रतेन तासां विप्रेन्द्र मदर्धासनभागिनी} %॥४६॥

\twolineshloka
{जायते मत्प्रिया सा हि यावदाभूतसम्प्लवम्}
{उद्यापनमथो वक्ष्ये व्रतसम्पूर्णहेतवे} %॥४७॥

\twolineshloka
{कलशे स्थापयेद्देवमुमया सहितं शिवम्}
{सुवर्णनिर्मितं देवं वृषभे रौप्यनिर्मिते} %॥४८॥

\twolineshloka
{विधिना पूजनं कृत्वा रात्रौ जागरणं चरेत्}
{ततः प्रभाते विमले स्नात्वा नद्यां विधानतः} %॥४९॥

\twolineshloka
{आचार्यं वरयेद्भक्त्या द्विजैरेकादशैः सह}
{होमश्चैव प्रकर्तव्यो घृतपायसबिल्वकैः} %॥५०॥

\twolineshloka
{रुद्रसूक्तेन गायत्र्या मूलमन्त्रेण वा पुनः}
{ततः पूर्णाहुतिं हुत्वा आचार्यादीन्प्रपूजयेत्} %॥५१॥

\twolineshloka
{तथैकादश सद्विप्रान्सपत्नीकांस्तु भोजयेत्}
{एवं या कुरुते नारी सर्वपापैः प्रमुच्यते} %॥५२॥

\twolineshloka
{कथां श्रुत्वा विधानेन स्थाप्यं सर्वं न्यवेदयेत्}
{अश्वमेधसहस्रस्य फलं भवति निश्चितम्} %॥५३॥


{॥इति श्रीस्कन्दपुराणे ईश्वरसनत्कुमारसंवादे श्रावणमासमाहात्म्ये धारणापारणा-मासोपवासरुद्रवर्तिकथनं नाम चतुर्थोऽध्यायः॥४॥}



\sect{पञ्चमोऽध्यायः}

\uvacha{ईश्वर उवाच}
\twolineshloka
{माहात्म्यं कोटिलिङ्गानां पुण्यं वक्तुं न शक्यते}
{एकैकस्यापि लिङ्गस्य किं पुनः कोटिसङ्ख्यया} %॥१॥

\twolineshloka
{अशक्तौ लक्षकं कुर्यात्सहस्त्रमथवा शतम्}
{एकस्यापि हि लिङ्गस्य कारणान्मम सन्निधिः} %॥२॥

\twolineshloka
{षडक्षरेण मन्त्रेण पूजा कार्या स्मरद्विषः}
{उपचारैः षोडशभिर्भक्तियुक्तेन चेतसा} %॥३॥

\twolineshloka
{उद्यापनं च कर्तव्यं ग्रहयज्ञपुरःसरम्}
{सम्पादनीयो होमश्च ब्राह्मणांश्चैव भोजयेत्} %॥४॥

\twolineshloka
{नाकालमरणं तस्य वन्ध्यत्वहरणं परम्}
{सर्वापत्तिक्षयकरं सर्वसम्पत्प्रवर्धनम्} %॥५॥

\twolineshloka
{प्रेत्य कैलासवासश्च आकल्पं मम सन्निधौ}
{पञ्चामृताभिषेकं च यः कुर्याच्छ्रावणे नरः} %॥६॥

\twolineshloka
{स पञ्चामृतभोजी स्यात्सम्पन्नो गोधनेन च}
{अत्यन्तं मधुरालापो प्रियश्च त्रिपुरद्विषः} %॥७॥

\twolineshloka
{अनोदनव्रती चैव हविष्याशी च यो नरः}
{व्रीह्यादिसर्वधान्यानामक्षय्योऽसौ निधिर्भवेत्} %॥८॥

\twolineshloka
{पत्रावल्यां तु भुञ्जानः स्वर्णभाजनभोजनः}
{शाकवर्जनतः स्याद्वै शाककर्ता नरोत्तमः} %॥९॥

\twolineshloka
{केवलं भूमिशायी तु कैलासे वा समाप्नुयात्}
{प्रातः स्नानान्नभोमासि अब्दं तत्फलभाङ्मतः} %॥१०॥

\twolineshloka
{जितेन्द्रियत्वान्मासेऽस्मिन्बलमैन्द्रियकं भवेत्}
{स्फाटिकेऽश्ममये वापि मार्ने मारकतेऽपि वा} %॥११॥

\twolineshloka
{स्वयम्भौ वास्वयम्भौ वा पैष्टे धातुमयेऽपि वा}
{चन्दने नावनीते वा अन्यस्मिन्वापि लिङ्गके} %॥१२॥

\twolineshloka
{सकृत्पूजां प्रकुर्वाणो ब्रह्महत्याशतं दहेत्}
{सूर्यचन्द्रोपरागेषु सिद्धिः क्षेत्रेऽपि वा क्वचित्} %॥१३॥

\twolineshloka
{सिद्धिर्या लक्षजाप्येन सकृत्स्याज्जपतोऽत्र सा}
{अन्यकाले कृता ये स्युर्नमस्काराः प्रदक्षिणाः} %॥१४॥

\twolineshloka
{सहस्रेण फलं यत्स्यान्मासेऽस्मिन्नेकवारतः}
{मत्प्रिये श्रावणे मासि वेदपारायणे कृते} %॥१५॥

\twolineshloka
{सर्वेषां वेदमन्त्राणां सिद्धिः सम्यक्प्रजायते}
{मासेऽस्मिन्पौरुषं सूक्तं जपते श्रद्धयान्वितः} %॥१६॥

\twolineshloka
{कृत्वा सङ्ख्यासहस्त्रं तु कलौ स्यात्तु चतुर्गुणम्}
{वर्णानां सङ्ख्यया वापि शतं कुर्यादतन्द्रितः} %॥१७॥

\twolineshloka
{अशक्तः सङ्ख्यया कर्तुं शतमानेन वा जपेत्}
{बह्महत्यादिपापेभ्यो मुच्यते नात्र संशयः} %॥१८॥

\twolineshloka
{गुरुतल्पे कृते पापे प्रायश्चित्तमिदं परम्}
{नास्त्येतत्सदृशं पुण्यं पवित्रं पापनाशनम्} %॥१९॥

\twolineshloka
{विना पौरुषजाप्येन न नयेदेकमप्यहः}
{अर्थवादमिमं ब्रूयात्स नरो निरयी भवेत्} %॥२०॥

\twolineshloka
{ग्रहयज्ञः प्रकर्तव्यः समिच्चरुतिलाज्यकैः}
{धूपगन्धप्रसूनादिनैवेद्यादिप्रभेदतः} %॥२१॥

\twolineshloka
{तद्रूपाणां च ध्यानादि सम्पाद्य च यथायथम्}
{कोटिहोमो लक्षहोमो ऽयुतहोमस्तु शक्तितः} %॥२२॥

\twolineshloka
{तिलैर्व्याहतिभिः कार्यों ग्रहयज्ञाभिधोऽप्यसौ}
{अथ वक्ष्यामि वाराणां व्रतानि शृणु साम्प्रतम्} %॥२३॥

\twolineshloka
{तत्रादौ रविवारस्य व्रतं वक्ष्यामि तेऽनघ}
{अत्राप्युदाहरन्तीममितिहासं पुरातनम्} %॥२४॥

\twolineshloka
{प्रतिष्ठानपुरे रम्ये सुकर्मा नाम वै द्विजः}
{आसीद्दरिद्रः कृपणो भैक्ष्यचर्यापरायणः} %॥२५॥

\twolineshloka
{एकदा स गतो धान्यं याचितुं पर्यटन्पुरम्}
{स्त्रियो ददर्श सदने कस्यचिद् गृहमेधिनः} %॥२६॥

\twolineshloka
{चरन्त्यो रविवारस्य मिलिता व्रतमुत्तमम्}
{तदोचुस्ताश्च तं दृष्ट्वा आच्छादयत सत्वरम्} %॥२७॥

\twolineshloka
{पूजाविधिं ततो विप्रः प्रार्थयामास ताः स्त्रियः}
{छाद्यते किं नु भो साध्व्यो भवतीभिरिदं व्रतम्} %॥२८॥

\twolineshloka
{कथयध्वं कृपां कृत्वा ममोपरि दयालवः}
{परोपकारसदृशो धर्मो नास्ति जगत्त्रये} %॥२९॥

\twolineshloka
{साधूनां समचित्तानां परार्थः स्वार्थ एव हि}
{दरिद्रः पीडितश्चाहं श्रुत्वेदं व्रतमुत्तमम्} %॥३०॥

\onelineshloka
{चरिष्यामि विधिं ब्रूत फलं चास्य व्रतस्य हि} %॥३१॥


\uvacha{स्त्रिय ऊचुः}
\twolineshloka
{उन्मादं वा प्रमादं वा विस्मृतिं वा करिष्यसि}
{अभक्तिं वाप्यनास्थां वा कथं देयं तव द्विज} %॥३२॥

\twolineshloka
{इति तासां वचः श्रुत्वा विप्रेन्द्रो वाक्यमब्रवीत्}
{ज्ञानवानस्मि भो साध्व्यो भक्तिमांश्चास्मि सुव्रताः} %॥३३॥

\twolineshloka
{एवं तद्वचनं श्रुत्वा प्रौढा तासु च याभवत्}
{सा प्रोवाच व्रतं तस्मै यथाभूतं च तद्विधिम्} %॥३४॥

\twolineshloka
{श्रावणे शुक्लपक्षे तु प्रथमे रविवासरे}
{मौनेनोत्थायावगाहं कुर्याच्छीतोदकेन तत्} %॥३५॥

\twolineshloka
{स्वनित्यकर्म सम्पाद्य नागवल्लीदले शुभे}
{परिधिद्वादशयुतं मण्डलं तत्र संलिखेत्} %॥३६॥

\twolineshloka
{अर्कवद्वर्तुले सम्यग् रक्तचन्दनतः शुभम्}
{तत्र संज्ञायुतं सूर्यं पूजयेद्रक्तचन्दनात्} %॥३७॥

\twolineshloka
{जानुभ्यामवनीं गत्वा अर्घ्यं द्वादशमण्डलैः}
{रक्तचन्दनमिश्रं च जपाकुसुमसंयुतम्} %॥३८॥

\twolineshloka
{दद्याद् गभस्तये सम्यक् श्रद्धाभक्तिपुरःसरम्}
{रक्ताक्षतैर्जपापुष्पैस्तथान्यैरुपचारकैः} %॥३९॥

\twolineshloka
{नारीकेलस्य बीजं तु खण्डशर्करया युतम्}
{नैवेद्यमर्पयित्वा तु मन्त्रैरादित्यलिङ्गकः} %॥४०॥

\twolineshloka
{स्तुवीत द्वादशवरैर्नमस्कारान्प्रदक्षिणाः}
{षट्तन्तुनिर्मितं सूत्रं षड्भिर्ग्रन्थिभिरन्वितम्} %॥४१॥

\twolineshloka
{अर्पयित्वा तु देवेशं बध्नीयात्तु गले च तत्}
{द्विजाय वायनं दद्यात्फलैर्द्वादशभिर्युतम्} %॥४२॥

\twolineshloka
{एतद्व्रतप्रकारं न श्रावयेत्कस्यचित्पुरा}
{एवं व्रते कृते विप्र निर्धनो धनमाप्नुयात्} %॥४३॥

\twolineshloka
{अपुत्रो लभते पुत्रं कुष्ठी कुष्ठात्प्रमुच्यते}
{बद्धः स्याद्बन्धरहितो रोगी रोगेण हीयते} %॥४४॥

\twolineshloka
{किं बहूक्तेन विप्रेन्द्र यद्यदिच्छति वाञ्छितम्}
{तत्तल्लभेत्साधकोऽसौ व्रतस्यास्य प्रभावतः} %॥४५॥

\twolineshloka
{एवं चतुर्षु वारेषु कदाचिदपि पञ्चसु}
{उद्यापनं ततः कार्य व्रतसम्पूर्णहेतवे} %॥४६॥

\twolineshloka
{एवं कुरुष्व विप्रेन्द्र सर्वसिद्धिर्भविष्यति}
{नमस्कृत्वा तु ताः साध्वीर्विप्रः स्वगृहमाययौ} %॥४७॥

\twolineshloka
{तथा चकार तत्सर्वं व्रतं चैव यथाश्रुतम्}
{स्वकन्यकाद्वयस्यापि श्रावयामास तद्विधिम्} %॥४८॥

\twolineshloka
{तस्य श्रवणमात्रेण दर्शनात्पूजनस्य च}
{स्वरङ्गनोपमे कन्ये जाते तस्य प्रभावतः} %॥४९॥

\twolineshloka
{तदाप्रभृति विप्रस्य गृहे लक्ष्मीर्विवेश ह}
{नानामार्गैर्निमित्तैश्च लक्ष्मीवानिति सोऽभवत्} %॥५०॥

\twolineshloka
{कदाचिद् गच्छता राज्ञा विप्रसद्म पुरोऽध्वना}
{वातायने स्थिते कन्ये दृष्टे निरुपमे शुभे} %॥५१॥

\twolineshloka
{देहावयवसंस्थानैर्वस्तु यद्यच्च सुन्दरम्}
{त्रैलोक्ये भर्त्सयन्त्यौ ते पद्मचन्द्रादिकं च यत्} %॥५२॥

\twolineshloka
{राजा मोहं समापेदे तत्रैवावस्थितः क्षणम्}
{आमन्त्र्य ब्राह्मणं सद्यः प्रार्थयामास कन्यके} %॥५३॥

\twolineshloka
{विप्रोऽपि हृषितो भूत्वा प्रादाद्राज्ञे सुताद्वयम्}
{राजानं प्राप्य भर्तारं तेऽपि कन्ये मुदान्विते} %॥५४॥

\twolineshloka
{पुत्रपौत्रादिसम्पन्ने चक्रतुश्च स्वयं व्रतम्}
{व्रतमेतत्समाख्यातं मुने तव महोदयम्} %॥५५॥

\twolineshloka
{यस्य श्रवणमात्रेण सर्वान्कामानवाप्नुयात्}
{अनुष्ठानफलं तस्य किं वर्ण्य विधिनन्दन} %॥५६॥


{॥इति श्रीस्कन्दपुराणे ईश्वरसनत्कुमारसंवादे श्रावणमासमाहात्म्ये प्रकीर्णकनानाव्रत-रविवारव्रतादिकथनं नाम पञ्चमोऽध्यायः॥५॥}



\sect{षष्ठोऽध्यायः}

\uvacha{सनत्कुमार उवाच}
\twolineshloka
{रविवारस्य माहात्म्यं श्रुतं मे हर्षकारकम्}
{सोमवारस्य माहात्म्यं श्रावणे मासि मे वद} %॥१॥


\uvacha{ईश्वर उवाच}
\twolineshloka
{रविर्मे नयनं तस्य माहात्म्यमिदमुत्तमम्}
{उमासहितमन्नाम्नस्तस्य सोमस्य किं पुनः} %॥२॥

\twolineshloka
{माहात्म्यं वर्णनीयं मे यत्किञ्चिदपि ते ब्रुवे}
{सोमश्चन्द्रो विप्रराजः सोमः स्याद्यज्ञसाधनम्} %॥३॥

\twolineshloka
{निमित्तानि च तन्नाम्नः शृणु मत्तः समाहितः}
{मत्स्वरूपो यतो वारस्ततः सोम इति स्मृतः} %॥४॥

\twolineshloka
{प्रदाता सर्वराज्यस्य श्रेष्ठश्चैव ततो हि सः}
{समस्तराज्यफलदो व्रतकर्तुर्यतो हि सः} %॥५॥

\twolineshloka
{तस्य शृणु विधिं विप्र विस्तरात्कथयामि ते}
{द्वादशेष्वपि मासेषु सोमवारः प्रशस्यते} %॥६॥

\twolineshloka
{तावत्कर्तुमशक्तश्चेछ्रावणे मासि कारयेत्}
{अस्मिन्मासे व्रतं कृत्वा अब्दव्रतफलं लभेत्} %॥७॥

\twolineshloka
{श्रावणे शुक्लपक्षे तु प्रथमे सोमवासरे}
{सङ्कल्पयेद् व्रतं सम्यक् शिवो मे प्रीयतामिति} %॥८॥

\twolineshloka
{एवं चतुर्षु वारेषु भवेयुः पञ्च वा यदि}
{प्रातः सङ्कल्पयेत्तत्र नक्तं च शिवपूजनम्} %॥९॥

\twolineshloka
{उपचारैः षोडशभिः सायं च पूजयेच्छिवम्}
{शृणुयाच्च कथां दिव्यामेकाग्रकृतमानसः} %॥१०॥

\twolineshloka
{सोमवारव्रतस्यास्य कथ्यमानं निबोध मे}
{श्रावणे प्रथमे सोमे गृह्णीयाद् व्रतमुत्तमम्} %॥११॥

\twolineshloka
{सुस्नातश्च शुचिर्भूत्वा शुक्लाम्बरधरो नरः}
{कामक्रोधाद्यहङ्कारद्वेषपैशून्यवर्जितः} %॥१२॥

\twolineshloka
{आहरेच्छ्रेतपुष्पाणि मालतीमल्लिकादिकाः}
{अन्यैश्च विविधैः पुष्पैरभीष्टैरुपचारकैः} %॥१३॥

\onelineshloka
{पूजयेन्मूलमन्त्रेण त्र्यम्बकेण ततः परम्} %॥१४॥

\twolineshloka
{शर्वाय भवनाशाय महादेवाय धीमहि}
{उग्राय चोग्रनाथाय भवाय शशिमौलिने} %॥१५॥

\twolineshloka
{रुद्राय नीलकण्ठाय शिवाय भवहारिणे}
{एवं सम्पूज्य देवेशमुपचारैर्मनोहरैः} %॥१६॥


\threelineshloka
{यथाविभवसारेण तस्य पुण्यफलं शृणु}
{सोमवारे यजन्ते ये पार्वत्या सहितं शिवम्}
{ते लभन्त्यक्षयांल्लोकान्पुनरावृत्तिदुर्लभान्} %१७

\twolineshloka
{अत्र नक्तेन यत्पुण्यं कथयामि समासतः}
{सप्तजन्मार्जितं पापमभेद्यं देवदानवैः} %॥१८॥

\twolineshloka
{प्रणश्येन्नक्तभुक्तेन नात्र कार्या विचारणा}
{उपवासेन वा कुर्याद् व्रतमेतदनुत्तमम्} %॥१९॥

\twolineshloka
{पुत्रार्थी लभते पुत्रान्धनार्थी लभते धनम्}
{यं यं चिन्तयते कामं तं तं प्राप्नोति मानवः} %॥२०॥

\twolineshloka
{इह लोके चिरं स्थित्वा भुक्त्वा भोगान्यथेप्सितान्}
{विमानवरमारुह्य रुद्रलोके महीयते} %॥२१॥

\twolineshloka
{चलं चित्तं चलं वित्तं चलं जीवितमेव च}
{एवं ज्ञात्वा प्रयत्नेन व्रतस्योद्यापनं चरेत्} %॥२२॥

\twolineshloka
{उमामहेश्वरौ हैमौ राजते वृषभे स्थितौ}
{यथाशक्त्या प्रकर्तव्यौ वित्तशाठ्यं न कारयेत्} %॥२३॥

\twolineshloka
{मण्डलं लिङ्गतोभद्रं दिव्यं वै कारयेच्छुभम्}
{तत्र संस्थापयेत्कुम्भं श्वेतवस्त्रयुगान्वितम्} %॥२४॥

\twolineshloka
{ताम्रपात्रं वैणवं वा कुम्भस्योपरि विन्यसेत्}
{तस्योपरि न्यसेद्देवमुमया सहितं शिवम्} %॥२५॥

\twolineshloka
{श्रुतिस्मृतिपुराणोक्तैर्मन्त्रैः सम्पूजयेच्छिवम्}
{पुष्पमण्डपिका कार्या वितानं चैव शोभनम्} %॥२६॥

\twolineshloka
{रात्रौ जागरणं कार्यं गीतवादित्रनिःस्वनैः}
{स्वगृह्येोक्तविधानेन ततोऽग्निं स्थापयेद् बुधः} %॥२७॥

\twolineshloka
{ततो होमं च शर्वाद्यैरेकादशसुनामभिः}
{पालाशाभिः समिद्भिश्च हुनेदष्टाधिकं शतम्} %॥२८॥

\twolineshloka
{यवव्रीहितिलाद्यैश्च आप्यायस्वेति मन्त्रतः}
{बिल्वपत्रैस्त्र्यम्बकेण षड्वर्णेनापि वा पुनः} %॥२९॥

\twolineshloka
{पूर्णाहुतिं ततो हुत्वा कृत्वा स्विष्टकृतादिकम्}
{आचार्यं पूजयेत्पश्चाद् गां च तस्मै प्रदापयेत्} %॥३०॥

\twolineshloka
{ब्राह्मणान्भोजयेत्पश्चादेकादश सुशोभनान्}
{एकादश घटास्तेभ्यो वंशपात्रसमन्विताः} %॥३१॥

\twolineshloka
{सम्पूजितं ततो देवं देवोपकरणानि च}
{आचार्याय ततो दद्यात्प्रार्थयेत्तदनन्तरम्} %॥३२॥

\twolineshloka
{परिपूर्ण व्रतं मे स्याच्छिवो मे प्रीयतामिति}
{बन्धुभिः सह भुञ्जीत ततो हर्षपुरःसरम्} %॥३३॥

\twolineshloka
{अनेनैव विधानेन य इदं व्रतमाचरेत्}
{यं यं चिन्तयते कामं तं तं प्राप्नोति मानवः} %॥३४॥

\twolineshloka
{शिवलोके ततो गत्वा तस्मिँल्लोके महीयते}
{कृष्णेनाचरितं पूर्वं सोमवारव्रतं शुभम्} %॥३५॥

\twolineshloka
{नृपैः श्रेष्ठैस्तथा चीर्णमास्तिकैर्धर्मतत्परैः}
{इदं यः शृणुयान्नित्यं सोऽपि तत्फलमाप्नुयात्} %॥३६॥


{॥इति श्रीस्कन्दपुराणे ईश्वरसनत्कुमारसंवादे श्रावणमासमाहात्म्ये सोमवारव्रतकथनं नाम षष्ठोऽध्यायः॥६॥}



\sect{सप्तमोऽध्यायः}

\uvacha{ईश्वर उवाच}
\twolineshloka
{सनत्कुमार वक्ष्यामि भौमव्रतमनुत्तमम्}
{यस्यानुष्ठानमात्रेण अवैधव्यं प्रजायते} %॥१॥

\twolineshloka
{विवाहानन्तरं पञ्चवर्षाणि व्रतमाचरेत्}
{नामास्य मङ्गलागौरीव्रतं पापप्रणाशनम्} %॥२॥

\twolineshloka
{विवाहानन्तरं चाद्ये श्रावणे शुक्लपक्षके}
{प्रथमं भौमवारस्य व्रतमेतत्तु कारयेत्} %॥३॥

\twolineshloka
{पुष्पमण्डपका कार्यां कदलीस्तम्भमण्डिता}
{नानाविधैः फलैश्चैव पट्टकूलैश्च भूषयेत्} %॥४॥

\twolineshloka
{तत्र संस्थापयेद्देव्याः प्रतिमां स्वर्णनिर्मिताम्}
{अन्यधातुमयीं वापि स्वशक्त्या तत्र पूजयेत्} %॥५॥

\twolineshloka
{उपचारैः षोडशभिर्मङ्गलागौरिसंज्ञिताम्}
{दूर्वादलैः षोडशभिरपामार्गदलैस्तथा} %॥६॥

\twolineshloka
{तावत्सङ्ख्यैस्तण्डुलैश्च चणानां शकलैस्तथा}
{षोडशोन्मितवर्तीभिस्तावद्दीपांश्च दीपयेत्} %॥७॥

\twolineshloka
{दध्योदनं च नैवेद्यं तत्र भक्त्या प्रकल्पयेत्}
{समीपं स्थापयेद्देव्या दृषदं चोपलं तथा} %॥८॥

\twolineshloka
{एवं कृत्वा तु पञ्चाब्दं तत उद्यापनं चरेत्}
{मात्रे दद्याद्वायनं तु प्रकारं शृणु तस्य च} %॥९॥

\twolineshloka
{प्रतिमां मङ्गलागौर्याः सुवर्णपलनिर्मिताम्}
{तदर्थेन तदर्धेन शक्त्या वाप्यथ कारयेत्} %॥१०॥

\twolineshloka
{तण्डुलैः पूरिते भाण्डे शक्त्या स्वर्णादिनिर्मिते}
{संस्थाप्य परिधानीयं रमणीयां च कञ्चुकीम्} %॥११॥

\twolineshloka
{तयोरुपरि देव्यास्तु प्रतिमां स्थापयेत्ततः}
{समीपभागे संस्थाप्य दृषदं चोपलं तथा} %॥१२॥

\twolineshloka
{रौप्येण निर्मितं मात्रे एवं दद्यात्तु वायनम्}
{षोडशापि सुवासिन्यो भोजनीयाः प्रयत्नतः} %॥१३॥

\twolineshloka
{एवं कृते व्रते व सौभाग्यं सप्तजन्मसु}
{पुत्रपौत्रादिभिश्चैव रमते सम्पदा युता} %॥१४॥


\uvacha{सनत्कुमार उवाच}
\twolineshloka
{केनेदं व्रतमाचीर्णं कस्य जातं फलं पुरा}
{यथा स्यात्प्रत्ययः शम्भो कृपां कृत्वा तथा वद} %॥१५॥


\uvacha{ईश्वर उवाच}
\twolineshloka
{कुरुदेशे पुरा राजा श्रुतकीर्तिरिति श्रुतः}
{बभूव श्रुतसम्पन्नः कीर्तिमान्हतशात्रवः} %॥१६॥

\twolineshloka
{चतुःषष्टिकलाभिज्ञो धनुर्विद्याविशारदः}
{पुत्रादन्यच्छुभं सर्वं तस्य राज्ञो बभूव ह} %॥१७॥

\twolineshloka
{सन्तानविषयेऽथासौ बहुचिन्ताकुलोऽभवत्}
{देव्या आराधनं चक्रे जपध्यानपुरःसरम्} %॥१८॥

\twolineshloka
{क्रूरेण तपसा तस्य देवी तुष्टा बभूव ह}
{उवाच वचनं तस्मै वरं वरय सुव्रत} %॥१९॥


\uvacha{श्रुतकीर्तिरुवाच}
\twolineshloka
{यदि देवि प्रसन्नासि पुत्रं मे देहि शोभनम्}
{अन्यद्देवि त्वत्प्रसादान्न न्यूनं किञ्चिदस्ति मे} %॥२०॥

\twolineshloka
{इति तस्य वचः श्रुत्वा देवी प्राह शुचिस्मिता}
{दुर्लभं याचितं राजन्दास्ये तुभ्यं कृपावशात्} %॥२१॥

\twolineshloka
{परं शृणुष्व राजेन्द्र पुत्रश्चेद् गुणवत्तरः}
{ईप्सितश्चेत्षोडशाब्दं जीविष्यति न चाधिकम्} %॥२२॥

\twolineshloka
{रूपविद्याविहीनश्चेच्चिरञ्जीवी भविष्यति}
{इति देव्या वचः श्रुत्वा नृपश्चिन्तातुरोऽभवत्} %॥२३॥

\twolineshloka
{भार्यया सह सम्मन्त्र्य ययाचे गुणभूषितम्}
{सर्वलक्षणसम्पन्नं षोडशाब्दायुषं सुतम्} %॥२४॥

\twolineshloka
{आज्ञापयामास तदा देवी भक्तं नराधिपम्}
{आम्रवृक्षो मम द्वारे वर्तते नृपनन्दन} %॥२५॥

\twolineshloka
{तस्यैकं फलमादाय पत्यै देहि ममाज्ञया}
{भक्षणार्थं च सा धर्त्री गर्भं सद्यो न संशयः} %॥२६॥

\twolineshloka
{हृष्टो राजा तथा चक्रे पत्नी गर्भं च सा दधौ}
{दशमे मासि सुषुवे पुत्रं देवसुतोपमम्} %॥२७॥

\twolineshloka
{जातकर्मादिकं चक्रे हर्षशोकसमन्वितः}
{चिरायुरिति नामास्य पिता चक्रे शिवं भजन्} %॥२८॥

\twolineshloka
{प्राप्ते तु षोडशे वर्षे चिन्तामाप सभार्यकः}
{ततश्चक्रे विचारं स कष्टलब्धो ह्ययं सुतः} %॥२९॥

\twolineshloka
{स्वसमीपे कथं मृत्युर्द्रष्टव्यो दुःखदोऽस्य तु}
{काशीं प्रस्थापयामास मातुलेन समं विभुः} %॥३०॥

\twolineshloka
{भ्रातरं प्रार्थयामास राजपत्नी यशस्विनी}
{धृत्वा कार्पटिकं वेषं काशीं प्रति सुतं नय} %॥३१॥

\twolineshloka
{मृत्युञ्जयः प्रार्थितोऽस्ति पुत्रार्थं तु मया पुरा}
{प्रेषयिष्यामि विश्वेश यात्रार्थं च जगत्पतेः} %॥३२॥

\twolineshloka
{तस्मान्नेयः सुतो मेऽद्य पालनीयश्च यत्नतः}
{इति श्रुत्वा स्वसुर्वाक्यं स्वस्त्रीयेण समं ययौ} %॥३३॥

\twolineshloka
{दिनानि कतिचिद् गच्छन्नानन्दनगरं ययौ}
{तत्र राजा वीरसेनो नाम्ना सर्वसमृद्धिमान्} %॥३४॥

\twolineshloka
{मङ्गलागौरी सर्वलक्षणसंयुता}
{वयोमध्यगता रम्या रूपलावण्यशालिनी} %॥३५॥

\twolineshloka
{उपमानानि सर्वाणि तुच्छीकृत्योदयं गता}
{नगरोपवने रम्ये सखीभिः परिवारिता} %॥३६॥

\twolineshloka
{ततस्तावपि सम्प्राप्तौ चिरायुर्मातुलश्च सः}
{विश्रान्तिं प्रापतुस्तत्र तासां दर्शनलालसौ} %॥३७॥

\twolineshloka
{क्रीडन्तीनां विनोदेन कुपिता तत्र काचन}
{उवाच राजतनयां सा रण्डेत्यति दुर्वचः} %॥३८॥

\twolineshloka
{श्रुत्वा तदशुभं वाक्यमुवाच नृपनन्दिनी}
{अयोग्यं भाषसे त्वं किं मत्कुले नैव तद्विधा} %॥३९॥

\twolineshloka
{प्रसादान्मङ्गलागौर्यांस्तद्व्रतस्य प्रभावतः}
{मत्करादक्षता यस्य प्रपतिष्यन्ति मूर्धनि} %॥४०॥

\twolineshloka
{विवाहे स चिरायुः स्यादल्पायुरपि चेत्सखि}
{ततः समस्तास्ताः कन्याः स्वं स्वं वेश्म ययुस्तदा} %॥४१॥

\twolineshloka
{स एव दिवसो राजकन्यायाः पाणिपीडने}
{राज्ञो बाह्लीकदेशस्य दृढधर्माभिधस्य वै} %॥४२॥

\twolineshloka
{सुकेतुनाने पुत्राय दातुं सा निश्चिताभवत्}
{स सुकेतुरविद्यश्च कुरूपो बधिरस्तथा} %॥४३॥

\twolineshloka
{ततस्ते मन्त्रयामासुर्नेयोऽन्योऽद्य वरः परः}
{अथ सिद्धे विवाहे च सुकेतुस्तत्र गच्छतु} %॥४४॥

\twolineshloka
{ततश्चिरायुषं गत्वा याचिरे मातुलं प्रति}
{देयोऽस्मभ्यमयं बालः कार्यसिद्धिर्हि नो भवेत्} %॥४५॥

\twolineshloka
{परोपकारतुल्यो हि धर्मो नास्त्यपरो भुवि}
{मातुलस्तद्वचः श्रुत्वा ह्यन्तर्हृष्टमना अभूत्} %॥४६॥

\twolineshloka
{पूर्व श्रुतं चोपवने कन्यावाक्यमनेन यत्}
{एकवारं तथाप्याह युष्माभिर्याच्यते कथम्} %॥४७॥

\twolineshloka
{वस्त्रालङ्करणादीनि याच्यं कार्यस्य साधने}
{न वरो याच्यते क्वापि दीयते गौरवाद्धि वः} %॥४८॥

\twolineshloka
{विवाहं साधयामासुनित्वा तत्र चिरायुषम्}
{सप्तपद्यादिके जाते रात्रौ गौरीहरान्तिके} %॥४९॥

\twolineshloka
{सहासौ मङ्गला गौर्या सुप्तो हर्षसमन्वितः}
{तदहिन षोडशाब्दानि समाप्तानि चिरायुषः} %॥५०॥

\twolineshloka
{निशीथे सर्परूपेण कालस्तत्र समीयिवान्}
{तदन्तरे भूपसुता जागृता दैवयोगतः} %॥५१॥

\twolineshloka
{सा ददर्श महासर्प चकम्पे भयविह्वला}
{धैर्यं कृत्वा तदा वाला पूजयामास सोरगम्} %॥५२॥

\twolineshloka
{उपचारैः षोडशभिर्दुग्धं पातुं ददौ बहु}
{प्रार्थयामास तं सर्प दीनवाण्या च तुष्टुवे} %॥५३॥

\twolineshloka
{ययाचे मङ्गलागौरी करिष्ये व्रतमुत्तमम्}
{जीव्यान्मे पतिरेतस्माच्चिरं जीवेत्तथा कुरु} %॥५४॥

\twolineshloka
{एतस्मिन्नन्तरे सर्पः करके प्रविवेश ह}
{कञ्चुक्या स्वीयया सा तु चक्रे तन्मुखबन्धनम्} %॥५५॥

\twolineshloka
{एतस्मिन्नन्तरे भर्ता अङ्गमोटनपूर्वकम्}
{जागृतश्चाब्रवीद्भार्यां क्षुधा मां बाधते प्रिये} %॥५६॥

\twolineshloka
{मातुः सकाशं गत्वा सा आनयामास पायसम्}
{लड्डुकादि च तद्दत्तं बुभुजे प्रीतमानसः} %॥५७॥

\twolineshloka
{हस्तक्षालनकाले तु तद्धस्तान्मुद्रिकापतत्}
{ताम्बूलं भक्षयित्वा तु प्रसुप्तः पुनरेव सः} %॥५८॥

\twolineshloka
{ततः सा करके त्यक्तुमगच्छत्तु विधेर्गतिः}
{हारकान्तिं बहिर्दृष्ट्वा स्फुरन्तीं विस्मयं ययौ} %॥५९॥

\twolineshloka
{दृष्ट्वा घटस्थं तं हारं स्वकण्ठे च दधार सा}
{किञ्चिन्निशावशेषे तु मातुलस्तं निनाय सः} %॥६०॥

\twolineshloka
{ततस्ते वरपक्षीयाः सुकेतुं तत्र चानयन्}
{दृष्ट्वा तं मङ्गलागौरी उवाचायं न मे पतिः} %॥६१॥

\twolineshloka
{तामूचुस्ते ततः सर्वे किमिदं भाषसे शुभे}
{परिचायकमस्तीह किञ्चित्ते तद्वदस्व नः} %॥६२॥


\uvacha{मङ्गलागौर्युवाच}
\twolineshloka
{मे दत्तं येन रात्रौ च नवरत्नाङ्गुलीयकम्}
{अस्याङ्गुलौ तन्निक्षिप्य पश्यध्वं परिचायकम्} %॥६३॥

\twolineshloka
{पत्या दत्तोऽस्ति में हारो रात्रौ तद्रत्नसञ्चयः}
{कीदृशोऽनेन वाच्योऽसौ प्रतिवीरपरान्वितम्} %॥६४॥

\twolineshloka
{किञ्चाम्रसेचने रात्रौ तत्पदं कुङ्कुमान्वितम्}
{ऊरौ मे वर्तते तच्च सर्वे पश्यन्तु माचिरम्} %॥६५॥

\twolineshloka
{किञ्च रात्रौ भाषणादि भक्षणादि च यत्कृतम्}
{तदनेन च वक्तव्यं तदा स्यान्मे पतिः स्वयम्} %॥६६॥

\twolineshloka
{एवं श्रुत्वा तु तद्वाक्यं साधु साध्विति चाब्रुवन्}
{एकस्यापि न योगोऽभूत्तदा सर्वैर्निषेधितः} %॥६७॥

\twolineshloka
{तदा ते वरपक्षीया जग्मुः सर्वे यथागतम्}
{जनको मङ्गलागौर्याः श्रुतकीर्तिः कुलोद्वहः} %॥६८॥

\twolineshloka
{अन्नपानादिकं सत्रं चकार सुमहामनाः}
{वरपक्षस्य वृत्तान्तः श्रुतः कर्णापकर्णतः} %॥६९॥

\twolineshloka
{स्वरूपाय कुरूपत्वादानीतः कश्चनादृतः}
{स्थापयामास सौधे तु कन्यां जवनिकावृताम्} %॥७०॥

\twolineshloka
{एवं गते हायने तु यात्रां कृत्वा समातुलः}
{चिरायुः प्रययौ तत्र किं जातमवलोकितुम्} %॥७१॥

\twolineshloka
{तं सा जालान्तराद् दृष्ट्वा लोकोत्तरमुदान्विता}
{पितरौ कथयामास मम भर्ता समागतः} %॥७२॥

\twolineshloka
{सुहृद्गणं समाहूय पूर्वोक्तं परिचायकम्}
{दृष्ट्वा सर्वमपि हास्मै ददौ कन्यां शुचिस्मिताम्} %॥७३॥

\twolineshloka
{शिष्टैः परिणयोत्साहं कारयामास भूपतिः}
{वस्त्राण्याभरणादीनि सेनामश्वान्गजान् रथान्} %॥७४॥

\twolineshloka
{प्रस्थापयामास नृपो दत्त्वान्यदपि भूरिशः}
{पल्या सह चिरायुः स मातुलेन समन्वितः} %॥७५॥

\twolineshloka
{स्वपुरं सेनया सार्धं जगाम कुलनन्दनः}
{श्रुत्वा जनमुखात्तं तु ह्यागतं पितरावुभौ} %॥७६॥

\twolineshloka
{विश्वासं लेभतुर्नैव स्यात्कथं दैवमन्यथा}
{एतस्मिन्नन्तरे प्राप्तः पित्रोरन्तिकमेव सः} %॥७७॥

\twolineshloka
{पपात पादयोर्भक्त्या पित्रोः स्नेहपरिप्लुतः}
{मूर्ध्न्यवनाय तं पुत्रं परमां मुदमापतुः} %॥७८॥

\twolineshloka
{स्नुषापि मङ्गलागौरी श्वशुरौ प्रणनाम सा}
{अङ्के निवेश्य तां श्वश्रूः पप्रच्छोदन्तमञ्जसा} %॥७९॥

\twolineshloka
{स्नुषापि मङ्गला गौर्या व्रतमाहात्म्यमुत्तमम्}
{कथयामास तत्सर्वं यथावृत्तं महामुने} %॥८०॥

\twolineshloka
{इत्येतत्कथितं तुभ्यं मङ्गलागौरिकाव्रतम्}
{य एतच्छृणुयात्कश्चिद्यश्चापि परिकीर्तयेत्} %॥८१॥

\onelineshloka
{मनोरथास्तस्य सर्वे सिध्यन्त्यत्र न संशयः} %॥८२॥


\uvacha{सूत उवाच}
\twolineshloka
{सनत्कुमारमित्येवं कथयामास धूर्जटिः}
{स चानन्दं परं लेभे श्रुत्वा कार्यकरं व्रतम्} %॥८३॥


{॥इति श्रीस्कन्दपुराणे ईश्वरसनत्कुमारसंवादे श्रावणमासमाहात्म्ये मङ्गलागौरीव्रतकथनं नाम सप्तमोऽध्यायः॥७॥}



\sect{अष्टमोऽध्यायः}

\uvacha{ईश्वर उवाच}
\twolineshloka
{बुधगुर्वोरथो वक्ष्ये व्रतं पापप्रणाशनम्}
{यत्कृत्वा श्रद्धया मर्त्यः परां सिद्धिमवाप्नुयात्} %॥१॥

\twolineshloka
{प्रजापतिः शीतरश्मिं द्विजराज्येऽभ्यषेचयत्}
{स कदाचिद् गुरोर्भार्यां तारानाम्नीं ददर्श ह} %॥२॥

\twolineshloka
{रूपयौवनसम्पन्नां लावण्यमदगर्विताम्}
{मोहितो रूपसम्पत्या कामबाणवशं गतः} %॥३॥

\twolineshloka
{स्वगृहे स्थापयित्वा तु बलात्स बुभुजे च ताम्}
{एवं बहुतिथे काले गते पुत्रो बभूव ह} %॥४॥

\twolineshloka
{बुधो बुधो रूपशाली सर्वलक्षणसंयुतः}
{अन्वेषयन्नगुरुः पत्नीं ज्ञातवाञ्छशिसद्यनि} %॥५॥

\twolineshloka
{ययाचे देहि मे भार्यां त्वं कथं गुरुतल्पगः}
{गुरुतल्पकृतात्पापान्निष्कृतिस्ते कथं भवेत्} %॥६॥

\twolineshloka
{महापातकसंयोगे कथं ते बुद्धिरादृता}
{गुप्तमेव प्रयच्छेमां गुरुभार्यां मम प्रियाम्} %॥७॥

\twolineshloka
{रहसि प्रायश्चित्तं च कृत्वा निष्कल्मषो भव}
{नोचेदिन्द्रसमीपे ते आगः सङ्कथयाम्यहम्} %॥८॥

\twolineshloka
{इत्येवं बहुधोक्तोऽपि न ददौ तां कलङ्कितः}
{तदा देवसभां गत्वा कथयामास गीष्पतिः} %॥९॥

\twolineshloka
{चन्द्रेण मे ह्यपहृता भार्या तां न ददाति सः}
{देवराजोऽसि शक्र त्वं दापनीया त्वयाज्ञया} %॥१०॥

\twolineshloka
{नोचेत्ते तत्कृतं पापं सङ्क्रमिष्यत्यसंशयम्}
{राजा राष्ट्रकृतं पापं भुङ्गे शास्त्रविनिर्णयात्} %॥११॥

\twolineshloka
{दुर्बलस्य बलं राजा पुराणे त्विति भण्यते}
{इति श्रुत्वा गुरोर्वाक्यं चन्द्रमाहूय वासवः} %॥१२॥

\twolineshloka
{आज्ञापयामास रुषा देहि भार्यां गुरोर्विधो}
{अन्यदाराभिगमनं केवलं पापसंज्ञितम्} %॥१३॥

\twolineshloka
{गुरुदाराभिगमनं महापातकसंज्ञितम्}
{तस्माच्चन्द्र गुरोर्भार्यां देहि त्वमविचारयन्} %॥१४॥

\twolineshloka
{देवेन्द्रवचनं श्रुत्वा निशापतिरथाब्रवीत्}
{दास्ये त्वदाज्ञया भार्यां पुत्रं नैव ददाम्यहम्} %॥१५॥

\twolineshloka
{मत्सकाशात्सुतो जातो वैभवयुग्यतः}
{गीष्पतिस्त्वाह मत्तोऽभूत्ततः संशयिताः सुराः} %॥१६॥

\twolineshloka
{ततस्ते निर्णयं चक्रुर्माता जानाति चाङ्गजम्}
{पप्रच्छुस्ते तदा तारां केनायं गर्भ आहितः} %॥१७॥

\twolineshloka
{सत्यं वदस्व कल्याणि न मिथ्या वक्तुमर्हसि}
{तदा लज्जान्विता तारा औरसोऽयं विधोः सुतः} %॥१८॥

\twolineshloka
{गीष्यतेः क्षेत्रजश्चातो योग्यः स्यात्तस्य दीयताम्}
{शास्त्रतस्ते विचार्याथ ददुश्चन्द्राय तं बुधम्} %॥१९॥

\twolineshloka
{तदा खिन्नं गुरुं दृष्ट्वा ददुर्देवा वरं तयोः}
{गच्छस्व त्वं चन्द्र गृहं तवाप्यस्ति सुतो ह्ययम्} %॥२०॥

\twolineshloka
{चन्द्रस्य गीष्पतेश्चायं ग्रहत्वं यात्वसौ सुतः}
{अन्यच्चापि सुराचार्य गृहाणेमं वरं शुभम्} %॥२१॥

\twolineshloka
{यः करिष्यति मेधावी मिलित्वा युवयोर्व्रतम्}
{तस्य स्यात्सकला सिद्धिः सत्यं सत्यं न संशयः} %॥२२॥

\twolineshloka
{श्रावणे मासि सम्प्राप्ते शङ्करस्य महाप्रिये}
{बुधगुर्वोर्वासरयोर्ये करिष्यन्ति पूजनम्} %॥२३॥

\twolineshloka
{नैवेद्यं दधिभक्तेन साधने मूलकं भवेत्}
{युवयोर्मूर्तिमालिख्य स्थानभेदात्फलं लभेत्} %॥२४॥

\twolineshloka
{बालान्दोलोपरिस्थाने लिखित्वा पूजयेद्यदि}
{स पुत्रं लभते दीर्घायुषं सर्वगुणान्वितम्} %॥२५॥

\twolineshloka
{कोशागारे लिखित्वा तु पूजयेद्यदि मानवः}
{तस्य कोशा विवर्धन्ते क्षीयन्ते न कदाचन} %॥२६॥

\twolineshloka
{पाकागारे पाकवृद्धिर्देवागारे तु तत्कृपा}
{शय्यागारे पूजने तु स्त्रीवियोगो न कर्हिचित्} %॥२७॥

\twolineshloka
{धान्यागारे धान्यवृद्धिरेवं तत्तत्फलं लभेत्}
{सप्तवर्षाणि कृत्वैवं तत उद्यापनं चरेत्} %॥२८॥

\twolineshloka
{अधिवास्याह्नि पूर्वस्मिन् रात्रौ जागरणं चरेत्}
{सुवर्णप्रतिमां कृत्वा पूजयित्वा यथाविधि} %॥२९॥

\twolineshloka
{उपचारैः षोडशभिस्ततो होमं समाचरेत्}
{तिलैराज्येन चरुणा तथैव च समिद्धजैः} %॥३०॥

\twolineshloka
{अपामार्गाश्वत्थमयैस्ततः पूर्णाहुतिं चरेत्}
{स्वस्त्रीयमातुलौ चैव भोजनीयौ प्रयत्नतः} %॥३१॥

\twolineshloka
{ब्राह्मणान्भोजयेदन्यान्भुञ्जीत स्वयमेव च}
{एवं कृते सप्तवर्षं सर्वान्कामानवाप्नुयात्} %॥३२॥

\twolineshloka
{विद्याकामनया कुर्याद्वेदशास्त्रार्थविद्भवेत्}
{बुधस्तु बुधतां दद्याद् गुरुस्तु गुरुतां तथा} %॥३३॥


\uvacha{सनत्कुमार उवाच}
\twolineshloka
{भगवन्यत्त्वया प्रोक्तं भोज्यौ स्वस्त्रीयमातुलौ}
{एतन्निमित्तं कथय यदि वक्तुं क्षमं भवेत्} %॥३४॥


\uvacha{ईश्वर उवाच}
\twolineshloka
{पुरा कौचिद् द्विजन्मानौ दीनौ स्वस्त्रीयमातुलौ}
{दरिद्रौ पर्यटन्तौ तावुदरार्थे कृतश्रमौ} %॥३५॥

\twolineshloka
{कस्मिंश्चिन्नगरे रम्ये गतौ धान्यं प्रयाचितुम्}
{गृहे गृहेऽपश्यतां तौ श्रावणे मासि तद्व्रतम्} %॥३६॥

\twolineshloka
{तत्तद्वारे व्रतं तत्तद् बुधगुर्वोर्न कुत्रचित्}
{अन्योन्यं तौ तदा तत्र विचारं चक्रतुश्चिरात्} %॥३७॥

\twolineshloka
{वासराणां तु सर्वेषां व्रतं सर्वत्र दृश्यते}
{बुधगुर्वोर्विना तस्मादावाभ्यां तद्व्रतं शुभम्} %॥३८॥

\twolineshloka
{अनुच्छिष्टं यतश्चास्ति तस्मात्कर्तव्यमादरात्}
{विध्यज्ञानात्परं तस्य संशयं प्रापतुः पुनः} %॥३९॥

\twolineshloka
{तावत्तस्यां निशायां तु स्वप्नोऽभूद्विधिदर्शनः}
{तथा तौ चक्रतुः पश्चात्परां सम्पदमापतुः} %॥४०॥

\twolineshloka
{प्रत्यहं वृद्धिगा चाभूत्सम्पत्तिः सर्वगोचरा}
{एवं कृत्वा सप्तवर्षं पुत्रपौत्रादिसंयुतौ} %॥४१॥

\twolineshloka
{साक्षाद्भूतौ बुधगुरू वरं च ददतुस्तयोः}
{आवाभ्यामावयोर्यस्माद् व्रतमेतत्प्रवर्तितम्} %॥४२॥

\twolineshloka
{इति चारभ्य तस्माद्यः करिष्यति शुभं व्रतम्}
{स्वस्त्रीयमातुलौ तेन भोजनीयौ प्रयत्नतः} %॥४३॥

\twolineshloka
{एतद्व्रतप्रभावेण सर्वसिद्धिः परा भवेत्}
{अन्ते चास्मल्लोकवासो यावच्चन्द्रदिवाकरौ} %॥४४॥


{॥इति श्रीस्कन्दपुराणे ईश्वरसनत्कुमारसंवादे श्रावणमासमाहात्म्ये बुधगुरुव्रतकथनं नामाष्टमोऽध्यायः॥८॥}



\sect{नवमोऽध्यायः}

\uvacha{ईश्वर उवाच}
\twolineshloka
{अतः परं प्रवक्ष्यामि शुक्रवारकथानकम्}
{यच्छ्रुत्वा श्रद्धया मर्त्यो मुच्यते सर्वसङ्कटात्} %॥१॥

\twolineshloka
{अत्रैवोदाहरन्तीममितिहासं पुरातनम्}
{सुशीलो नाम राजासीत्पाण्ड्यवंशसमुद्भवः} %॥२॥

\twolineshloka
{बहुप्रयत्नशीलोऽपि अपत्यं नैव चाप्तवान्}
{तस्य भार्या सुकेशीति नाम्ना सर्वगुणान्विता} %॥३॥

\twolineshloka
{अपत्यं न यदा लेभे महाचिन्तामवाप सा}
{स्त्रीस्वभावात्तदा वस्त्रखण्डानि प्रतिमासिके} %॥४॥

\twolineshloka
{बध्वोदरं महच्चक्रे महासाहसमानसा}
{अन्वेषयद्गर्भिणीं सा स्वप्रसूत्यनुसारिणीम्} %॥५॥

\twolineshloka
{भाविना दैवयोगेन गृहिणी तत्पुरोधसः}
{गर्भिण्यासीत्तदा राज्ञः पत्नी कपटकारिणी} %॥६॥

\twolineshloka
{प्रसूतिकारिणी काञ्चित्तत्कार्ये सा न्ययोजयत्}
{दत्त्वा बहुधनं तस्यै सूतिकायै रहो गता} %॥७॥

\twolineshloka
{राजा चक्रे पुंसवनं तथैवानवलोभनम्}
{सीमन्तोन्नयने काले महाहर्षसमन्वितः} %॥८॥

\twolineshloka
{तस्याः प्रसूतिसमयं श्रुत्वा सापि तथाकरोत्}
{आद्यगर्भवती यस्मात्सा पुरोधः कुटुम्बिनी} %॥९॥

\twolineshloka
{अज्ञा प्रसूतिकाकृत्ये सूतिकावचने स्थिता}
{तां सूतिका वञ्चयन्ती चक्रे तन्नेत्रबन्धनम्} %॥१०॥

\twolineshloka
{प्रेषयामास तं पुत्रं सा राजमहिषीं प्रति}
{कस्यचिद्धस्ततः शीघ्रमज्ञातमपि केनचित्} %॥११॥

\twolineshloka
{राज्ञी गृहीत्वा तं पुत्रं प्रसूतास्मीत्यघोषयत्}
{पुरोधः स्त्रीनेत्रबन्धं मोक्षयामास सूतिका} %॥१२॥

\twolineshloka
{सहानीतं मांसपिण्डं तस्यै प्रादर्शयच्च सा}
{विस्मयं चैव खेदं च स्वयं चक्रे तदग्रतः} %॥१३॥

\twolineshloka
{किमरिष्टमिदं जातं पत्या कार्य च शान्तिकम्}
{सन्ततिर्नास्ति चेन्मास्तु स्वदिष्ट्या जीवितासि भोः} %॥१४॥

\onelineshloka
{परं संशयिता सासीत्प्रसवस्पर्शचिन्तनात्} %॥१५॥


\uvacha{ईश्वर उवाच}
\twolineshloka
{राजा श्रुत्वा पुत्रजन्म जातकर्माद्यकारयत्}
{गजानश्वान् रथांश्चैव ब्राह्मणेभ्यो ददौ नृपः} %॥१६॥


\threelineshloka
{बद्धान् कारागृहे सर्वान्मोचयामास हर्षितः}
{सूतकान्ते नामकर्म संस्कारान्सर्वतोऽकरोत्}
{चक्रे प्रियव्रत इति नाम पुत्रस्य भूमिपः} %१७

\twolineshloka
{श्रावणे मासि सम्प्राप्ते पुरोधोदयिता सती}
{जीवन्तिकां शुक्रवारे पूजयामास भक्तितः} %॥१८॥

\twolineshloka
{कुड्ये विलिख्य तन्मूर्ति बहुबालसमन्विताम्}
{पुष्पमालिकया पूज्य पञ्चदीपैरदीपयत्} %॥१९॥

\twolineshloka
{गोधूमपिष्टसम्भूतैस्तानभक्षयत स्वयम्}
{अक्षतांश्चैव चिक्षेप यत्र मे बालको भवेत्} %॥२०॥

\twolineshloka
{तत्र त्वया रक्षणीयो जीवन्ति करुणार्णवे}
{इति प्रार्थ्य कथां श्रुत्वा नमश्चक्रे यथाविधि} %॥२१॥

\twolineshloka
{जीवन्तिकाप्रसादेन दीर्घायुर्बालकोऽभवत्}
{ररक्ष तमहोरात्रं देवी तन्मातृगौरवात्} %॥२२॥

\twolineshloka
{एवं काले गते राजा कालधर्ममुपेयिवान्}
{पितृभक्तोऽथ तत्पुत्रश्चक्रे तत्साम्परायिकम्} %॥२३॥

\twolineshloka
{प्रियव्रतोऽभिषिक्तोऽभूद्राज्ये मन्त्रिपुरोहितैः}
{पालयित्वा प्रजा राज्यं भुक्त्वा स कतिचित्समाः} %॥२४॥

\twolineshloka
{पित्रर्णस्य विमोक्षाय गयां गन्तुं प्रचक्रमे}
{राज्यभारममात्येषु स्थाप्य वृद्धेषु भक्तितः} %॥२५॥

\twolineshloka
{राजभावं परित्यज्य वेषं कार्पटिकं दधे}
{मार्गमध्ये क्वचित्पुर्यां कस्यचिद् गृहमेधिनः} %॥२६॥

\twolineshloka
{चक्रे वासं गृहे तस्य प्रसूता गृहिणी त्वभूत्}
{पुरा षष्ठ्या पञ्चमेऽह्नि तत्पुत्राः पञ्च मारिताः} %॥२७॥

\twolineshloka
{तदापि पञ्चमदिनेऽप्यासीत्तत्र नृपो गतः}
{रात्रौ सुप्ते नृपे षष्ठी बालं नेतुं समागता} %॥२८॥

\twolineshloka
{जीवन्त्या वारिता सा तु नृपमुल्लङ्घ्य मा व्रज}
{षष्ठी निषेधाज्जीवन्त्याः सा जगाम यथागता} %॥२९॥

\twolineshloka
{जीवितं पञ्चमदिने बालं लेभे गृहाधिपः}
{एतत्प्रभावः प्रायोऽयं प्रार्थयामास तं नृपम्} %॥३०॥

\twolineshloka
{राजन्नद्यतने चाह्नि तव वासोऽस्तु मे गृहे}
{तव प्रसादान्मे बालः षष्ठोऽयं जीवितः प्रभो} %॥३१॥

\twolineshloka
{एवं सम्प्रार्थितस्तेन उवाच करुणानिधिः}
{ततो गतो गयां राजा प्रवृत्तः पिण्डपातने} %॥३२॥

\twolineshloka
{विष्णुपादे तत्र किञ्चिद्दत्त्वाश्चर्यमभूत्तदा}
{पिण्डस्य ग्रहणार्थं हि निःसृतं तु करद्वयम्} %॥३३॥

\twolineshloka
{परं विस्मयमापन्नः संशयं प्राप भूपतिः}
{ब्राह्मणानुमतः पश्चात्पिण्डं विष्णुपदे ददौ} %॥३४॥

\twolineshloka
{पप्रच्छ ब्राह्मणं कञ्चिज्ज्ञानिनं सत्यवादिनम्}
{स चाह ब्राह्मणस्तस्मै पितृद्वयकराविमौ} %॥३५॥

\twolineshloka
{किमिदं तद् गृहे गत्वा मात्रे पृच्छ वदिष्यति}
{ततश्चिन्तापरो दुःखी हृदि नाना व्यचारयत्} %॥३६॥

\twolineshloka
{यात्रां कृत्वा तत्र यातो यत्रासौ जीवितः शिशुः}
{तदापि पञ्चमदिनमासीत्सैव प्रसूतिका} %॥३७॥

\twolineshloka
{द्वितीयोऽप्यभवत्पुत्रो रात्रौ षष्ठी समाययौ}
{पुनश्च जीवन्तिकया निषिद्धा साब्रवीच्च ताम्} %॥३८॥

\twolineshloka
{एतस्यावश्यकं किं ते एतन्मात्रा च किं व्रतम्}
{क्रियते हि यतस्त्वं च एनं रक्षस्यहर्निशम्} %॥३९॥

\twolineshloka
{षष्ठीवाक्यमिति श्रुत्वा जीवन्ती प्राह सुस्मिता}
{तन्निमित्तं निशि द्रष्टुं जाग्रदासीन्मृषा स्वपन्} %॥४०॥

\twolineshloka
{संवादमुभयो राजा शुश्राव सकलं तदा}
{श्रावणे भृगुवारे तु एतन्माता ममार्चने} %॥४१॥

\twolineshloka
{व्रतस्य नियमं सर्वं कुरुते तं वदामि ते}
{परिधत्ते न वसनं हरितं कञ्चुकीं तथा} %॥४२॥

\twolineshloka
{न धारयति तद्वर्णं काचकङ्कणकं करे}
{कदापि नोल्लङ्घयति तन्दुलक्षालनोदकम्} %॥४३॥

\twolineshloka
{नैव गच्छत्यधस्ताच्च हरित्पल्लवमण्डपम्}
{कृकलस्य च शाकं सा नाश्नाति हरिवर्णतः} %॥४४॥

\twolineshloka
{सर्वमेव मम प्रीत्यै मारयिष्यामि मा सुतम्}
{श्रुत्वा सर्वं नृपः प्रातर्जगाम स्वपुरं प्रति} %॥४५॥

\twolineshloka
{प्रत्युद्गता नागरिका दैशिकाः सर्व एव हि}
{पप्रच्छ मातरं राजा त्वया जीवन्तिकाव्रतम्} %॥४६॥

\twolineshloka
{क्रियते तु कथं मातर्न वेद्मीति च साब्रवीत्}
{साद्गुण्यार्थं तु यात्राया ब्राह्मणांश्च सुवासिनीः} %॥४७॥

\twolineshloka
{इच्छन्नृपो भोजयितुं व्रतं चापि परीक्षितुम्}
{सुवासिनीभ्यो वस्त्राणि कञ्चुक्यः कङ्कणानि च} %॥४८॥

\twolineshloka
{आगन्तव्यं भोजनार्थं सर्वाभी राजसद्मनि}
{ततः पुरोधसः पत्नी तत्र दूतमुवाच ह} %॥४९॥

\twolineshloka
{हरिद्वर्णं मया किञ्चिद् गृह्यते न कदाचन}
{दूतो निवेदयामास राज्ञे तस्याः प्रभाषितम्} %॥५०॥

\twolineshloka
{राजा सर्वं रक्तवर्णं तस्यै सम्प्रेषयच्छुभम्}
{अङ्गीकृत्य च तत्सर्वं सापि राजगृहं ययौ} %॥५१॥

\twolineshloka
{पूर्वद्वारे तण्डुलानां दृष्ट्वा क्षालनजं जलम्}
{मण्डपं च हरिद्वर्णं दृष्ट्वान्यद्द्वारतो ययौ} %॥५२॥

\twolineshloka
{राजा पुरोधसः पत्नीं नत्वा पप्रच्छ चाखिलम्}
{निमित्तं नियमस्यास्य सा प्रोवाच व्रतं भृगोः} %॥५३॥

\twolineshloka
{तं दृष्ट्वा तु तदात्यन्तं प्रस्तुतौ तत्पयोधरौ}
{सिषिञ्चतुस्तं राजानं धाराभिः सर्वतः स्तनौ} %॥५४॥

\twolineshloka
{गयायां करयुग्मेन देव्योः संवादतस्तथा}
{स्तनयोः प्रस्रवाच्चैव राजा प्रत्ययमाप सः} %॥५५॥

\twolineshloka
{पालिकां मातरं गत्वा पप्रच्छ विनयान्वितः}
{मा भैर्मातब्रूहि सत्यं वृत्तान्तं मम जन्मनः} %॥५६॥

\twolineshloka
{श्रुत्वा सुकेशिनी प्राह याथातथ्येन सर्वशः}
{हृष्टो भूत्वा नमश्चक्रे पितरौ स्वस्य जन्मदौ} %॥५७॥

\twolineshloka
{सम्पत्त्या वर्धयामास तौ परां मुदमापतुः}
{एकस्मिन्दिवसे राजा जीवन्तीं प्रार्थयन्निशि} %॥५८॥

\twolineshloka
{जीवन्त्ययं जन्मदो मे गयायां च करौ कथम्}
{तदा स्वप्नगता देवी प्राह संशयनाशकम्} %॥५९॥

\twolineshloka
{मया त्वत्प्रत्ययार्थेयं कृता माया न संशयः}
{एतत्ते सर्वमाख्यातं श्रावणे भृगुवासरे} %॥६०॥

\onelineshloka
{एतद् व्रतमनुष्ठाय सर्वान्कामानवाप्नुयात्} %॥६१॥


{॥इति श्रीस्कन्दपुराणे ईश्वरसनत्कुमारसंवादे श्रावणमासमाहात्म्ये शुक्रवारजीवन्तिकाव्रतकथनं नाम नवमोऽध्यायः॥९॥}




\sect{दशमोऽध्यायः}
\twolineshloka
{अतः ईश्वर उवाच परं प्रवक्ष्यामि मन्दवारविधिं तव}
{सनत्कुमार यत्कृत्वा मन्दत्वं नैव जायते} %॥१॥

\twolineshloka
{श्रावणे मासि देवानां त्रयाणां पूजनं शनौ}
{नृसिंहस्य शनेश्चैव अञ्जनीनन्दनस्य च} %॥२॥

\twolineshloka
{कुड्ये स्तम्भेऽथवालिख्य नृसिंहप्रतिमां शुभाम्}
{हरिद्रायुक्चन्दनेन लक्ष्म्या सह जगत्पतिम्} %॥३॥

\twolineshloka
{सम्पूज्य नीलपुष्पैश्च रक्तैः पीतैश्च शोभनैः}
{नैवेद्यं खिच्चडीसंज्ञं शाकं कुञ्जरसंज्ञितम्} %॥४॥

\twolineshloka
{स्वयं चैव तदश्नीयाद् ब्राह्मणांश्चैव भोजयेत्}
{तिलतैलं घृतस्नानं नृसिंहस्य प्रियं भवेत्} %॥५॥

\twolineshloka
{शनैश्चरदिने तैलं प्रशस्तं सर्वकर्मसु}
{अभ्यज्या ब्राह्मणास्तद्वत्सुवासिन्यस्तु तैलतः} %॥६॥

\twolineshloka
{स्वयमभ्यज्य च स्नायात्कुटुम्बसहितः शनौ}
{माषान्नं च प्रकर्तव्यं प्रीणाति नरकेसरी} %॥७॥

\twolineshloka
{एवं चतुर्षु वारेषु श्रावणे मासि तद् व्रतम्}
{कुर्वीत तस्य सदने लक्ष्मीः स्थिरतरा भवेत्} %॥८॥

\twolineshloka
{धनधान्यसमृद्धिश्च अपुत्रः पुत्रवान्भवेत्}
{इह लोके सुखं भुक्त्वा अन्ते वैकुण्ठमाप्नुयात्} %॥९॥

\twolineshloka
{दिग्व्यापिनी च सत्कीर्तिर्नृसिंहस्य प्रसादतः}
{एतत्ते कथितं सौम्य नृसिंहव्रतमुत्तमम्} %॥१०॥

\twolineshloka
{एवं शनिप्रीणना कर्तव्यं तच्छृणुष्व भोः}
{खञ्जं ब्राह्मणमेकं तु तदभावे तु कञ्चन} %॥११॥

\twolineshloka
{तिलतैलेन स्नापयेदुष्णवारिणा}
{नृसिंहोक्तेन चान्नेन भोजयेच्छ्रद्धयान्वितः} %॥१२॥

\twolineshloka
{तैलं लोहं तिलान्माषान्दद्यात्कम्बलमेव च}
{शनैश्चरप्रीणनाय शनिर्मे प्रीयतामिति} %॥१३॥

\twolineshloka
{अभ्यज्य शनैश्चरस्याभिषेकं तिलतैलेन कारयेत्}
{प्रशस्ता अक्षतास्तस्य पूजने तिलमाषयोः} %॥१४॥

\twolineshloka
{ध्यानं तस्य च वक्ष्यामि शृणुष्वावहितो मुने}
{शनैश्चरः कृष्णवर्णो मन्दः काश्यपगोत्रजः} %॥१५॥

\twolineshloka
{सौराष्ट्रदेशसम्भूतः सूर्यपुत्रो वरप्रदः}
{दण्डाकृतिर्मण्डले स्यादिन्द्रनीलसमद्युतिः} %॥१६॥

\twolineshloka
{बाणबाणासनधरः शूलधृग्गृध्रवाहनः}
{यमाधिदैवतश्चैव ब्रह्मप्रत्यधिदैवतः} %॥१७॥

\twolineshloka
{कस्तूर्यगुरुगन्धः स्यात्तथा गुग्गुलुधूपकः}
{कृसरान्नप्रियश्चैवं विधिरस्य प्रकीर्तितः} %॥१८॥

\twolineshloka
{प्रतिमा पूजने चास्य कार्या लोहमयी शुभा}
{अस्योद्देशेन पूजायां दानं कृष्णं द्विजोत्तम} %॥१९॥

\twolineshloka
{कृष्णवस्त्रयुगं दद्याद्दद्याद् गां कृष्णवत्सकाम्}
{एवं सम्पूज्य विधिवत्प्रार्थयेच्च स्तुवीत च} %॥२०॥

\twolineshloka
{यः पुनर्नष्टराज्याय नीलाय शनिं नीलाञ्जनप्रख्यं परितोषितः}
{ददौ निजं महाराज्यं स मे सौरिः प्रसीदतु} %॥२१॥

\twolineshloka
{मन्दचेष्टाप्रसारिणम्}
{छायामार्तण्डसम्भूतं तन्नमामि शनैश्चरम्} %॥२२॥

\twolineshloka
{नमस्ते कोणसंस्थाय पिङ्गलाय नमोऽस्तु ते}
{प्रसादं कुरु देवेश दीनस्य प्रणतस्य च} %॥२३॥

\twolineshloka
{एवं स्तुत्या प्रार्थयित्वा प्रणमेच्च पुनः पुनः}
{पूजने वैदिको मन्त्रः शन्नोदेवीरिति स्मृतः} %॥२४॥

\twolineshloka
{त्रैवर्णानां च शूद्राणां नाममन्त्रः प्रकीर्तितः}
{य एवं विधिना मन्दं पूजयेत्सुसमाहितः} %॥२५॥

\twolineshloka
{तदीयं तु भयं तस्य स्वजेऽपि न भविष्यति}
{एवमेतद् व्रतं विप्र ये करिष्यन्ति मानवाः} %॥२६॥

\twolineshloka
{वासरे वासरे प्राप्ते श्रावणे मासि भक्तितः}
{तेषां शनैश्चरकृतः पीडालेशोऽपि नो भवेत्} %॥२७॥

\twolineshloka
{प्रथमो वा द्वितीयो वा चतुर्थः पञ्चमोऽपि वा}
{सप्तमश्चाष्टमो वापि नवमो द्वादशोऽपि वा} %॥२८॥

\twolineshloka
{जन्मराशेः स्थितो मन्दः पीडां च कुरुते सदा}
{शमग्निरिति मन्त्रस्य तत्प्रसादे जपो मतः} %॥२९॥

\twolineshloka
{इन्द्रनीलमणेर्दानं प्रदद्यात्तस्य तुष्टये}
{अतः परं प्रवक्ष्यामि हनुमत्तुष्टये विधिम्} %॥३०॥

\twolineshloka
{शनिवारे श्रावणे च च अभिषेकं समाचरेत्}
{रुद्रमन्त्रेण तैलेन हनुमत्प्रीणनाय च} %॥३१॥

\twolineshloka
{तैलमिश्रितसिन्दूरलेपं तस्य समर्पयेत्}
{जपाकुसुममालाभिरर्कमालाभिरेव च} %॥३२॥

\twolineshloka
{मालाभिर्मान्दराभिश्च वटकानां तथैव च}
{पूजयेदञ्जनीपुत्रं तथान्यैरुपचारकैः} %॥३३॥

\twolineshloka
{यथाविधि यथावित्तं श्रद्धाभक्तिसमन्वितः}
{जपेद् द्वादश नामानि हनुमत्प्रीतये बुधः} %॥३४॥

\twolineshloka
{हनुमानञ्जनीसूनुर्वायुपुत्रो महाबलः}
{रामेष्टः फाल्गुनसखः पिङ्गाक्षोऽमितविक्रमः} %॥३५॥

\twolineshloka
{उदधिक्रमणश्चैव सीताशोकविनाशकः}
{लक्ष्मणप्राणदाता च दशग्रीवस्य दर्पहा} %॥३६॥

\twolineshloka
{द्वादशैतानि नामानि प्रातरुत्थाय यः पठेत्}
{नाशुभं जायते तस्य सर्वसम्पत्प्रजायते} %॥३७॥

\twolineshloka
{श्रावणे मन्दवारे तु एवमाराध्य वायुजम्}
{वज्रतुल्यशरीरः स्यादरोगो बलवान्नरः} %॥३८॥

\twolineshloka
{वेगवान्कार्यकरणे बुद्धिवैभवभूषितः}
{शत्रुः सङ्क्षयमाप्नोति मित्रवृद्धिः प्रजायते} %॥३९॥

\twolineshloka
{वीर्यवान्कीर्तिमांश्चैव प्रसादादञ्जनीजनेः}
{आञ्जनेयालये लक्षं हनुमत्कवचं पठेत्} %॥४०॥

\twolineshloka
{अणिमाद्यष्टसिद्धीनां साधकः स्वामितामियात्}
{यक्षराक्षसवेताला दर्शनात्तस्य वेगतः} %॥४१॥

\twolineshloka
{पलायन्ते दशदिशः कम्पिता भयविह्वलाः}
{अश्वत्थालिङ्गनं चैव ह्यश्वत्थस्य च पूजनम्} %॥४२॥


\threelineshloka
{मन्दभिन्ने न कर्तव्यः स्पर्शोऽश्वत्थस्य सत्तम}
{शनावालिङ्गनं तस्य सर्वसम्पत्समृद्धिदम्}
{पूजनं सप्तवारेषु तत्रापि श्रावणेऽधिकम्} %४३


{॥इति श्रीस्कन्दपुराणे ईश्वरसनत्कुमारसंवादे श्रावणमासमाहात्म्ये शनैश्चरनृसिंहहनुमत्पूजनादि-शनैश्चरकृत्यकथनं नाम दशमोऽध्यायः॥१०॥}




\sect{एकादशोऽध्यायः}

\uvacha{सनत्कुमार उवाच}
\twolineshloka
{वारव्रतानि सर्वाणि त्वत्तो देव श्रुतानि मे}
{तव वागमृतं पीत्वा तृप्तिर्मे नैव जायते} %॥१॥

\twolineshloka
{श्रावणेन समो मासो नास्त्यन्यः प्रतिभाति मे}
{अथातस्तिथिमाहात्म्यं कथयस्व जगत्प्रभो} %॥२॥


\uvacha{ईश्वर उवाच}
\twolineshloka
{मासानां कार्तिकः श्रेष्ठस्तस्मान्माघः परो मतः}
{ततोऽपि माधवः श्रेष्ठः सहश्चापि हरिप्रियः} %॥३॥

\twolineshloka
{विश्वरूपेण चत्वारो मासाश्चैते मम प्रियाः}
{द्वादशस्वपि मासेषु श्रावणः शिवरूपकः} %॥४॥

\twolineshloka
{तिथयः श्रावणे मासि सर्वाश्च व्रतसंयुताः}
{प्राधान्यतस्तथापि त्वां वच्मि काश्चित्सुशोभनाः} %॥५॥

\twolineshloka
{तिथिवारविमिश्रं तु व्रतमाद्यं वदामि ते}
{प्रतिपच्छ्रावणे मासि यदा सोमयुता भवेत्} %॥६॥

\twolineshloka
{सोमवारास्तदा पञ्च पतन्त्यत्र हि मासिके}
{रोटकाख्यं व्रतं तत्र कर्तव्यं श्रावणे नरैः} %॥७॥

\twolineshloka
{सार्धमासत्रयं वापि रोटकाख्यं व्रतं भवेत्}
{लक्ष्मीवृद्धिकरं प्रोक्तं सर्वकामार्थसिद्धिदम्} %॥८॥

\twolineshloka
{विधानं तस्य वक्ष्यामि शृणुष्वावहितो मुने}
{श्रावणस्य सिते पक्षे प्रतिपत्सोमवासरे} %॥९॥

\twolineshloka
{प्रातः सङ्कल्पयेद् विद्वान् करिष्ये रोटकव्रतम्}
{अद्यारभ्य सुरश्रेष्ठ कृपां कुरु जगद्गुरो} %॥१०॥

\twolineshloka
{दिने दिने प्रकर्तव्या पूजा देवस्य शूलिनः}
{बिल्वपत्रैरखण्डैश्च तुलसीपत्रकैस्तथा} %॥११॥

\twolineshloka
{नीलोत्पलैश्च कमलैः कहारैः कुसुमैस्तथा}
{चम्पकैर्मालतीपुष्पैः कुविन्दैरकपुष्पकैः} %॥१२॥

\twolineshloka
{अन्यैर्नानाविधैः पुष्पैर्ऋतुकालोद्भवैः शुभैः}
{धूपैर्दीपैश्च नैवेद्यैः फलैर्नानाविधैरपि} %॥१३॥

\twolineshloka
{नैवेद्यमर्पयेन्मुख्यं रोटकानां विशेषतः}
{कर्तव्या रोटकाः पञ्च पुरुषाहारमानतः} %॥१४॥

\twolineshloka
{द्वौ तु विप्राय दातव्यौ द्वाभ्यां च भोजनं मतम्}
{एको देवाय दातव्यो नैवेद्यार्थं सदा बुधैः} %॥१५॥

\twolineshloka
{शेषपूजां विधायाथ अर्घ्यं दद्याद्विचक्षणः}
{रम्भाफलं नारिकेलं जम्बीरं बीजपूरकम्} %॥१६॥

\twolineshloka
{खर्जूरी कर्कटी द्राक्षा नारङ्गं मातुलिङ्गकम्}
{अक्षोटकं च दाडिम्बं यच्चान्यदृतुसम्भवम्} %॥१७॥

\twolineshloka
{प्रशस्तमर्घ्यदाने पुण्यफलं शृणु}
{सप्तसागरसंयुक्तां भूमिं दत्त्वा तु यत्फलम्} %॥१८॥

\twolineshloka
{तत्फलं समवाप्नोति व्रतं कृत्वा विधानतः}
{पञ्चवर्षं प्रकर्तव्यमतुलं धनमीप्सुभिः} %॥१९॥

\twolineshloka
{पश्चादुद्यापनं कुर्याद्रोटकाख्यव्रतस्य तु}
{| उद्यापने तु कर्तव्यौ हेमरूप्यौ च रोटकौ} %॥२०॥

\twolineshloka
{पूर्वेद्युरधिवास्याथ प्रातर्होमं समाचरेत्}
{सर्पिषा शिवमन्त्रेण बिल्वपत्रैश्च शोभनैः} %॥२१॥

\twolineshloka
{एवं कृते व्रते तात सर्वान्कामानवाप्नुयात्}
{सनत्कुमार वक्ष्यामि द्वितीयायां व्रतं शुभम्} %॥२२॥

\twolineshloka
{यत्कृत्वा श्रद्धया मर्त्यो लक्ष्मीवान्पुत्रवान्भवेत्}
{औदुम्बराभिधं चैव तद्वतं पापनाशनम्} %॥२३॥

\twolineshloka
{श्रावणे मासि सम्प्राप्ते द्वितीयायां शुभे तिथौ}
{प्रातः सङ्कल्प्य विधिवद् व्रतं कुर्याद्विचक्षणः} %॥२४॥

\twolineshloka
{नारी वाथ नरो वापि पात्रं स्यात्सर्वसम्पदाम्}
{साक्षादुदुम्बरः पूज्यस्तदभावे तु कुड्यके} %॥२५॥

\twolineshloka
{लिखित्वा पूजयेत्तत्र चतुर्भिर्नाममन्त्रकैः}
{उदुम्बर नमस्तुभ्यं नमस्ते हेमपुष्पक} %॥२६॥

\twolineshloka
{नमो रक्ताण्डशालिने}
{तत्राधिदेवते पूज्ये शिवः शुक्रस्तथैव च} %॥२७॥

\twolineshloka
{सजन्तुफलयुक्ताय त्रयस्त्रिंशत्फलान्यस्य गृहीत्वा भागमाचरेत्}
{एकादश ब्राह्मणाय दद्यात्तावन्ति दैवते} %॥२८॥

\twolineshloka
{तावन्ति स्वयमश्नीयान्नान्नाहारस्तु तद्दिने}
{शिवं शुक्रं च सम्पूज्य रात्रौ जागरणं चरेत्} %॥२९॥

\twolineshloka
{एवं कृत्वा व्रतं तात वर्षाण्येकादशैव तु}
{पश्चादुद्यापनं कुर्याद् व्रतसम्पूर्णहेतवे} %॥३०॥

\twolineshloka
{उदुम्बरः सुवर्णेन फलपुष्पदलान्वितः}
{तत्र सम्पूजयेद्विद्वान्प्रतिमे शिवशुक्रयोः} %॥३१॥

\twolineshloka
{प्रातर्होमं चरेच्चैव ह्युदुम्बरफलैः शुभैः}
{कोमलैरल्पमात्रैश्च सङ्ख्ययाऽष्टोत्तरं शतम्} %॥३२॥

\twolineshloka
{उदुम्बरसमिद्भिश्च तिलैराज्यैश्च होमयेत्}
{एवं समाप्य होमं तु आचार्यं पूजयेत्ततः} %॥३३॥

\twolineshloka
{ब्राह्मणान्भोजयेत्पश्चाच्छतं शक्तौ दशाथ वा}
{एवं व्रते कृते वत्स फलं यत्स्याच्छृणुष्व तत्} %॥३४॥

\twolineshloka
{बहुजन्तुफलो वृक्षो यथायं साधकस्तथा}
{भवेदनेकसुतवान्वंशवृद्धिस्तथा भवेत्} %॥३५॥

\twolineshloka
{हेमपुष्यैर्यथा वृक्षस्तथा लक्ष्मीप्रदो भवेत्}
{अद्यावधि न कस्यापि व्रतमेतत्प्रकाशितम्} %॥३६॥

\twolineshloka
{गोप्याद् गोप्यतरं चैव तवाग्रे कथितं मया}
{नैवात्र संशयः कार्यो भक्त्या चैतद् व्रतं चरेत्} %॥३७॥


{॥इति श्रीस्कन्दपुराणे ईश्वरसनत्कुमारसंवादे श्रावणमासमाहात्म्ये प्रतिपोटकव्रतद्वितीयोदुम्बरव्रतकथनं नामैकादशोऽध्यायः॥११॥}




\sect{द्वादशोऽध्यायः}

\uvacha{ईश्वर उवाच}
\twolineshloka
{परं प्रवक्ष्यामि स्वर्णगौरीव्रतं शुभम्}
{श्रावणे शुक्लपक्षे तु तृतीयायां विधातृज} %॥१॥

\twolineshloka
{प्रातः स्नात्वा नित्यकर्म कृत्वा सङ्कल्पमाचरेत्}
{पार्वतीशङ्करौ पूज्यौ षोडशैरुपचारकैः} %॥२॥

\twolineshloka
{देवदेव समागच्छ प्रार्थयेऽहं जगत्पते}
{इमां मया कृतां पूजां गृहाण सुरसत्तम} %॥३॥

\twolineshloka
{वायनानि प्रदेयानि दम्पतीभ्यस्तु षोडश}
{भवान्याश्च महादेव्या व्रतसम्पूर्णहेतवे} %॥४॥

\twolineshloka
{प्रीतये द्विजवर्याय वायनं प्रददाम्यहम्}
{धानाषोडशपक्वान्नैर्वेणुपात्राणि षोडश} %॥५॥

\twolineshloka
{कुर्याद्वस्त्रादिभिर्युक्तान्याहूय द्विजदम्पतीन्}
{व्रतसम्पूर्णतार्थं तु ब्राह्मणेभ्यो ददाम्यहम्} %॥६॥

\twolineshloka
{स्वलङ्कृताः सुवासिन्यः पातिव्रत्येन भूषिताः}
{मम कार्यसमृद्ध्यर्थं प्रतिगृह्णन्तु शोभनाः} %॥७॥

\twolineshloka
{एवं षोडशवर्षाणि ह्यष्टौ चत्वारि वा पुनः}
{एकवर्षं तु सद्यो वा कृत्वा चोद्यापनं चरेत्} %॥८॥

\onelineshloka
{पूजान्ते च कथां श्रुत्वा वाचकं सम्प्रपूजयेत्} %॥९॥


\uvacha{सनत्कुमार उवाच}
\twolineshloka
{केन चीर्णं व्रतमिदं माहात्म्यं चास्य कीदृशम्}
{उद्यापनं कथं कार्यं तत्सर्वं वद मे प्रभो} %॥१०॥


\uvacha{ईश्वर उवच}
\twolineshloka
{साधु पृष्टं महाभाग कथयामि तवाग्रतः}
{स्वर्णगौरीव्रतं नाम सर्वसम्पत्करं नृणाम् } %॥११॥

\twolineshloka
{पुरा सरस्वतीतीरे सुविलाख्या महापुरी}
{तत्र चन्द्रप्रभो नाम राजासीद्धनदोपमः} %॥१२॥

\twolineshloka
{तस्यास्तां रूपलावण्ये सौन्दर्यस्मेरविभ्रमे}
{महादेवीविशालाख्ये द्विभार्ये कमलेक्षणे } %॥१३॥

\twolineshloka
{तयोः प्रियतरा ज्येष्ठा तस्यासीन्नृपतेर्मता}
{स कदाचिद्वनं भेजे मृगयासक्तमानसः} %॥१४॥

\twolineshloka
{सिंहशार्दूलवाराहवनमाहिषकुञ्जरान्}
{हत्वा बभ्राम तृष्णार्तः स राजा विपिनं महत्} %॥१५॥

\twolineshloka
{चक्रकारण्डवाकीर्णं चञ्चरीकपिकाकुलम्}
{उत्फुल्लमल्लिकाजातिकुमुदोत्पलमण्डितम्} %॥१६॥

\twolineshloka
{अपूर्वमवनीशोऽसौ ददर्शाप्सरसां सरः}
{समासाद्य सरस्तीरं पीत्वा जलमनुत्तमम्} %॥१७॥

\twolineshloka
{भक्त्या गौरीमर्चयन्तं ददर्शाप्सरसां गणम्}
{किमेतदिति पप्रच्छ राजा राजीवलोचनः} %॥१८॥

\twolineshloka
{स्वर्णगौरीव्रतमिदं क्रियतेऽस्माभिरुत्तमम्}
{सर्वसम्पत्करं नृणां तत्कुरुष्व नृपोत्तम} %॥१९॥


\uvacha{राजोवाच}
\twolineshloka
{विधानं कीदृशं ब्रूत किं फलं विस्तरान्मम}
{ता ऊचुर्योषितः सर्वास्तृतीयायां नभोयुजि} %॥२०॥

\twolineshloka
{कर्तव्यं व्रतमेतद्धि स्वर्णगौरीतिसंज्ञितम्}
{पार्वतीशङ्करौ पूज्यौ भक्त्या परमया मुदा} %॥२१॥

\twolineshloka
{बध्नीयाद्दक्षिणे करे}
{नरो वामे तु नारीणां गले वा बन्धनं मतम्} %॥२२॥

\twolineshloka
{दोरकं षोडशगुणं तत्कृत्वा सोऽपि जग्राह व्रतं नियतमानसः}
{गुणैः षोडशभिर्युक्तं दोरकं दक्षिणे करे} %॥२३॥

\twolineshloka
{बध्नामि देवदेवेशि प्रसादं कुरु मे वरम्}
{एवं देव्या व्रतं कृत्वा आजगाम निजं गृहम्} %॥२४॥

\twolineshloka
{पप्रच्छ दोरकं हस्ते दृष्ट्वा ज्येष्ठातिकोपना}
{त्रोटयित्वा च चिक्षेप बाह्यं शुष्कतरूपरि} %॥२५॥

\twolineshloka
{न कर्तव्यं न कर्तव्यमिति राज्ञि वदत्यपि}
{तेन संस्पृष्टमात्रेण तरुः पल्लवितोऽभवत्} %॥२६॥

\twolineshloka
{तद् द्वितीया ततो दृष्ट्वा विस्मयाकुलिताभवत्}
{तत्रस्थं दोरकं छिन्नं गृहीत्वा सा बबन्ध ह} %॥२७॥

\twolineshloka
{ततस्तन्मासमाहात्म्यात्पत्युः प्रियतराभवत्}
{ज्येष्ठा व्रतापचारेण सा त्यक्ता दुःखिता वनम्} %॥२८॥

\twolineshloka
{प्रययौ सा महादेवीं ध्यायन्ती मनसा च ह}
{मुनीनामाश्रमे पुण्ये निवसन्ती क्वचित् क्वचित्} %॥२९॥

\twolineshloka
{निवारिता मुनिवरैर्गच्छ पापे यथासुखम्}
{धावन्ती विपिनं घोरं निर्विण्णा निषसाद ह} %॥३०॥

\twolineshloka
{ततस्तत्कृपया देवी प्रादुरासीत्तदग्रतः}
{तां दृष्ट्वा दण्डवद्भूमौ नत्वा स्तुत्वा नृपप्रिया} %॥३१॥

\twolineshloka
{जय देवि नमस्तुभ्यं जय भक्तवरप्रदे}
{जय शङ्करवामाङ्गे जय मङ्गलमङ्गले} %॥३२॥

\twolineshloka
{ततो भक्त्या वरं लब्ध्वा गौरीमभ्यर्च्य यद् व्रतम्}
{चक्रे तस्य प्रभावेण भर्ता तां चानयद् गृहम्} %॥३३॥

\twolineshloka
{ततो देव्याः प्रसादेन सर्वान्कामानवाप सा}
{ततस्ताभ्यां नृपो राज्यं चक्रे सर्वं समृद्धिमान्} %॥३४॥

\onelineshloka
{अन्ते शिवपदं प्राप्तः कान्ताभ्यां सहितो नृपः} %॥३५॥

\twolineshloka
{यः शोभनं व्रतमिदं कुरुते शिवायाः कुर्यान्मम प्रियतरो भविता च गौर्याः}
{प्राप्य श्रियं समधिकां भुवि शत्रुसङ्घ निर्जित्य निर्मलपदं स शिवस्य याति} %॥३६॥

\twolineshloka
{एतस्योद्यापनविधिं सावधानमनाः शृणु}
{शुभे तिथौ शुभे वारे चन्द्रे ताराबलान्विते} %॥३७॥

\twolineshloka
{मण्डपेऽष्टदले पद्मे कुम्भं धान्योपरि न्यसेत्}
{पूर्णपात्रं ताम्रमयं पलषोडशनिर्मितम्} %॥३८॥

\twolineshloka
{तिलपूर्णं तत्र देवीशङ्करप्रतिमे न्यसेत्}
{श्वेतवस्त्रयुगच्छन्नं श्वेतयज्ञोपवीतिनम्} %॥३९॥

\twolineshloka
{वेदोक्तेन प्रतिष्ठा च कर्तव्या तु यथाविधि}
{सम्यक्पूजां तु सम्पाद्य रात्रौ जागरणं चरेत्} %॥४०॥

\twolineshloka
{प्रातः पूजां ततः कृत्वा ततो होमं समाचरेत्}
{ग्रहहोमं पुरा कृत्वा प्रधानं जुहुयात्ततः} %॥४१॥

\twolineshloka
{तिलाश्च यवसम्मिश्रा आज्येन च परिप्लुताः}
{द्रव्यप्रधाने सङ्ख्या तु सहस्रमथ वा शतम्} %॥४२॥

\twolineshloka
{आचार्यं पूजयेत्पश्चाद्वस्त्रालङ्कारधेनुभिः}
{वायनानि च देयानि ब्राह्मणांश्चैव भोजयेत्} %॥४३॥


\threelineshloka
{दम्पतीन् भोजयेच्चैव सङ्ख्यया षोडशैव तु}
{भूयसीं दक्षिणां दद्यात् स्वस्य वित्तानुसारतः}
{बन्धुभिः सह भुञ्जीत हर्षोत्सवसमन्वितः} %४४


{॥इति श्रीस्कन्दपुराणे ईश्वरसनत्कुमारसंवादे श्रावणमासमाहात्म्ये तृतीयायां स्वर्णगौरीव्रतकथनं नाम द्वादशोऽध्यायः॥१२॥}




\sect{त्रयोदशोऽध्यायः}

\uvacha{सनत्कुमार उवाच}

\threelineshloka
{केन व्रतेन भगवन्सौभाग्यमतुलं भवेत्}
{पुत्रपौत्रधनैश्वर्यं मनुजः सुखमेधते}
{तन्मे वद महादेव व्रतानामुत्तमं व्रतम्} %१


\uvacha{ईश्वर उवाच}
\twolineshloka
{अस्ति दूर्वा गणपतेर्व्रतं त्रैलोक्यविश्रुतम्}
{भगवत्या पुरा चीर्णं पार्वत्या श्रद्धया सह} %॥२॥


\threelineshloka
{सरस्वत्या महेन्द्रेण विष्णुना धनदेन च}
{अन्यैश्च देवैर्मुनिभिर्गन्धर्वैः किन्नरैस्तथा}
{चीर्णमेतद् व्रतं सर्वैः पुराभून्मुनिसत्तम} %३

\twolineshloka
{चतुर्थी या भवेच्छुद्धा नभोमासि सुपुण्यदा}
{तस्यां व्रतमिदं कुर्यात्सर्वपापौघनाशनम्} %॥४॥

\twolineshloka
{गजाननं चतुर्थ्यां तु एकदन्तविपाटितम्}
{विधाय हेम्ना विघ्नेशं हेमपीठासने स्थितम्} %॥५॥

\twolineshloka
{तदा हेममयी दूर्वा तदाधारे व्यवस्थितम्}
{संस्थाप्य विघ्नहर्तारं कलशे ताम्रभाजने} %॥६॥

\twolineshloka
{वेष्टिते रक्तवस्त्रेण सर्वतोभद्रमण्डले}
{पूजयेद्रक्तकुसुमैः पत्रिकाभिश्च पञ्चभिः} %॥७॥

\twolineshloka
{अपामार्गशमीदूर्वातुलसीबिल्वपत्रकैः}
{अन्यैः सुगन्धैः कुसुमैर्यथालब्धैः सुगन्धिभिः} %॥८॥

\twolineshloka
{फलैश्च मोदकैः पश्चादुपहारं प्रकल्पयेत्}
{यथावदुपचारैश्च पूजयेद् गिरिजासुतम्} %॥९॥

\twolineshloka
{प्रतिमायां स्वर्णमय्यां निर्मितायां यथाविधि}
{आवाहयामि विघ्नेशमागच्छतु कृपानिधिः} %॥१०॥

\twolineshloka
{रत्नबद्धमिदं हैमं सिंहासनमनुत्तमम्}
{आसनार्थमिदं दत्तं प्रतिगृह्णातु विश्वराट्} %॥११॥

\twolineshloka
{उमासुत नमस्तुभ्यं विश्वव्यापिन्सनातन}
{विघ्नौघं छिन्धि सकलं मम पाद्यं ददामि ते} %॥१२॥

\twolineshloka
{गणेश्वराय देवाय उमापुत्राय वेधसे}
{अर्घ्यमेतत्प्रयच्छामि गृहाण भगवन्मम} %॥१३॥

\twolineshloka
{विनायकाय शूराय वरदाय नमो नमः}
{इदमाचमनीयं ते ददामि प्रतिगृह्यताम्} %॥१४॥

\twolineshloka
{गङ्गादिसर्वतीर्थेभ्यो मया प्रार्थनयाहृतम्}
{स्नानार्थं ते मया दत्तं गृहाण सुरपुङ्गव} %॥१५॥

\twolineshloka
{सिन्दूरेण यथा लक्ष्म कुङ्कुमै रञ्जितं मया}
{वस्त्रयुग्ममिदं दत्तं गृहाण च नमोऽस्तु ते} %॥१६॥

\twolineshloka
{लम्बोदराय देवाय सर्वविघ्नापहारिणे}
{उमाङ्गमलसम्भूत चन्दनं प्रतिगृह्यताम्} %॥१७॥

\twolineshloka
{अक्षताश्च सुरश्रेष्ठ रक्तचन्दनचर्चिताः}
{मया निवेदिता भक्त्या गृहाण सुरसत्तम} %॥१८॥

\twolineshloka
{चम्पकैः केतकीपत्रैर्जपाकुसुमसङ्घकैः}
{गौरीपुत्रं पूजयामि प्रसीद तु ममोपरि} %॥१९॥

\twolineshloka
{अनुग्रहाय लोकानां दानवानां वधाय च}
{अवतीर्णः स्कन्दगुरुर्धूपं गृह्णातु वै मुदा} %॥२०॥

\twolineshloka
{परं ज्योतिः सर्वसिद्धिप्रदाय च}
{दीपं तुभ्यं प्रदास्यामि महादेवात्मने नमः} %॥२१॥

\twolineshloka
{गणानां त्वेति नैवेद्यमर्पयेन्मोदकादिकम्}
{अन्नं चतुर्विधं चैव पायसं लड्डुकादिकम्} %॥२२॥

\twolineshloka
{कर्पूरैलादिसंयुक्तं नागवल्लीदलान्वितम्}
{ताम्बूलं ते प्रदास्यामि मुखवासार्थमादरात्} %॥२३॥

\twolineshloka
{हिरण्यगर्भगर्भस्थं हेमबीजं विभावसोः}
{दक्षिणां ते प्रदास्यामि ह्यतः शान्तिं प्रयच्छ मे} %॥२४॥

\twolineshloka
{गणेश्वर गणाध्यक्ष गौरीपुत्र गजानन}
{व्रतं सम्पूर्णतां यातु त्वत्प्रसादादिभानन} %॥२५॥

\twolineshloka
{एवं सम्पूज्य विघ्नेशं यथाविभवविस्तरैः}
{सोपस्करं गणाध्यक्षमाचार्याय निवेदयेत्} %॥२६॥

\twolineshloka
{गृहाण भगवन्ब्रह्मन्गणराजं सदक्षिणम्}
{एतत्त्वद्वचनादद्य पूर्णतां यातु मे व्रतम्} %॥२७॥

\twolineshloka
{एवं यः पञ्चवर्षाणि कृत्वोद्यापनमाचरेत्}
{ईप्सितांल्लभते कामान्देहान्ते शाङ्करं पदम्} %॥२८॥

\twolineshloka
{यद्वा वर्षत्रयं कुर्यात्सर्वसिद्धिमवाप्नुयात्}
{उद्यापनं विना यस्तु करोति व्रतमुत्तमम्} %॥२९॥

\twolineshloka
{सर्वं निष्फलतां याति यथाविध्यपि यत्कृतम्}
{उद्यापनदिने प्रातस्तिलैः स्नानं समाचरेत्} %॥३०॥

\twolineshloka
{हेम्नः पलात्तदर्धार्धात्कृत्वा गणपतिं बुधः}
{पञ्चगव्यैस्तु संस्नाप्य दूर्वाभिस्तु प्रपूजयेत्} %॥३१॥

\twolineshloka
{मन्त्रैस्तु दशभिर्भक्त्या श्रद्धया सहितो नरः}
{गणाधीश नमस्तुभ्यमुमापुत्राघनाशन} %॥३२॥

\twolineshloka
{विनायकेशपुत्रेति सर्वसिद्धिप्रदायक}
{एकदन्तेभवक्त्रेति तथा मूषकवाहन} %॥३३॥

\twolineshloka
{कुमारगुरवे तुभ्यमितिनामपदैः पृथक्}
{पूर्वेद्युरधिवास्यैव प्रातर्होमं समाचरेत्} %॥३४॥

\twolineshloka
{दूर्वाभिर्मोदकैश्चैव ग्रहहोमपुरःसरम्}
{पूर्णाहुतिं ततो हुत्वा आचार्यादीन्प्रपूजयेत्} %॥३५॥

\twolineshloka
{गां सवत्सां घटोध्नीं च दद्याद्वित्तानुसारतः}
{एवं कृते व्रते वत्स सर्वान्कामानवाप्नुयात्} %॥३६॥

\twolineshloka
{मदीयप्रियपुत्रस्य व्रतेनाहं च तोषितः}
{भुवि दत्त्वा सर्वभोगं ददाम्यन्ते च सद्गतिम्} %॥३७॥

\twolineshloka
{यथा शाखाप्रशाखाभिर्दूर्वा वृद्धिं गता भवेत्}
{तथैव पुत्रपौत्रादिसन्ततिर्वृद्धिगामिनी} %॥३८॥

\twolineshloka
{इत्येतत्कथितं गुह्यं दूर्वागणपतिव्रतम्}
{श्रेष्ठाच्छ्रेष्ठतरं चैव कर्तव्यं सुखमीप्सुभिः} %॥३९॥


{॥इति श्रीस्कन्दपुराणे ईश्वरसनत्कुमारसंवादे श्रावणमासमाहात्म्ये दूर्वागणपतिव्रतकथनं नाम त्रयोदशोऽध्यायः॥१३॥}




\sect{चतुर्दशोऽध्यायः}

\uvacha{ईश्वर उवाच}
\twolineshloka
{अतः परं प्रवक्ष्यामि श्रावणे शुक्लपक्षके}
{पञ्चम्यां यच्च कर्तव्यं तच्छृणुष्व महामुने} %॥१॥

\twolineshloka
{चतुर्थ्यामेकभुक्तं तु नक्तं स्यात्पञ्चमीदिने}
{कृत्वा स्वर्णमयं नागमथवा रौप्यसम्भवम्} %॥२॥

\twolineshloka
{कृत्वा दारुमयं वापि अथवा मृण्मयं शुभम्}
{पञ्चम्यामर्चयेद्भक्त्या नागं पञ्चफणान्वितम्} %॥३॥

\twolineshloka
{द्वारस्योभयतो लेख्या गोमयेन विषोल्बणाः}
{पूजयेद् विधिवच्चैव दधिदूर्वाङ्करैः शुभैः} %॥४॥

\twolineshloka
{करवीरैर्मालतीभिर्जातिपुष्पैश्च चम्पकैः}
{तथा गन्धैरक्षतैश्च धूपैर्दीपैर्मनोहरैः} %॥५॥

\twolineshloka
{ब्राह्मणान्भोजयेत्पश्चाद् घृतमोदकपायसैः}
{अनन्तं वासुकिं शेषं पद्मनाभं च कम्बलम्} %॥६॥

\twolineshloka
{तथा कर्कोटकं नाम नागमश्वं तथाष्टमम्}
{धृतराष्ट्रं शङ्खपालं कालीयं तक्षकं तथा} %॥७॥

\twolineshloka
{हरिद्रया चन्दनेन कुड्ये नागकुलाधिपान्}
{नवकडूंश्च संलिख्य पूजयेत्कुसुमादिभिः} %॥८॥

\twolineshloka
{वल्मीके पूजयेन्नागान्दुग्धं चैव तु पाययेत्}
{घृतयुक्तं शर्कराढ्यं यथेष्टं चार्पयेद् बुधः} %॥९॥

\twolineshloka
{लोहपात्रे पोलिकादि न कुर्यात्तद्दिने नरः}
{गोधूमपायसं कुर्यान्नैवेद्यार्थं तु भक्तितः} %॥१०॥

\twolineshloka
{भर्जिताश्चणकाश्चैव व्रीहयो यावनालिकाः}
{अर्पणीयाश्च सर्पेभ्यः स्वयं चैव तु भक्षयेत्} %॥११॥

\twolineshloka
{बालकेभ्योऽर्पणीयाश्च दृढा दन्ता भवन्ति हि}
{वल्मीकस्य समीपे च गायनं वाद्यमेव च} %॥१२॥

\twolineshloka
{स्त्रीभिः कार्यं भूषिताभिः कार्यश्चैवोत्सवो महान्}
{एवं कृते कदाचिच्च सर्पतो न भयं भवेत्} %॥१३॥

\twolineshloka
{अन्यच्च शृणुयाद्विप्र लोकानां हितकाम्यया}
{कथयिष्यामि किञ्चित्ते तच्छृणुष्व महामुने} %॥१४॥

\twolineshloka
{नागदष्टो नरो वत्स प्राप्य मृत्युं व्रजत्यधः}
{अधो गत्वा भवेत्सर्पस्तामसो नात्र संशयः} %॥१५॥

\twolineshloka
{पूर्वोक्तविधिना सर्वमेकभुक्तादि कारयेत्}
{नागनिर्माणपूजादि विप्रैः सह तथादरात्} %॥१६॥

\twolineshloka
{एवं द्वादशमासेषु मासि मासि व्रतं चरेत्}
{पञ्चम्यां शुक्लपक्षस्य पूर्णे संवत्सरे पुनः} %॥१७॥

\twolineshloka
{ब्राह्मणांश्च यतींश्चैव नागानुद्दिश्य भोजयेत्}
{इतिहासविदे नागं काञ्चनं रत्नचित्रितम्} %॥१८॥

\twolineshloka
{गां च दद्यात्सवत्सां वै सर्वोपस्करसंयुताम्}
{दानकाले पठेदेतत्स्मरन्नारायणं विभुम्} %॥१९॥

\twolineshloka
{सर्वगं सर्वदातारमनन्तमपराजितम्}
{ये केचिन्मे कुले सर्पदष्टाः प्राप्ता ह्यधोगतिम्} %॥२०॥

\twolineshloka
{व्रतदानेन गोविन्द मुक्तिभाजो भवन्तु ते}
{इत्युच्चार्याक्षतैर्युक्तं सितचन्दनमिश्रितम्} %॥२१॥

\twolineshloka
{वासुदेवाग्रतो भक्त्या तोये तोयं विनिक्षिपेत्}
{अनेन विधिना सर्वे ये मरिष्यन्ति वा मृताः} %॥२२॥

\twolineshloka
{सर्पतस्तेऽभियास्यन्ति स्वर्गतिं मुनिसत्तम}
{एवं सर्वान्समुद्धृत्य कुलजान्कुलनन्दन} %॥२३॥

\twolineshloka
{प्रयाति शिवसान्निध्यं सेव्यमानोऽप्सरोगणैः}
{वित्तशाठ्यविहीनो यः सर्वमेतत्फलं लभेत्} %॥२४॥

\twolineshloka
{नक्तेन भक्तिसहिताः सितपञ्चमीषु ये हि भवन्ति पूजयन्ति सुभगान्कुसुमोपहारैः}
{तेषां गृहेष्वभयदा हि भवन्ति सर्पा हर्षान्विता मणिमयूखविभासिताङ्गाः} %॥२५॥

\twolineshloka
{प्रतिग्रहं ये च कुर्युर्गृहदानस्य वाडवाः}
{प्रयान्ति सर्पतां तेऽपि घोरां भुक्त्वा तु यातनाम्} %॥२६॥

\twolineshloka
{अन्तकाले च ये केचिन्नागहत्यावशादिह}
{मृतापत्या अपुत्रा वा भवन्ति मुनिसत्तम} %॥२७॥

\twolineshloka
{कार्पण्यवशतः स्त्रीणां सर्पतां यान्ति केचन}
{निक्षेपानृतवादाच्च केचित्सर्पा भवन्ति हि} %॥२८॥

\twolineshloka
{अन्यैश्चापि निमित्तैर्ये सर्पतां यान्ति मानवाः}
{उपायोऽयं विनिर्दिष्टः सर्वेषां निष्कृतौ परः} %॥२९॥

\twolineshloka
{वित्तशाठ्यविहीनेन कृता चेन्नागपञ्चमी}
{तद्धितार्थं हरिं शेषः सर्वनागाधिपो विभुम्} %॥३०॥

\twolineshloka
{बद्धाञ्जलिः प्रार्थयते वासुकिश्च सदाशिवम्}
{शेषवासुकिविज्ञप्त्या शिवविष्णू प्रसादितौ} %॥३१॥

\twolineshloka
{मनोरथांस्तस्य सर्वान्कुरुतः परमेश्वरौ}
{नागलोके तु तान्भोगान्भुक्त्वा तु विविधान्बहून्} %॥३२॥

\twolineshloka
{ततो वैकुण्ठमासाद्य कैलासं वापि शोभनम्}
{शिवविष्णुगणो भूत्वा लभते परमं सुखम्} %॥३३॥

\twolineshloka
{एतत्ते कथितं वत्स नागानां पञ्चमीव्रतम्}
{अतः परं किमन्यत्त्वं श्रोतुमिच्छसि तद्वद} %॥३४॥


{॥इति श्रीस्कन्दपुराणे ईश्वरसनत्कुमारसंवादे श्रावणमासमाहात्म्ये नागपञ्चमीव्रतकथनं नाम चतुर्दशोऽध्यायः॥१४॥}




\sect{पञ्चदशोऽध्यायः}

\uvacha{सनत्कुमार उवाच}
\twolineshloka
{श्रुतमाश्चर्यजनकं नागानां पञ्चमीव्रतम्}
{षष्ठ्यां कथय देवेश किं व्रतं कीदृशो विधिः} %॥१॥


\uvacha{ईश्वर उवाच}
\twolineshloka
{शुक्लपक्षे श्रावणे तु षष्ठ्यां कार्यं व्रतं शुभम्}
{सूपौदनाख्यं विप्रेन्द्र महामृत्युविनाशनम्} %॥२॥

\twolineshloka
{शिवालये गृहे वापि शिवं सम्पूज्य यत्नतः}
{सूपौदनस्य नैवेद्यमर्पयेद्विधिसंयुतः} %॥३॥

\twolineshloka
{आम्रस्य लवणं शाकं साधने परिकल्पयेत्}
{नैवेद्यस्य पदार्थैस्तु वायनं ब्राह्मणस्य च} %॥४॥

\twolineshloka
{य एतद्विधिना कुर्यात्तस्य पुण्यमनन्तकम्}
{अत्रैवोदाहरन्तीममितिहासं पुरातनम्} %॥५॥

\twolineshloka
{राजासीद्रोहितो नाम बहुकालमपुत्रवान्}
{पुत्रार्थी स तपश्चक्रे महत्परमदारुष्णम्} %॥६॥

\twolineshloka
{प्रारब्धे नास्ति ते पुत्रो बोधितोऽपि स वेधसा}
{निर्बन्धान्न निवृत्तोऽभूत्तपसः सोऽतिलालसः} %॥७॥

\twolineshloka
{ततः सङ्कटमापन्नो वेधाः प्रादुरभूत्पुनः}
{पुत्रो दत्तस्तव मया अल्पायुः स भविष्यति} %॥८॥

\twolineshloka
{पत्नी राजा मन्त्रयेतां वन्ध्यात्वं तु गमिष्यति}
{अपुत्रत्वापवादश्च अलमित्येव जायताम्} %॥९॥

\twolineshloka
{ततो ब्रह्मवरात्पुत्रो हर्षशोकपरोऽभवत्}
{जातकर्मादिसंस्कारांश्चक्रे राजा यथाविधि} %॥१०॥

\twolineshloka
{राज्ञी सा दक्षिणा नाम राजा चैव स रोहितः}
{शिवदत्त इति प्रेम्णा चक्रतुर्नाम तस्य तौ} %॥११॥

\twolineshloka
{उपनीतश्च तनयो राज्ञा तु भयचेतसा}
{विवाहं न चकारास्य भूमिपालो मृतेर्भयात्} %॥१२॥

\twolineshloka
{तदा षोडशवर्षेऽसौ मरणं प्राप पुत्रकः}
{चिन्तामाप परां राजा ब्रह्मचारिमृतिं स्मरन्} %॥१३॥

\twolineshloka
{येषां कुले ब्रह्मचारी निधनं प्राप्नुयाद्यदि}
{तत्कुलं क्षयमायाति सोऽपि दुर्गतिमापतेत्} %॥१४॥


\uvacha{सनत्कुमार उवाच}
\twolineshloka
{देवदेव जगन्नाथ परिहारोऽस्ति वा न वा}
{अस्ति चेच्च वदस्वाद्य दोषशान्तिर्यदा भवेत्} %॥१५॥


\uvacha{ईश्वर उवाच}
\twolineshloka
{स्नातको ब्रह्मचारी च निधनं प्राप्नुयाद्यदि}
{स योज्यश्चार्कविधिना संयोज्यौ तौ ततः परम्} %॥१६॥

\twolineshloka
{देशकालौ तु सङ्कीर्त्यांमुकगोत्रादिनामतः}
{व्रतं वैसर्गिकं कुर्वे मृतस्य ब्रह्मचारिणः} %॥१७॥

\twolineshloka
{हेम्नाभ्युदयिकं कृत्वा प्रतिष्ठाप्य च पावकम्}
{आघारान्तं च सम्पाद्य चतुर्व्याहृतिभिर्हुनेत्} %॥१८॥

\twolineshloka
{व्रतपत्यग्नये चैव व्रतानुष्ठानसत्फलम्}
{सम्पादनाय विश्वेभ्यो देवेभ्यश्च हुनेद् घृतम्} %॥१९॥

\twolineshloka
{ततः स्विष्टकृते हुत्वा होमशेषं समापयेत्}
{देशकालौ पुनः स्मृत्वा करिष्येऽर्कविवाहकम्} %॥२०॥

\twolineshloka
{हेम्नाभ्युदयिकं कृत्वा अर्कशाखां शवं तथा}
{लिप्त्वा तैलहरिद्राभ्यां पीतसूत्रेण वेष्टयेत्} %॥२१॥

\twolineshloka
{पीतवस्त्रयुगेनापि अग्निं संस्थापयेत्ततः}
{आघारान्तेऽग्नये चैव विवाहविधियोजकम्} %॥२२॥

\twolineshloka
{बृहस्पतये कामाय चतुर्व्याहृतिभिस्तथा}
{आज्यं स्विष्टकृतं हुत्वा कर्म चैवं समापयेत्} %॥२३॥

\twolineshloka
{अर्कशाखां शवं चैव दाहयेच्च यथाविधि}
{मृतस्य म्रियमाणस्य षडब्दं व्रतमाचरेत्} %॥२४॥

\twolineshloka
{त्रिंशद्भ्यो ब्रह्मचारिभ्यो दद्यात्कौपीनकान्नवान्}
{हस्तमात्राः कर्णमात्रा दद्यात्कृष्णाजिनानि च} %॥२५॥

\twolineshloka
{पादुकां छत्रमाल्यानि गोपीचन्दनमेव च}
{मणिप्रवालमाल्यं च भूषणानि समर्पयेत्} %॥२६॥

\onelineshloka*
{एवं कृते विधानेन विघ्नः कोऽपि न जायते}

\uvacha{ईश्वर उवाच}

\onelineshloka
{एवं श्रुत्वा ब्राह्मणेभ्यो राजा हृद्यविचारयत्} %॥२७॥

\twolineshloka
{भाति मेऽर्कविवाहोऽयमनुकल्पो न मुख्यकः}
{न ददाति प्रमीतस्य कन्यां कश्चिद्वधूं यतः} %॥२८॥

\twolineshloka
{अहं राजास्मि प्रददे रत्नानि न धनं बहु}
{ददामि तस्मै यः कश्चिद्दास्यतेऽस्य वधूं यदि} %॥२९॥

\twolineshloka
{विप्रः कश्चित्पुरे तस्मिन्नासीद्देशान्तरं गतः}
{तस्य पूर्वं मृतायास्तु भार्यायाः कन्यका शुभा} %॥३०॥

\twolineshloka
{आसीद् द्वितीया भार्या तु दुष्टचित्ताविचारयत्}
{सपत्नीद्वेषतश्चापि बहुद्रव्यस्य लोभतः} %॥३१॥

\twolineshloka
{दशवर्षा तु सा बाला दीना मातृवशं गता}
{सापत्नमाता सा लक्षं गृहीत्वा प्रददौ सुताम्} %॥३२॥

\twolineshloka
{कन्यां गृहीत्वा जग्मुस्ते श्मशानं सरितस्तटे}
{विवाहं चक्रतुश्चैव शवेन सह कन्यकाम्} %॥३३॥

\twolineshloka
{योजयित्वा विधानेन दग्धुं समुपचक्रमुः}
{ततः सा कन्यकापृच्छत्किमिदं क्रियते जनाः} %॥३४॥

\twolineshloka
{ततस्ते दुःखिताः प्रोचुर्दह्यतेऽयं पतिस्तव}
{ततः प्रोवाच सा भीता रुदती बालभावतः} %॥३५॥

\twolineshloka
{पतिः किं दह्यते मेऽसौ दग्धुं नैव ददाम्यहम्}
{गच्छध्वं सहिताः सर्वे तिष्ठाम्यत्राहमेकिका} %॥३६॥

\twolineshloka
{पत्या सह गमिष्यामि उत्तिष्ठति यदा ह्यसौ}
{दृष्ट्वा तस्यास्तु निर्बन्धं करुणादीनचेतसः} %॥३७॥

\twolineshloka
{प्रारब्धवादिनो वृद्धाः केचित्तत्रैवमूचिरे}
{अहो किं वा भावि कर्म ज्ञायते नैव कस्यचित्} %॥३८॥

\twolineshloka
{दीनपालः कृपालुश्च भगवान् किं करिष्यति}
{निराश्रिता च कन्येयं मात्रा सापत्नभावतः} %॥३९॥

\twolineshloka
{विक्रीता स्यादतो देवः कदाचित्पालको भवेत्}
{अतोऽस्माभिरशक्येयं दग्धुं चायं तथा शवः} %॥४०॥

\twolineshloka
{अतोऽस्माभिश्च गन्तव्यं सर्वेषां रोचते यदि}
{सम्मन्त्र्यैवं तु सर्वेऽपि गतास्ते नगरं प्रति} %॥४१॥

\twolineshloka
{सैका शिवं पार्वतीं च स्मरन्ती भयविह्वला}
{अजानन्ती बालभावात्किमेतदिति विह्वला} %॥४२॥

\twolineshloka
{तस्याः संस्मरणाद्देव्याः सर्वज्ञौ पार्वतीशिवौ}
{करुणापूर्णहृदयौ तत्राजग्मतुरञ्जसा} %॥४३॥

\twolineshloka
{वृषारूढौ तु तौ दृष्ट्वा दम्पती तेजसां निधी}
{ननाम दण्डवद्भूमौ न जानन्त्यपि देवते} %॥४४॥

\twolineshloka
{आश्वासनं परं लेभे आगता सङ्गतिस्त्विति}
{उवाच च पतिः किं मे जागृतो नैव जायते} %॥४५॥

\twolineshloka
{प्रसन्नौ बालभावेन दयया च परिप्लुतौ}
{उचतुस्ते जनन्यास्तु व्रतं सूपौदनाभिधम्} %॥४६॥

\twolineshloka
{व्रतं सङ्कल्प्य सतिलं गृहीत्वास्य प्रयच्छ मे}
{ब्रूहि यन्मज्जनन्यास्ति व्रतं सूपौदनाभिधम्} %॥४७॥

\twolineshloka
{कृतं तस्य प्रभावेण उत्तिष्ठतु पतिर्मम}
{तया कृतं तथा सर्वं शिवदत्तस्तथोत्थितः} %॥४८॥

\twolineshloka
{उपदिश्य व्रतं तस्यास्तदान्तर्दधतुः शिवौ}
{शिवदत्तस्तु पप्रच्छ का त्वं क्वेहागतोऽस्म्यहम्} %॥४९॥

\twolineshloka
{सा चाह किञ्चिद् वृत्तान्तं रात्रिश्चापि गताभवत्}
{प्रातर्नदीतीरगता जना राज्ञे न्यवेदयन्} %॥५०॥

\twolineshloka
{राजन् पुत्रः स्नुषा चैव नदीतीरेऽवतिष्ठतः}
{प्रामाणिकेभ्यः श्रुत्वासौ हर्षं लोकोत्तरं ययौ} %॥५१॥

\twolineshloka
{हर्षभेरीं वादयन्स नदीतीरे समाययौ}
{जनाश्च मुदिताः सर्वे प्रशशंसुर्जनाधिपम्} %॥५२॥

\twolineshloka
{राजन् गतः कालगृहं पुत्रस्ते पुनरागतः}
{प्रशशंस स्नुषां राजा किमहं शस्यते जनैः} %॥५३॥

\twolineshloka
{दुरदृष्टोऽधमश्चाहं धन्येयं सुभगा स्नुषा}
{एतत्पुण्यप्रभावेण पुत्रोऽयं जीवितो मम} %॥५४॥

\twolineshloka
{एवं स्नुषां सुसम्भाव्य राजा ब्राह्मणसत्तमान्}
{पूजयामास विभवैर्दानमानपुरः सुरम्} %॥५५॥

\twolineshloka
{बहिनतप्रमीतस्य पुनर्ग्रामप्रवेशने}
{विधिं ब्राह्मणसन्दिष्टं शान्तिकं विधिनाचरत्} %॥५६॥

\twolineshloka
{एतत्ते कथितं वत्स व्रतं सूपौदनाभिधम्}
{पञ्चवर्षाणि कृत्वैतत्पश्चादुद्यापनं चरेत्} %॥५७॥

\twolineshloka
{प्रतिमां पार्वतीशस्य अर्चयेत्प्रतिवासरे}
{प्रातर्होमं प्रकुर्वीत चरुणाम्रदलैस्तथा} %॥५८॥

\twolineshloka
{नैवेद्यं वायनं चैव व्रतोक्तविधिनाचरेत्}
{पुत्रं चिरायुषं लब्ध्वा अन्ते शिवपुरं व्रजेत्} %॥५९॥


{॥इति श्रीस्कन्दपुराणे ईश्वरसनत्कुमारसंवादे श्रावणमासमाहात्म्ये सूपौदनषष्ठीव्रतकथनं नाम पञ्चदशोऽध्यायः॥१५॥}




\sect{षोडशोऽध्यायः}

\uvacha{ईश्वर उवाच}
\twolineshloka
{अतः परं प्रवक्ष्यामि शीतलासप्तमीव्रतम्}
{श्रावणे शुक्लपक्षे तु सप्तम्यामाचरेद् व्रतम्} %॥१॥

\twolineshloka
{कुड्ये लिखित्वा वापीं तु तथा सलिलदेवताः}
{सप्तसङ्ख्या दिव्यरूपा अशरीरिणसंज्ञकाः} %॥२॥

\twolineshloka
{बालद्वययुता नारी पुरुषत्रयसंज्ञिता}
{अश्वश्च वृषभश्चैव शिबिका नरवाहना} %॥३॥

\twolineshloka
{पूजा वार्देवतानां स्यात्षोडशैरुपचारकैः}
{दध्योदनस्य नैवेद्यं साधने कर्कटीफलम्} %॥४॥

\twolineshloka
{द्विजाय वायनं दद्यान्नैवेद्यस्य पदार्थकैः}
{सप्तवर्षाणि कृत्वैवं सुवासिन्यश्च सप्त वै} %॥५॥

\twolineshloka
{प्रत्यब्दं भोजनीयाः स्युः पश्चादुद्यापनं चरेत्}
{वार्देवतानां प्रतिमा एकस्मिन् स्वर्णपात्रके} %॥६॥

\twolineshloka
{बालेन सहिताः पूज्याः सायं पूर्वेऽह्नि भक्तितः}
{प्रातर्होमं च चरुणा ग्रहहोमपुरःसरम्} %॥७॥

\twolineshloka
{व्रतमेतत्पुरा चीर्णं फलितं च तथा तथा शृणु}
{सौराष्ट्रदेशे नगरमासीच्छोभनसंज्ञितम्} %॥८॥

\twolineshloka
{तत्रासीद्धनिकः कश्चित्सर्वधर्मपरायणः}
{स वापीं खानयामास निर्जले विजने वने} %॥९॥

\twolineshloka
{पादमार्गां शुभां रम्यां बहुद्रव्यव्ययेन सः}
{पशूनां जलपानाय अपि योग्यां दृढाश्मभिः} %॥१०॥

\twolineshloka
{बद्धां चिरस्थायिनीं च बहिः प्रान्ते द्रुमैर्युतम्}
{आरामं कारयामास श्रान्तपान्थसुखाय च} %॥११॥

\twolineshloka
{परं शुष्कं जलं तत्र न लब्धं बिन्दुमात्रकम्}
{प्रयासो मे वृथा जातो द्रव्यं च व्ययितं वृथा} %॥१२॥

\twolineshloka
{इति चिन्तापरश्चासीद्धनिको धनदाभिधः}
{रात्रौ तत्रैव सुष्वाप स्वप्ने तं जलदेवताः} %॥१३॥

\twolineshloka
{आगत्य कथयामासुः शृणूपायं जलागमे}
{दास्यसे यदि ते पौत्रं बलिमस्माकमादृतः} %॥१४॥

\twolineshloka
{तदैव वापिकेयं ते जलपूर्णा भविष्यति}
{दृष्ट्वैवं गृहमागत्य पुत्रायाकथयद्धनी} %॥१५॥

\twolineshloka
{द्रविणो नाम तत्पुत्रः सोऽपि धर्मपरायणः}
{शृणुष्व मम वत्सस्य भवान्मज्जनको यतः} %॥१६॥

\twolineshloka
{तत्राप्येतद्धर्मकार्यं किं विचार्यमिह त्वया}
{स्थावरश्चास्ति धर्मोऽयं नश्वरं च सुतादिकम्} %॥१७॥

\twolineshloka
{अल्पमौल्यं महावस्तु लाभोऽयं दुर्लभः क्रयः}
{शीतांशुश्चैव चण्डांशुर्वर्तेते तनयौ मम} %॥१८॥

\twolineshloka
{शीतांशुर्नाम ज्येष्ठोऽयं बलिर्देयोऽविचारतः}
{मन्त्रोऽयं सर्वथा स्त्रीभिर्ज्ञातव्यो नैव भोः पितः} %॥१९॥

\twolineshloka
{उपायस्तत्र मत्पत्नी गर्भिणी वर्ततेऽधुना}
{आसन्नप्रसवा चैव गन्त्र्यसौ स्वपितुर्गृहे} %॥२०॥

\twolineshloka
{प्रसूत्यर्थं कनिष्ठोऽसौ तया सह गमिष्यति}
{तदा कार्यमिदं तात निर्विघ्नेन भविष्यति} %॥२१॥

\twolineshloka
{इति श्रुत्वा पुत्रवाक्यं पिता तं स तुतोष ह}
{धन्योऽसि पुत्र धन्योऽहं त्वया पुत्रेण पुत्रवान्} %॥२२॥

\twolineshloka
{एतस्मिन्नन्तरे तस्याः सुशीलायाः पितुर्गृहात्}
{आकारणं समगमत्तदा सा च जगाम ह} %॥२३॥

\twolineshloka
{ज्येष्ठोऽस्माकं समीपेऽस्तु कनिष्ठो नीयतां त्वया}
{सा तथैव सती चक्रे भर्तृश्वशुरवाक्यतः} %॥२४॥

\twolineshloka
{तदा तौ पुत्रपितरौ तैलेनाभ्यज्य बालकम्}
{स्नापयित्वा सुवस्त्रैश्च भूषणैः समलङ्कृतम्} %॥२५॥

\twolineshloka
{पूर्वाषाढावारुणर्क्षे स्थापयामासतुर्मुदा}
{वाप्या वार्देवतास्तुष्टा भवन्त्विति समूचतुः} %॥२६॥

\twolineshloka
{तदैव वापी पूर्णाभूत्सुधातुल्येन वारिणा}
{उभौ गृहं जग्मतुस्तौ हर्षशोकसमन्वितौ} %॥२७॥

\twolineshloka
{सा सुशीला पितुर्गेहेऽसूत पुत्रं तृतीयकम्}
{मासत्रयोत्तरं गेहं निजं गन्तुं च निर्गता} %॥२८॥

\twolineshloka
{वापीसमीपं प्राप्तासौ वापीं पूर्णां ददर्श च}
{विस्मयं परमं प्राप तत्र स्नानं चकार ह} %॥२९॥

\twolineshloka
{श्वशुरस्य प्रयासो मे सार्थकश्च धनव्ययः}
{तद्दिने सप्तमी चासीच्छ्रावणे शुक्लपक्षके} %॥३०॥

\twolineshloka
{सुशीलाया व्रतं चासीच्छीतलासंज्ञितं शुभम्}
{सा तत्र पाकमकरोदोदनं चानयद् दधि} %॥३१॥

\twolineshloka
{वार्देवताश्च सम्पूज्य दध्यन्नं कर्कटीफलम्}
{नैवेद्यं कल्पयामास दत्त्वा विप्राय वायनम्} %॥३२॥

\twolineshloka
{स्वयं तदेव बुभुजे सहिता सहवासिभिः}
{ततो योजनमात्रं तु तस्या ग्रामो बभूव ह} %॥३३॥

\twolineshloka
{ततः सा निर्गता चासीदारुह्य शिबिकां शुभाम्}
{बालकद्वयसंयुक्ता ततस्ता जलदेवताः} %॥३४॥

\twolineshloka
{ऊचुः परस्परं चास्याः पुत्रो देयो यतोऽनया}
{अस्माकं व्रतमाचीर्णं प्रज्ञा च विहिता परा} %॥३५॥

\twolineshloka
{एतद्व्रतप्रभावेण नूतनो दीयते सुतः}
{पूर्वजातो यदि ग्राह्यो ह्यस्मत्तोषस्य किं फलम्} %॥३६॥

\twolineshloka
{विसर्जयामासुरिति चोक्त्वान्योन्यं दयालवः}
{मातरं दर्शयामासुर्वाप्या निष्कास्य बाह्यतः} %॥३७॥

\twolineshloka
{अधावत्पृष्ठतो मातुर्मातरित्याह्वयञ्छिशुः}
{संश्रुत्य पुत्रशब्दं सा परावृत्यावलोकयत्} %॥३८॥

\twolineshloka
{दृष्ट्वा सा नन्दनं स्वीयं चकिता साभवद्धृदि}
{स्थाप्याङ्के मूर्ध्न्यवघ्राय किञ्चित्पप्रच्छ नो सुतम्} %॥३९॥

\twolineshloka
{बिभेष्यतीति बुद्ध्या सा हृदये त्वन्वचिन्तयत्}
{तस्करैर्यदि वा नीतस्तर्ह्यलङ्कारवान्कथम्} %॥४०॥

\twolineshloka
{पिशाचैर्यदि वा नीतो मोक्षितश्च पुनः कथम्}
{चिन्तासमुद्रे मग्नाः स्युर्गृहसम्बन्धिनो जनाः} %॥४१॥

\twolineshloka
{इत्येवं चिन्तयन्ती सा नगरद्वारमाप सा}
{जनाः सङ्कथयामासुः सुशीला सुसमागता} %॥४२॥

\twolineshloka
{श्रुत्वा तु पितृपुत्रौ तौ परां चिन्तामवापतुः}
{किं वदिष्यति चास्माकमस्माभिर्वा किमुच्यताम्} %॥४३॥

\twolineshloka
{एतस्मिन्नन्तरे प्राप्ता पुत्रत्रयसमन्विता}
{ज्येष्ठं दृष्ट्वा तु तं बालं श्वशुरश्च पतिश्च सः} %॥४४॥

\twolineshloka
{आश्चर्यं परमं प्राप परां मुदमवाप च}
{त्वया किं पुण्यमाचीर्णं व्रतं वापि शुचिस्मिते} %॥४५॥

\twolineshloka
{पतिव्रतासि धन्यासि पुण्यवत्यसि भामिनि}
{मासद्वयं तु सञ्जातमकस्मान्नास्त्यभूच्छिशुः} %॥४६॥

\twolineshloka
{स च त्वया पुनर्लब्धो वापी पूर्णापि चाभवत्}
{एकपुत्रा गतासीस्त्वमागतासि त्रयान्विता} %॥४७॥

\twolineshloka
{त्वयोद्धृतं कुलं सुभ्रु किं त्वां स्तौमि शुभानने}
{श्वशुरेण स्तुतैवं सा पत्या प्रेम्णा च वीक्षिता} %॥४८॥

\twolineshloka
{श्वश्र्वा चानन्दितोवाच पुण्यं मार्गस्य सर्वशः}
{प्रापुः सर्वेऽपि चानन्दं भुक्त्वा भोगान्यथेप्सितान्} %॥४९॥

\twolineshloka
{इत्येतत्कथितं वत्स शीतलासप्तमीव्रतम्}
{दध्योदनं शीतलं च शीतलं कर्कटीफलम्} %॥५०॥

\twolineshloka
{वापीजलं शीतलं तु शीतलाश्चापि देवताः}
{तापत्रयस्य सन्तापाच्छीतलाव्रतिनस्ततः} %॥५१॥

\onelineshloka
{अतो हेतोः सप्तमीयं शीतलेति यथार्थिका} %॥५२॥


{॥इति श्रीस्कन्दपुराणे ईश्वरसनत्कुमारसंवादे श्रावणमासमाहात्म्ये शीतलासप्तमीव्रतकथनं नाम षोडशोऽध्यायः॥१६॥}




\sect{सप्तदशोऽध्यायः}

\uvacha{ईश्वर उवाच}
\twolineshloka
{अथ वक्ष्यामि देवेश पवित्रारोपणं शुभम्}
{सप्तम्यामधिवास्याथ ह्यष्टम्यामर्पयेत्तु तत्} %॥१॥

\twolineshloka
{पवित्रं कारयेद्यस्तु तस्य पुण्यफलं शृणु}
{सर्वयज्ञव्रतं दानं सर्वतीर्थाभिषेचनम्} %॥२॥

\twolineshloka
{प्राप्नुयान्नात्र सन्देहो यस्मात्सर्वगता शिवा}
{नाधनो न च दुःखानि न पीडा व्याधयोऽपि च} %॥३॥

\twolineshloka
{न भयं शत्रुजं तस्य न ग्रहैः पीड्यते क्वचित्}
{सिध्यन्ति सर्वकार्याणि ह्यल्पानि च महान्ति च} %॥४॥

\twolineshloka
{नातः परतरं वत्स ह्यन्यत्पुण्यविवृद्धये}
{नराणां च नृपाणां च स्त्रीणां चैव विशेषतः} %॥५॥

\twolineshloka
{सौभाग्यजननं तात तव स्नेहात्प्रकाशितम्}
{श्रावणे शुक्लसप्तम्यामधिवास्य विधातृज} %॥६॥

\twolineshloka
{सर्वोपस्करसंयुक्तो देव्यां सद्भक्तिमांश्च सः}
{सर्वाणि पूजाद्रव्याणि गन्धपुष्पफलानि च} %॥७॥

\twolineshloka
{नैवेद्यान् विविधांश्चैव वस्त्राद्याभरणानि च}
{सम्पाद्य शोधयेदेतान्प्राशयेत्पञ्चगव्यकम्} %॥८॥

\twolineshloka
{चरुणा दिग्बलिं दद्यात्कार्यं चैवाधिवासनम्}
{छादयेत्सदृशैर्वस्त्रैः पत्रैश्चैतत्पवित्रकम्} %॥९॥

\twolineshloka
{देव्यास्तन्मूलमन्त्रेण शतवाराभिमन्त्रितम्}
{स्थापयेत्पुरतो देव्याः सर्वशोभासमन्वितम्} %॥१०॥

\twolineshloka
{देव्यास्तु मण्डपं कृत्वा रात्रौ जागरणं चरेत्}
{नटनर्तकवेश्यानां कुशलान्विविधान्गणान्} %॥११॥

\twolineshloka
{स्थायेद्वाद्यगीतादी नृत्यविद्याविशारदान्}
{प्रत्यूषे विधिवत् स्नात्वा दिग्भ्यो दद्यात्पुनर्बलीन्} %॥१२॥

\twolineshloka
{देवीं सम्पूज्य विधिवत् स्त्रियो भोज्यास्तथा द्विजाः}
{पवित्रमर्पयेद्देव्या आदावन्ते च दक्षिणाम्} %॥१३॥

\twolineshloka
{यथाशक्त्या भवेद्वत्स नियमः कार्यसाधकः}
{स्त्रियोऽक्षा मृगया मांसं राज्ञा वर्ज्य प्रयत्नतः} %॥१४॥

\twolineshloka
{स्वाध्यायश्च द्विजाचार्यैर्न कार्यं कर्षणं कृषेः}
{वणिग्भिर्न च वाणिज्यं सप्तपञ्चदिनानि वा} %॥१५॥

\twolineshloka
{अथवा त्रीणि चैकं वा दिनं तस्यार्धमेव वा}
{देव्या व्यापार आसक्तिः कर्तव्या सततं हृदि} %॥१६॥

\twolineshloka
{न करोति विधानेन पवित्रारोपणं बुधः}
{तस्य सांवत्सरी पूजा निष्फला मुनिसत्तम} %॥१७॥

\twolineshloka
{तस्माद्भक्तिसमायुक्तैर्नरैर्देवीपरायणैः}
{वर्षे वर्षे प्रकर्तव्यं पवित्रारोपणं शुभम्} %॥१८॥

\twolineshloka
{कर्काटकगते सूर्ये तथा सिंहगतेऽपि वा}
{अष्टम्यां शुक्लपक्षस्य दद्याद्देव्याः पवित्रकम्} %॥१९॥

\onelineshloka
{एतस्याकरणे दोषो नित्यमेतत्प्रकीर्तितम्} %॥२०॥


\uvacha{सनत्कुमार उवाच}
\twolineshloka
{देवदेव महादेव पवित्रं यत्त्वयोदितम्}
{निर्मितव्यं कथं स्वामिंस्तद्विधिं वद सर्वशः} %॥२१॥


\uvacha{ईश्वर उवाच}
\twolineshloka
{हेमताम्रक्षौमरूप्यैः सूत्रैः कौशेयपट्टजैः}
{कुशैः काशैश्च कार्पासैर्ब्राह्मण्या कर्तितैः शुभैः} %॥२२॥

\twolineshloka
{कृत्वा त्रिगुणितं सूत्रं त्रिगुणीकृत्य साधयेत्}
{ततोत्तमं पवित्रं तु षष्ट्या सह शतैस्त्रिभिः} %॥२३॥

\twolineshloka
{सप्तत्या सहितं द्वाभ्यां शताभ्यां मध्यमं स्मृतम्}
{साशीतिना शतेनैव कनिष्ठं तत्समाचरेत्} %॥२४॥

\twolineshloka
{उत्तमं तु शतग्रन्थि पञ्चाशद्ग्रन्थि मध्यमम्}
{पवित्रकं कनिष्ठं स्यात्षट्त्रिंशद्ग्रन्थि शोभनम्} %॥२५॥

\twolineshloka
{अथ अथवाङ्गगुणैर्वेदैर्द्वाभ्यां द्वादशतोऽपि वा}
{चतुर्विंशद्वा दशाष्टग्रन्थिभिर्वा पवित्रकम्} %॥२६॥

\twolineshloka
{चाष्टोत्तरशतं चतुःपञ्चाशदेव वा}
{सप्तविंशतिरेवैवं ज्येष्ठमध्यकनीयसम्} %॥२७॥

\twolineshloka
{अधमं नाभिमात्रं स्यादूरुमात्रं तु मध्यमम्}
{उत्तमं जानुमात्रं तत्प्रतिमाया निगद्यते} %॥२८॥

\twolineshloka
{रञ्ज्याः सर्वाः कुङ्कुमेन पवित्रग्रन्थयः शुभाः}
{देवीं पूज्य पुरोभागे सर्वतोमण्डले शुभे} %॥२९॥

\twolineshloka
{कलशे वेणुपटले पवित्राणि निधापयेत्}
{त्रिसूत्र्यां ब्रह्मविष्ण्वीशानावाह्य च ततः शृणु} %॥३०॥

\twolineshloka
{नवसूत्र्यां तथोङ्कारं सोमं वनिं विधिं तथा}
{नागांश्चन्द्ररवीशांश्च विश्वेदेवांश्च स्थापयेत्} %॥३१॥

\twolineshloka
{अतः परं प्रवक्ष्यामि स्थाप्या ग्रन्थिषु देवताः}
{क्रिया च पौरुषी वीरा विजया चापराजिता} %॥३२॥

\twolineshloka
{मनोन्मनी जया भद्रा मुक्तिरीशा तथैव च}
{प्रणवादिनमोन्तैश्च नामभिर्ग्रन्थिसङ्ख्यया} %॥३३॥

\twolineshloka
{आवर्त्यमानैरावाह्य पूजयेच्चन्दनादिभिः}
{धूपितं प्रणवेनाभिमन्त्र्य देव्यै समर्पयेत्} %॥३४॥

\twolineshloka
{एतत्ते कथितं देव्याः पवित्रारोपणं शुभम्}
{अन्येषां चैव देवानां प्रतिपत्प्रभृतिष्वपि} %॥३५॥

\twolineshloka
{पवित्रारोपणं कार्यं देवतास्ता वदामि ते}
{धनदः श्रीस्तथा गौरी गणेशः सोमराड् गुरुः} %॥३६॥

\twolineshloka
{भास्करश्चण्डिकाम्बा च वासुकिश्च तथर्षयः}
{चक्रपाणिर्ह्यनन्तश्च शिवः कः पितरस्तथा} %॥३७॥

\twolineshloka
{प्रतिपत्प्रभृतिष्वेताः पूज्यास्तिथिषु देवताः}
{मुख्याया देवतायास्तु पवित्रारोपणं त्विदम्} %॥३८॥

\onelineshloka
{तदङ्गदेवतायास्तु त्रिसूत्रं स्यात्पवित्रकम्} %॥३९॥


\uvacha{ईश्वर उवाच}
\twolineshloka
{अतः परं प्रवक्ष्यामि कर्तव्यं नवमीदिने}
{श्रावणे मासि विप्रेन्द्र पक्षयोरुभयोरपि} %॥४०॥

\twolineshloka
{कुमारी नामिका दुर्गा पूजनीया यथाविधि}
{कुर्यान्नक्तव्रतं तत्र क्षीरमाक्षिक भोजनम् } %॥४१॥

\twolineshloka
{उपवासपरो वा स्यान्नवम्यां पक्षयोर्द्वयोः}
{कुमारी वेति नाम्ना वै चण्डिकामर्चयेत्सदा} %॥४२॥

\twolineshloka
{कृत्वा रौप्यमय भक्त्या दुर्गां वै पापनाशिनीम्}
{करवीरस्य पुष्पैस्तु गन्धैरगरुचन्दनैः} %॥४३॥

\twolineshloka
{धूपेन च दशाङ्गेन मोदकैश्चापि पूजयेत्}
{कुमारीं भोजयेत्पश्चात्स्त्रियो विप्रांश्च भक्तितः} %॥४४॥

\twolineshloka
{भुञ्जीत वाग्यतः पश्चाद् बिल्वपत्रकृताशनः}
{एवं यः पूजयेद्दुर्गां श्रद्धया परया युतः} %॥४५॥

\twolineshloka
{स याति परमं स्थानं यत्र देवो गुरुः स्थितः}
{एतत्ते नवमीकृत्यं कथितं विधिनन्दन} %॥४६॥

\onelineshloka
{सर्वपापप्रशमनं सर्वसम्पत्करं नृणाम् पुत्रपौत्रादिजननमन्ते सद्गतिदायकम्} %॥४७॥


{॥इति श्रीस्कन्दपुराणे ईश्वरसनत्कुमारसंवादे श्रावणमासमाहात्म्येऽष्टम्यां देवीपवित्रारोपणं नाम सप्तदशोऽध्यायः॥१७॥}




\sect{अष्टादशोऽध्यायः}

\uvacha{सनत्कुमार उवाच}
\twolineshloka
{भगवन्पार्वतीनाथ भक्तानुग्रहकारक}
{कथयस्व दयासिन्धो माहात्म्यं दशमीतिथेः} %॥१॥


\uvacha{ईश्वर उवाच}
\twolineshloka
{श्रावणे शुक्लपक्षे तु दशम्यां प्रारभेद् व्रतम्}
{प्रतिमासे दशम्यां तु शुक्लायां व्रतमाचरेत्} %॥२॥

\twolineshloka
{एवं द्वादशमासेषु कृत्वा व्रतमनुत्तमम्}
{नभःशुक्लदशम्यां तु तत उद्यापनं चरेत्} %॥३॥

\twolineshloka
{राज्याशयो राजपुत्रः कृष्यर्थं च कृषीबलः}
{वाणिज्यार्थं वणिक्पुत्रः पुत्रार्थं गुर्विणी तथा} %॥४॥

\twolineshloka
{धर्मार्थकामसिद्ध्यर्थं लोकः कन्या वरार्थिनी}
{यष्टुकामो द्विजवरो रोग्यारोग्यार्थमेव च} %॥५॥

\twolineshloka
{चिरप्रवसिते कान्ते पत्नी तस्यागमाय च}
{एतेष्वन्येषु कर्तव्यमाशाव्रतमिदं तदा} %॥६॥

\twolineshloka
{यस्माद्यस्य भवेदार्तिः कार्यं तेन तदा व्रतम्}
{नभःशुक्लदशम्यां तु स्नात्वा सम्पूज्य देवताम्} %॥७॥

\twolineshloka
{नक्तमाशासु पूज्या वै पुष्पपल्लवचन्दनैः}
{गृहाङ्गणे लेखयित्वा यवपिष्टातकेन वा} %॥८॥

\twolineshloka
{स्त्रीरूपाश्चाधिदेवस्य शस्त्रवाहनचिह्निताः}
{दत्त्वा घृताक्तं नैवेद्यं पृथग्दीपांश्च दापयेत्} %॥९॥

\twolineshloka
{फलानि कालजातानि ततः कार्यं निवेदयेत्}
{आशाः स्वाशाः सदा सन्तु सिध्यन्तु मे मनोरथाः} %॥१०॥

\twolineshloka
{भवतीनां प्रसादेन सदा कल्याणमस्त्विति}
{एवं सम्पूज्य विधिवद् दद्याद्विप्राय दक्षिणाम्} %॥११॥

\twolineshloka
{अनेन क्रमयोगेन मासि मासि सदाचरेत्}
{वर्षमेकं मुनिश्रेष्ठ तत उद्यापनं चरेत्} %॥१२॥

\twolineshloka
{सौवर्णीः कारयेदाशा रौप्याः पिष्टातकेन वा}
{ज्ञातिबन्धुजनैः सार्धं स्नातः सम्यगलङ्कृतः} %॥१३॥

\twolineshloka
{पूजयेद्भक्तियुक्तेन चेतसा दश देवताः}
{स्थापयेत्क्रमयोगेन मन्त्रैरेभिर्गृहाङ्गणे} %॥१४॥

\twolineshloka
{त्वयि सन्निहितः शक्रः सुरासुरनमस्कृतः}
{स्वामी च भुवनस्यास्य ऐन्द्रीदिग्देवते नमः} %॥१५॥

\twolineshloka
{अग्नेः परिग्रहादाशे त्वमाग्नेयीति पठ्यसे}
{तेजोरूपा पराशक्तिरतस्त्वं वरदा भव} %॥१६॥

\twolineshloka
{धर्मराजः समाश्रित्य लोकान्संयमयत्यसौ}
{तेन संयमिनी चासि याम्ये सत्कामदा भव} %॥१७॥

\twolineshloka
{खड्गहस्तातिविक्रान्ता निरृतिस्थानमाश्रिता}
{तेन निर्ऋतिरूपासि त्वमाशां पूरयस्व मे} %॥१८॥

\twolineshloka
{त्वय्यास्ते भुवनाधारो वरुणो यादसां पतिः}
{कार्यार्थं मम धर्मार्थं वारुणि प्रवणा भव} %॥१९॥

\twolineshloka
{अधिष्ठितासि यस्मात्त्वं वायुना जगदादिना}
{वायव्ये त्वमतः शान्तिं नित्यं यच्छ ममालये} %॥२०॥

\twolineshloka
{धनाधिपाधिष्ठितासि प्रख्याता त्वमिहोत्तरा}
{निरुत्तरा भवास्मासु दत्त्वा सद्यो मनोरथम्} %॥२१॥

\twolineshloka
{ऐशानि जगदीशेन शम्भुना त्वमलङ्कृता}
{पूरयस्व शुभे देवि वाञ्छितानि नमो नमः} %॥२२॥

\twolineshloka
{सर्वलोकोपरिगता सर्वदा त्वं शिवप्रदा}
{सनकाद्यैः परिवृता मां त्राहि त्राहि सर्वदा} %॥२३॥

\twolineshloka
{नक्षत्राणि च सर्वाणि ग्रहास्तारागणास्तथा}
{नक्षत्रमातरो याश्च भूतप्रेतविनायकाः} %॥२४॥

\twolineshloka
{पूजितास्तु मया भक्त्या भक्तिप्रवणचेतसा}
{सर्वे ममेष्टसिद्ध्यर्थं भवन्तु प्रवणाः सदा} %॥२५॥

\twolineshloka
{भुजङ्गनकुलेन त्वं सेवितासि यतो ह्यधः}
{नागाङ्गनाभिः सहिता तुष्टा भव ममाद्य वै} %॥२६॥

\twolineshloka
{एभिर्मन्त्रैः समभ्यर्च्य पुष्पधूपादिना ततः}
{अलङ्कारांश्च वासांसि फलानि च निवेदयेत्} %॥२७॥

\twolineshloka
{ततो वाद्यादिनादेन गीतनृत्यादिमङ्गलैः}
{नृत्यन्तीभिर्वरस्त्रीभिर्जागरेण निशां नयेत्} %॥२८॥

\twolineshloka
{कुङ्कुमाक्षतताम्बूलदानमानादिभिः सुखम्}
{अतिवाह्य च तां रात्रिं हर्षयुक्तेन चेतसा} %॥२९॥

\twolineshloka
{प्रभाते प्रतिमा अर्च्य ब्राह्मणाय निवेदयेत्}
{अनेन विधिना कृत्वा क्षमाप्य प्रणिपत्य च} %॥३०॥

\twolineshloka
{भुञ्जीत मित्रैः सहितः सुहृद्बन्धुजनेन च}
{एवं यः कुरुते तात दशमीव्रतमादरात्} %॥३१॥

\twolineshloka
{सर्वान्कामानवाप्नोति मनसोऽभिमतान्नरः}
{स्त्रीभिर्विशेषतः कार्यं व्रतमेतत्सनातनम्} %॥३२॥

\twolineshloka
{प्राणिवर्गे यतो नार्यः श्रद्धाकामपरायणाः}
{धन्यं यशस्यमायुष्यं सर्वकामफलप्रदम्} %॥३३॥

\twolineshloka
{कथितं च मुनिश्रेष्ठ मया व्रतमिदं तव}
{नानेन सदृशं चान्यद् व्रतमस्ति जगत्त्रये} %॥३४॥

\fourlineindentedshloka
{ये मानवा विधिजपुङ्गव कामकामाः}
{सम्पूजयन्ति दशमीषु सदा दशाशाः}
{तेषामशेषनिहितान्हृदयेऽतिकामानाशाः}
{फलन्ति किमिहास्ति बहूदितेन} %॥३५॥

\twolineshloka
{मोक्षप्रदं व्रतं ह्येतन्नात्र कार्या विचारणा}
{व्रतं चानेन सदृशं न भूतं न भविष्यति} %॥३६॥


{॥इति श्रीस्कन्दपुराणे ईश्वरसनत्कुमारसंवादे श्रावणमासमाहात्म्ये आशादशमीव्रतकथनं नामाष्टादशोऽध्यायः॥१८॥}




\sect{एकोनविंशोऽध्यायः}

\uvacha{ईश्वर उवाच}
\twolineshloka
{अथ वक्ष्ये नभोमासि पक्षयोरुभयोरपि}
{एकादश्यां तु यत्कृत्यं तच्छृणुष्व महामुने} %॥१॥

\twolineshloka
{न कस्यचिन्मयाख्यातं गुह्यमेतदनुत्तमम्}
{महापुण्यप्रदं वत्स महापातकनाशनम्} %॥२॥

\twolineshloka
{वाञ्छितार्थप्रदं नृणां श्रुतं पापापहारकम्}
{श्रेष्ठं व्रतानां सर्वेषां शुभमेकादशीव्रतम्} %॥३॥

\twolineshloka
{तत्तेऽहं सम्प्रवक्ष्यामि समाहितमनाः शृणु}
{दशम्यामुषसि स्नात्वा कृतसन्ध्यादिकः शुचिः} %॥४॥

\twolineshloka
{प्राप्याज्ञां वेदविदुषः पुराणज्ञान् जितेन्द्रियान्}
{सम्पूज्य देवदेवेशं षोडशैरुपचारकैः} %॥५॥

\twolineshloka
{एकादश्यां निराहारः स्थित्वाहमपरेऽहनि}
{भोक्ष्यामि पुण्डरीकाक्ष शरणं मे भवाच्युत} %॥६॥

\twolineshloka
{कुर्याच्च नियमं वत्स गुरुदेवाग्निसन्निधौ}
{तद्दिने भूमिशायी स्यात्कामक्रोधविवर्जितः} %॥७॥

\twolineshloka
{ततः प्रभाते विमले केशवार्पितमानसः}
{श्रीधरेति तदा वाक्यं क्षुतप्रस्खलनादिषु} %॥८॥

\twolineshloka
{पाखण्डादिभिरालापं दर्शनं श्रवणं तथा}
{त्यजेद्दिनत्रयं वत्स व्रतं कैवल्यकारकम्} %॥९॥

\twolineshloka
{ततो मध्याह्नसमये नद्यादौ विमले जले}
{स्नानं कुर्याज्जितक्रोधः पञ्चगव्यपुरःसरम्} %॥१०॥

\twolineshloka
{आदित्याय नमस्कृत्य श्रीधरं शरणं व्रजेत्}
{स्ववर्णाचारविधिना कृतकृत्यो गृहं व्रजेत्} %॥११॥

\twolineshloka
{पूजयेच्छ्रीधरं तत्र श्रद्धाभक्तिपुरः सरम्}
{पुष्पधूपैस्तथा दीपैर्नैवेद्यैर्विविधैरपि} %॥१२॥

\twolineshloka
{गीतवाद्यैः कथाभिश्च जागरं कारयेन्निशि}
{कुम्भं संस्थापयित्वा तु रत्नगर्भं सकाञ्चनम्} %॥१३॥

\twolineshloka
{छादितं वस्त्रयुग्मेन सितचन्दनचर्चितम्}
{प्रतिमां देवदेवस्य शङ्खचक्रगदाभृताम्} %॥१४॥

\twolineshloka
{कृत्वा यथावत्सम्पूज्य प्रभाते विमले सति}
{द्वादश्यां कृतकृत्यस्तु श्रीधरेति जपेद् बुधः} %॥१५॥

\twolineshloka
{पूजयेद्देवदेवेशं शङ्खचक्रगदाधरम्}
{विप्राय दद्यात्कलशं हेमदक्षिणयान्वितम्} %॥१६॥

\twolineshloka
{विशेषान्नवनीतं तु तत्र देयं द्विजातये}
{श्रीधरः प्रीयतां मेऽद्य श्रियं पुष्णात्वनुत्तमाम्} %॥१७॥

\twolineshloka
{इत्युच्चार्य मुनिश्रेष्ठ समभ्यर्च्य जगद्गुरुम्}
{सम्भोज्य विप्रमुख्यांश्च दद्याच्छक्त्या च दक्षिणाम्} %॥१८॥

\twolineshloka
{भृत्यादीन्भोजयित्वा तु यवसं गोषु दापयेत्}
{स्वयं भुञ्जीत च ततः सुहृद्बन्धुसमन्वितः} %॥१९॥

\twolineshloka
{सनत्कुमार कथितस्ते शुक्लैकादशीविधिः}
{एवमेव नभोमासि कृष्णायामपि साधयेत्} %॥२०॥

\twolineshloka
{अनुष्ठानं तुल्यमेव देवनाम्नि परं भिदा}
{जनार्दनः प्रीयतां मे वाक्यमेतदुदीरयेत्} %॥२१॥

\twolineshloka
{शुक्लायां श्रीधरो देवः कृष्णायां तु जनार्दनः}
{एतत्ते सम्यगाख्यातमुभयैकादशीव्रतम्} %॥२२॥

\twolineshloka
{नानेन सदृशं पुण्यं न भूतं न भविष्यति}
{इदं त्वया गोपनीयं न देयं दुष्टमानसे} %॥२३॥


\uvacha{ईश्वर उवाच}
\twolineshloka
{अथ वक्ष्यामि द्वादश्यां पवित्रारोपणं हरेः}
{उक्तः प्रायो विधिर्देव्याः पवित्रारोपणे तव} %॥२४॥

\twolineshloka
{विशेषो यश्च तं वक्ष्ये सावधानमनाः शृणु}
{अत्राधिकारी सन्दिष्टस्तं शृणुष्व महामुने} %॥२५॥

\twolineshloka
{ब्राह्मणः क्षत्रियो वैश्यस्तथा स्त्री शूद्र एव च}
{स्वधर्मावस्थिताः सर्वे भक्त्या कुर्युः पवित्रकम्} %॥२६॥

\twolineshloka
{अतो देवेति मन्त्रेण द्विजो विष्णोर्निवेदयेत्}
{स्त्रीशूद्राणां नाममन्त्रो येन सम्पूजयेद्धरिम्} %॥२७॥

\twolineshloka
{कद्रुद्रायेति मन्त्रेण द्विजः शम्भोर्निवेदयेत्}
{स्त्रीशूद्राणां नाममन्त्रो येन सम्पूजयेद्धरम्} %॥२८॥

\twolineshloka
{कृते मणिमयं कार्यं त्रेतायां हेमसम्भवम्}
{पट्टजं द्वापरे सूत्रं कार्पासं तु कलौ स्मृतम्} %॥२९॥

\twolineshloka
{यतिभिर्मानसं कार्यं पवित्रारोपणं शुभम्}
{कृतानि च पवित्राणि वैणवे पट्टले शुभे} %॥३०॥

\twolineshloka
{संस्थाप्य शुचिवस्त्रेण पिधाय्य पुरतो न्यसेत्}
{क्रियालोपविधानार्थं यत्त्वया पिहितं प्रभो} %॥३१॥

\twolineshloka
{मयैतत्क्रियते देव तव तुष्ट्यै पवित्रकम्}
{न मे विघ्नो भवेद्देव कुरु नाथ दयां मयि} %॥३२॥

\twolineshloka
{सर्वथा सर्वदा देव मम त्वं परमा गतिः}
{एतत्पवित्रतोऽहं त्वां तोषयामि जगत्पते} %॥३३॥

\twolineshloka
{कामक्रोधादयोऽप्येते न मे स्युर्व्रतघातकाः}
{अद्यप्रभृति देवेश यावत् स्याद्वार्षिकं दिनम्} %॥३४॥

\twolineshloka
{तावद्रक्षा त्वया कार्या त्वद्भक्तस्य नमोऽस्तु ते}
{देवं सम्प्रार्थ्य कलशे पात्रे वेणुमये शुभे} %॥३५॥

\twolineshloka
{संस्थितस्य पवित्रस्य कुर्यात्प्रार्थनमादृतः}
{संवत्सरकृतार्चायाः पवित्रीकरणाय भोः} %॥३६॥

\twolineshloka
{विष्णुलोकात्पवित्राद्य आगच्छेह नमोऽस्तु ते}
{विष्णुतेजोद्भवं रम्यं सर्वपापप्रणाशनम्} %॥३७॥

\twolineshloka
{सर्वकामप्रदं देव तवाङ्गे धारयाम्यहम्}
{आमन्त्रितोऽसि देवेश पुराणपुरुषोत्तम} %॥३८॥

\twolineshloka
{अतस्त्वां पूजयिष्यामि सान्निध्यं कुरु ते नमः}
{निवेदयाम्यहं तुभ्यं प्रातरेतत्पवित्रकम्} %॥३९॥

\twolineshloka
{ततः पुष्पाञ्जलिं दत्त्वा रात्रौ जागरणं चरेत्}
{एकादश्यामधिवसेद् द्वादश्यामर्चयेदुषः} %॥४०॥

\twolineshloka
{गन्धदूर्वाक्षतैर्युक्तं समादाय पवित्रकम्}
{देवदेव नमस्तुभ्यं गृहाणेदं पवित्रकम्} %॥४१॥

\twolineshloka
{पवित्रीकरणार्थाय वर्षपूजाफलप्रदम्}
{पवित्रं मां कुरुष्वाद्य यन्मया दुष्कृतं कृतम्} %॥४२॥

\twolineshloka
{शुद्धो भवाम्यहं देव त्वत्प्रसादात्सुरेश्वर}
{मूलसम्पुटितैरेतैर्मन्त्रैर्दद्यात्पवित्रकम्} %॥४३॥

\twolineshloka
{महानैवेद्यकं दत्त्वा नीराज्य प्रार्थयेत्ततः}
{मूलमन्त्रेण जुहुयाद्वह्नौ सघृतपायसम्} %॥४४॥

\twolineshloka
{विसर्जयित्वा मन्त्रेण अनेनैव पवित्रकम्}
{सांवत्सरीं शुभां पूजां सम्पाद्य विधिवन्मम} %॥४५॥

\twolineshloka
{व्रजेदानीं पवित्र त्वं विष्णुलोकं विसर्जितम्}
{उत्तार्य ब्राह्मणे दद्यात्तोये वाथ विसर्जयेत्} %॥४६॥

\twolineshloka
{एतत्ते कथितं वत्स पवित्रारोपणं हरेः}
{इह लोके सुखं भुक्त्वा ह्यन्ते वैकुण्ठमाप्नुयात्} %॥४७॥


{॥इति श्रीस्कन्दपुराणे ईश्वरसनत्कुमारसंवादे श्रावणमासमाहात्म्ये उभयैकादशीव्रतकथनं द्वादश्यां विष्णुपवित्रारोपणकथनं नामैकोनविंशोऽध्यायः॥१९॥}




\sect{विंशोऽध्यायः}

\uvacha{ईश्वर उवाच}
\twolineshloka
{त्रयोदशीदिने कृत्यं कथयामि तवाग्रतः}
{अत्रानङ्गः पूजनीयः षोडशैरुपचारकैः} %॥१॥

\twolineshloka
{अशोकैर्मालतीपुष्पैः पद्यैर्देवप्रियैस्तथा}
{कौसुम्भैर्बकुलैः पुष्पैस्तथान्यैरपि मादकैः} %॥२॥

\twolineshloka
{रक्ताक्षतैः पीतगन्धैर्द्रव्यैः सौगन्धिकैः शुभैः}
{पुष्टिकाजनकैर्द्रव्यै रेतोवृद्धिकरैः परैः} %॥३॥

\twolineshloka
{नैवेद्यमर्पयेच्चैव ताम्बूलं मुखरोचकम्}
{ताम्बूले योजयेद् द्रव्यं चिक्कणं क्रमुकं शुभम्} %॥४॥

\twolineshloka
{खादिरं चूर्णकं जातित्वचं जातिफलं तथा}
{लवङ्गैलानारिकेलबीजस्य शकलं लघु} %॥५॥

\twolineshloka
{स्वर्णरूप्याणि पत्राणि कर्पूरं केसरं तथा}
{जातानि मगधे देशे नागवल्लीदलानि च} %॥६॥

\twolineshloka
{श्वेतवर्णानि पक्वानि जीर्णानि च दृढानि च}
{रसयुक्तानि देयानि प्रीतये शम्बरद्विषः} %॥७॥

\twolineshloka
{माक्षीकमलसारेण निर्मिताभिश्च वर्तिभिः}
{नीराजयेच्चित्तभवं पुष्पाञ्जलिमथार्पयेत्} %॥८॥

\twolineshloka
{प्रार्थयेन्नामभिस्तस्य तानि ते कथयाम्यहम्}
{सर्वोपमानसौन्दर्यः प्रद्युम्नाख्यो हरेः सुतः} %॥९॥

\twolineshloka
{मीनकेतनकन्दर्पकानङ्गा मन्मथस्तथा}
{मारः कामात्मसम्भूतो झषकेतुर्मनोभवः} %॥१०॥

\twolineshloka
{रतिपीनघनोत्तुङ्गस्तनयोः पत्रवल्लिका}
{यस्य वक्षसि कस्तूर्याः शोभते परिरम्भणात्} %॥११॥

\twolineshloka
{पुष्पधन्वञ्छम्बरारे कुसुमेषो रतेः पते}
{मकरध्वज पञ्चेषो मदन स्मर सुन्दर} %॥१२॥

\twolineshloka
{देवानां कार्यसिद्धयर्थं शिवक्षिप्तहुताशन}
{परोपकारसीमानं ध्वनयंस्तेन कर्मणा} %॥१३॥

\twolineshloka
{निमित्तमात्रं विजये वसन्तस्य सहायता}
{त्वन्मनोरञ्जने शक्रस्तिष्ठत्येव दिवानिशि} %॥१४॥

\twolineshloka
{स्वपदभ्रंशने यस्मात्तपस्विभ्यो बिभेति सः}
{त्वदन्यः शम्भुना कोऽन्यो विरुध्येद् दृढमानसः} %॥१५॥

\twolineshloka
{परब्रह्मानन्दसमानन्ददस्त्वदृतेऽत्र कः}
{महामोहस्य सैन्येषु त्वादृशः कोऽस्ति वीर्यवान्} %॥१६॥

\twolineshloka
{अनिरुद्धपतिः कृष्णात्मजो कृष्णात्मजो यश्च सुरप्रभुः}
{मलयाचलसम्भूतचन्दनागरुवासितः} %॥१७॥

\twolineshloka
{दक्षिणादिङ्मातरिश्वा सहायस्ते जगज्जये}
{शरत्सुधांशुसन्मित्र जगत्सर्जनकारण} %॥१८॥

\twolineshloka
{नाथ त्वदस्त्रं परमममोघमतिदूरगम्}
{मर्मच्छिदामकरुणं रहितं प्रतिकारतः} %॥१९॥

\twolineshloka
{सुकुमारं श्रुतमपि निःसीमक्षोभकारणम्}
{स्वतुल्यस्य पदार्थस्य दर्शनादपि साधकम्} %॥२०॥

\twolineshloka
{प्रवृत्तिर्मुख्यालङ्कारः सहायेन जगज्जये}
{सर्वे श्रेष्ठास्त्वया देवा उपहास्याः कृता विभो} %॥२१॥

\twolineshloka
{ब्रह्मा कन्यालम्पटोऽभूद् वृन्दासक्तो हरिः स्मृतः}
{परदारकलङ्केन अस्पृष्टव्यः शिवो यतः} %॥२२॥

\twolineshloka
{स्वशक्त्यामेव निरतो बहुकालं व्यवायवान्}
{दुष्कर्मनिरतश्चेन्द्रो गौतमस्य वधूं प्रति} %॥२३॥

\twolineshloka
{द्विजराजो गुरोर्भार्यां बलादेवापहारवान्}
{विश्वामित्रस्तपोभ्रष्टः केनाकारि च भूयसा} %॥२४॥

\twolineshloka
{उक्ताः प्राधान्यतस्त्वेते किं बहूक्तेन मानद}
{विरलाः सन्ति लोकेऽस्मिन्ब्राह्मणा वशवर्तिनः} %॥२५॥

\twolineshloka
{तस्मात्प्रसीद भगवन्कृतया पूजयानया}
{पूजितः श्रावणे शुक्लत्रयोदश्यां मनोभवः} %॥२६॥

\twolineshloka
{प्रवृत्तिलम्पटस्यातिवीर्यं पुष्टिं ददात्यलम्}
{निवृत्तिमार्गनियतः स्वविकारं हरत्यपि} %॥२७॥

\twolineshloka
{सकामस्य स्त्रियो रम्याः पीनोत्तुङ्गपयोधराः}
{शरत्पूर्णसुधारश्मिवदनाः कमलेक्षणाः} %॥२८॥

\twolineshloka
{लम्बातिनीलकुरलस्निग्धकेश्यः सुनासिकाः}
{रम्भोरूर्वा गुप्तगुल्फा गतिनिर्जितकुञ्जराः} %॥२९॥

\twolineshloka
{कामागारा जिताश्वत्थपलाशा अतिशोभनाः}
{बृहच्छ्रोण्यः कम्बुकण्ठ्यो बृहज्जघनशोभिताः} %॥३०॥

\twolineshloka
{बिम्बोष्ठ्यः सिंहकट्यश्च नानालङ्कारभूषिताः}
{मनोरमा ददात्येष सन्तुष्टः श्रावणेऽर्चया} %॥३१॥

\twolineshloka
{शुक्लपक्षे त्रयोदश्यां ददाति च सुतान्बहून्}
{चिरायुषो गुणाढ्यांश्च सुखरूपान्सुसन्ततीन्} %॥३२॥

\twolineshloka
{कर्तव्यं यत् त्रयोदश्यामेतत्ते कथितं शुभम्}
{अतः परं चतुर्दश्यां कर्तव्यं शृणु मानद} %॥३३॥

\twolineshloka
{अष्टम्यां कथितं देव्याः पवित्रारोपणं तव}
{तत्र चेन्न कृतं तर्हि चतुर्दश्यां तु कारयेत्} %॥३४॥

\twolineshloka
{पवित्रं तु त्रिनेत्रस्य चतुर्दश्यां समर्पयेत्}
{पवित्रसाधनं सर्वं देवीविष्णुपवित्रवत्} %॥३५॥

\twolineshloka
{ऊहः परं प्रकर्तव्यः प्रार्थनादिषु नामसु}
{शैवागमे मया प्रोक्तं जाबालादिषु यत्परम्} %॥३६॥

\twolineshloka
{विकल्पात्कश्चिदस्तीह विशेषस्तं वदामि ते}
{एकादशाथ वा सूत्रैस्त्रिंशता चाष्टयुक्तया} %॥३७॥

\twolineshloka
{पञ्चाशता वा कर्तव्यं तुल्यग्रन्थ्यन्तरालकम्}
{द्वादशाङ्गुलमानानि तथा चाष्टाङ्गुलानि वा} %॥३८॥

\twolineshloka
{लिङ्गविस्तारमानानि चतुरङ्गुलिकानि वा}
{अर्पयेच्छिवतुष्ट्यर्थं विधिः पूर्वोक्त एव हि} %॥३९॥

\twolineshloka
{फलादि पूर्वमेवोक्तमन्ते कैलासमाप्नुयात्}
{एतत्ते कथितं वत्स किमन्यच्छ्रोतुमिच्छसि} %॥४०॥


{॥इति श्रीस्कन्दपुराणे ईश्वरसनत्कुमारसंवादे श्रावणमासमाहात्म्ये त्रयोदशीचतुर्दशीकर्तव्यकथनं नाम विंशोऽध्यायः॥२०॥}




\sect{एकविंशोऽध्यायः}

\uvacha{सनत्कुमार उवाच}
\twolineshloka
{पौर्णमास्या विधिं ब्रूहि कृपां कृत्वा दयानिधे}
{माहात्म्यं शृण्वतां स्वामिञ्छ्रवणेच्छा प्रवर्धते} %॥१॥


\uvacha{ईश्वर उवाच}
\twolineshloka
{उत्सर्जनमुपाकर्म अध्यायानां भवेदिह}
{पौषपूर्णा माघपूर्णा अथवोत्सर्जने तिथिः} %॥२॥

\twolineshloka
{पौषस्य प्रतिपद्वापि माघमासस्य वा भवेत्}
{ऋक्षं वा रोहिणीसंज्ञमुत्सर्जनकृतौ भवेत्} %॥३॥

\twolineshloka
{अथवान्येषु कालेषु स्वस्वशाखानुसारतः}
{सहप्रयोगो युक्तः स्यादुत्सर्गप्रकृतिद्वये} %॥४॥

\twolineshloka
{अतो नभःपौर्णमास्यामुत्सर्जनमिहेष्यते}
{उपाकर्मणि चैवं स्याच्छ्रवणर्थं तु बवृचाम्} %॥५॥

\twolineshloka
{चतुर्दश्यां पौर्णमास्यां प्रतिपद्दिवसेऽपि वा}
{यत्र वा श्रवणर्थं स्याद् बह्वृचानां तु तद्दिने} %॥६॥

\twolineshloka
{यजुषां पौर्णमास्यां स्यात्सामगानां तु हस्तभे}
{शुक्रगुर्वीरस्तमये उपाकर्म चरेत्सुखम्} %॥७॥

\twolineshloka
{आरम्भः प्रथमो न स्यादिति शास्त्रविदां मतम्}
{ग्रहसङ्क्रान्तदुष्टे तु काले कालान्तरे भवेत्} %॥८॥

\twolineshloka
{पञ्चम्यां हस्तयुक्तायां पूर्णायां वा नभस्यके}
{स्वस्वगृह्यानुसारेण उत्सर्जनमुपाकृतिः} %॥९॥

\twolineshloka
{मलमासे तु सम्प्राप्ते शुद्धे मासि तु सा भवेत्}
{नित्यं कर्मद्वयं चेदं प्रत्यब्दं नियमाच्चरेत्} %॥१०॥

\twolineshloka
{उपाकर्मसमाप्तौ तु संस्थितेषु द्विजातिषु}
{अर्पणीयः सभादीपो योषिद्भिस्तत्र संसदि} %॥११॥

\twolineshloka
{आचार्यः प्रतिगृह्णाति दद्याद्वान्यद्विजातये}
{सौवर्णे राजते वापि पात्रे ताम्रमयेऽपि वा} %॥१२॥

\twolineshloka
{प्रस्थमात्रं तु गोधूमा दीपं तत्पिष्टसम्भवम्}
{दीपपात्रं संविधाय ज्वालयेत्तत्र दीपकम्} %॥१३॥

\twolineshloka
{आज्येन वाथ तैलेन वर्तित्रयसमन्वितम्}
{सदक्षिणं सताम्बूलं ब्राह्मणाय निवेदयेत्} %॥१४॥

\twolineshloka
{दीपं सम्पूज्य विप्रं च मन्त्रमेतमुदीरयेत्}
{सदक्षिणः सताम्बूलः सभादीपोऽयमुत्तमः} %॥१५॥

\twolineshloka
{अर्पितो देवदेवस्य मम सन्तु मनोरथाः}
{सभादीपप्रदानेन पुत्रपौत्रादिकं कुलम्} %॥१६॥

\twolineshloka
{सर्वं ह्युज्ज्वलतां याति वर्धते यशसा सह}
{स्वरङ्गनाभिः सदृशं रूपं जन्मान्तरे लभेत्} %॥१७॥

\twolineshloka
{सौभाग्यं चैव लभते भर्तुः प्रियतरा भवेत्}
{एवं कृत्वा पञ्चवर्षं तत उद्यापनं चरेत्} %॥१८॥

\twolineshloka
{विप्राय दक्षिणां दद्याद्यथाशक्ति च भक्तितः}
{सभादीपस्य माहात्म्यमेतत्ते कथितं शुभम्} %॥१९॥

\twolineshloka
{श्रवणाकर्मसंस्था च तस्यामेव निशि स्मृता}
{तदुत्तरं सर्पबलिस्तत्रैव च विधीयते} %॥२०॥

\twolineshloka
{इदं संस्थाद्वयं कुर्यात्स्वस्वगृह्यमवेक्ष्य च}
{हयग्रीवस्यावतारस्तस्यामेव तिथौ मतः} %॥२१॥

\twolineshloka
{हयग्रीवजयन्त्यास्तु अतोऽत्रैव महोत्सवः}
{उपासनावतां तस्य नित्यस्तु परिकीर्तितः} %॥२२॥

\twolineshloka
{श्रावण्यां श्रवणे पूर्वं जातो हयशिरा हरिः}
{जगौ स सामवेदं तु सर्वकिल्बिषनाशनम्} %॥२३॥

\twolineshloka
{सिन्धूनदीवितस्तायां प्रवृत्तस्तत्र सङ्गमे}
{श्रवणर्क्षे ततस्तत्र स्नानं सर्वार्थसिद्धिदम्} %॥२४॥

\twolineshloka
{तत्र सम्पूजयेद्विष्णुं शार्ङ्गचक्रगदाधरम्}
{श्रोतव्यान्यथ सामानि पूज्या विप्राश्च सर्वथा} %॥२५॥

\twolineshloka
{क्रीडितव्यं च भोक्तव्यं तत्रैव स्वजनैः सह}
{जलक्रीडा च कर्तव्या नारीभिर्भर्तृलब्धये} %॥२६॥

\twolineshloka
{स्वस्वदेशे स्वस्वगृहे अपि कुर्यान्महोत्सवम्}
{पूजयेच्च हयग्रीवं जपेन्मन्त्रं च तं शृणु} %॥२७॥

\twolineshloka
{प्रणवादिर्नमः शब्दस्ततो भगवते इति}
{धर्मायाथ चतुर्थ्यन्तं योज्यं चात्मविशोधनम्} %॥२८॥

\twolineshloka
{पुनरन्ते नमः शब्दो मन्त्रश्चाष्टादशाक्षरः}
{सर्वसिद्धिकरश्चायं षट्प्रयोगैकसाधकः} %॥२९॥

\twolineshloka
{पुरश्चरणमेतस्य अक्षराणां तु सङ्ख्यया}
{लक्षं वाथ सहस्रं वा कलौ तु स्याच्चतुर्गुणम्} %॥३०॥

\twolineshloka
{एवं कृते हयग्रीवस्तुष्टः सत्कामदो भवेत्}
{एतस्यामेव पूर्णायां रक्षाबन्धनमिष्यते} %॥३१॥

\twolineshloka
{सर्वरोगोपशमनं सर्वाशुभविनाशनम्}
{शृणु त्वं मुनिशार्दूल इतिहासं पुरातनम्} %॥३२॥

\twolineshloka
{इन्द्राण्या यत्कृतं पूर्वमिन्द्रस्य जयसिद्धये}
{देवासुरमभूद्युद्धं पुरा द्वादशवार्षिकम्} %॥३३॥

\twolineshloka
{शक्रं दृष्ट्वा तदा श्रान्तं देवी प्राह सुरेश्वरम्}
{अद्य भूतदिनं देव प्रातः सर्वं भविष्यति} %॥३४॥

\twolineshloka
{अहं रक्षां विधास्यामि तेनाजेयो भविष्यसि}
{इत्युक्त्वा पौर्णमास्यां सा पौलोमी कृतमङ्गला} %॥३५॥

\twolineshloka
{बबन्ध दक्षिणे पाणौ रक्षां मोदप्रदां ततः}
{बद्धरक्षस्ततः शक्रः कृतस्वस्त्ययनो द्विजैः} %॥३६॥

\twolineshloka
{दुद्राव दानवानीकं क्षणाजिग्ये प्रतापवान्}
{वासवो विजयी भूत्वा पुनरेव जगत्त्रये} %॥३७॥

\twolineshloka
{एष प्रभावो रक्षायाः कथितस्ते मुनीश्वर}
{जयदः सुखदश्चैव पुत्रारोग्यधनप्रदः} %॥३८॥


\uvacha{सनत्कुमार उवाच}
\twolineshloka
{क्रियते केन विधिना रक्षाबन्धः सुरोत्तम}
{कस्यां तिथौ कदा देव एतन्मे वक्तुमर्हसि} %॥३९॥

\twolineshloka
{यथा यथा हि भगवन्विचित्राणि प्रभाषसे}
{तथा तथा न मे तृप्तिर्बह्वर्थाः शृण्वतः कथाः} %॥४०॥


\uvacha{ईश्वर उवाच}
\twolineshloka
{सम्प्राप्ते श्रावणे मासि पौर्णमास्यां दिनोदये}
{स्नानं कुर्वीत मतिमान् श्रुतिस्मृतिविधानतः} %॥४१॥

\twolineshloka
{सन्ध्याजपादि सम्पाद्य पितॄन्देवानृषींस्तथा}
{तर्पयित्वा ततः कुर्यात्स्वर्णपात्रविनिर्मिताम्} %॥४२॥

\twolineshloka
{हेमसूत्रैश्च सम्बद्धां मौक्तिकादिविभूषिताम्}
{कौशेयतन्तुभिः कीर्णैर्विचित्रैर्मलवर्जितैः} %॥४३॥

\twolineshloka
{विचित्रग्रन्थिसंयुक्तां पदगुच्छैश्च राजिताम्}
{सिद्धार्थैश्चाक्षतैश्चैव गर्भितां सुमनोहराम्} %॥४४॥

\twolineshloka
{संस्थाप्य कलशं तत्र पूर्णपात्रे तु तां न्यसेत्}
{उपविश्यासने रम्ये सुहृद्भिः परिवारितः} %॥४५॥

\twolineshloka
{वेश्यानर्तनगानादिकृतकौतुकमङ्गलः}
{ततः पुरोधसा कार्यो रक्षाबन्धः समन्त्रकः} %॥४६॥

\twolineshloka
{येन बद्धो बली राजा दानवेन्द्रो महाबलः}
{तेन त्वामनुबध्नामि रक्षे मा चल मा चल} %॥४७॥

\twolineshloka
{ब्राह्मणैः क्षत्रियैर्वैश्यैः शूद्रैश्चैवान्यमानवैः}
{रक्षाबन्धः प्रकर्तव्यो द्विजान्सम्पूज्य यत्नतः} %॥४८॥

\twolineshloka
{अनेन विधिना यस्तु रक्षाबन्धनमाचरेत्}
{स सर्वदोषरहितः सुखी संवत्सरं भवेत्} %॥४९॥

\twolineshloka
{यः श्रावणे विमलमासि विधानविज्ञो रक्षाविधानमिदमाचरते मनुष्यः}
{आस्ते सुखेन परमेण स वर्षमेकं पुत्रैश्च पौत्रसहितः ससुहृज्जनश्च} %॥५०॥

\twolineshloka
{भद्रायां च न कर्तव्यो रक्षाबन्धः शुचिव्रतैः}
{बद्धा रक्षा तु भद्रायां विपरीतफलप्रदा} %॥५१॥


{॥इति श्रीस्कन्दपुराणे ईश्वरसनत्कुमारसंवादे श्रावणमासमाहात्म्ये उपाकर्मोत्सर्जन-श्रवणाकर्मसर्पबलिसभादीपहयग्रीवजयन्ती-रक्षाबन्धविधिकथनं नामैकविंशोऽध्यायः॥२१॥}




\sect{विंशोऽध्यायः}

\uvacha{ईश्वर उवाच}
\twolineshloka
{त्रयोदशीदिने कृत्यं कथयामि तवाग्रतः}
{अत्रानङ्गः पूजनीयः षोडशैरुपचारकैः} %॥१॥

\twolineshloka
{अशोकैर्मालतीपुष्पैः पद्यैर्देवप्रियैस्तथा}
{कौसुम्भैर्बकुलैः पुष्पैस्तथान्यैरपि मादकैः} %॥२॥

\twolineshloka
{रक्ताक्षतैः पीतगन्धैर्द्रव्यैः सौगन्धिकैः शुभैः}
{पुष्टिकाजनकैर्द्रव्यै रेतोवृद्धिकरैः परैः} %॥३॥

\twolineshloka
{नैवेद्यमर्पयेच्चैव ताम्बूलं मुखरोचकम्}
{ताम्बूले योजयेद् द्रव्यं चिक्कणं क्रमुकं शुभम्} %॥४॥

\twolineshloka
{खादिरं चूर्णकं जातित्वचं जातिफलं तथा}
{लवङ्गैलानारिकेलबीजस्य शकलं लघु} %॥५॥

\twolineshloka
{स्वर्णरूप्याणि पत्राणि कर्पूरं केसरं तथा}
{जातानि मगधे देशे नागवल्लीदलानि च} %॥६॥

\twolineshloka
{श्वेतवर्णानि पक्वानि जीर्णानि च दृढानि च}
{रसयुक्तानि देयानि प्रीतये शम्बरद्विषः} %॥७॥

\twolineshloka
{माक्षीकमलसारेण निर्मिताभिश्च वर्तिभिः}
{नीराजयेच्चित्तभवं पुष्पाञ्जलिमथार्पयेत्} %॥८॥

\twolineshloka
{प्रार्थयेन्नामभिस्तस्य तानि ते कथयाम्यहम्}
{सर्वोपमानसौन्दर्यः प्रद्युम्नाख्यो हरेः सुतः} %॥९॥

\twolineshloka
{मीनकेतनकन्दर्पकानङ्गा मन्मथस्तथा}
{मारः कामात्मसम्भूतो झषकेतुर्मनोभवः} %॥१०॥

\twolineshloka
{रतिपीनघनोत्तुङ्गस्तनयोः पत्रवल्लिका}
{यस्य वक्षसि कस्तूर्याः शोभते परिरम्भणात्} %॥११॥

\twolineshloka
{पुष्पधन्वञ्छम्बरारे कुसुमेषो रतेः पते}
{मकरध्वज पञ्चेषो मदन स्मर सुन्दर} %॥१२॥

\twolineshloka
{देवानां कार्यसिद्धयर्थं शिवक्षिप्तहुताशन}
{परोपकारसीमानं ध्वनयंस्तेन कर्मणा} %॥१३॥

\twolineshloka
{निमित्तमात्रं विजये वसन्तस्य सहायता}
{त्वन्मनोरञ्जने शक्रस्तिष्ठत्येव दिवानिशि} %॥१४॥

\twolineshloka
{स्वपदभ्रंशने यस्मात्तपस्विभ्यो बिभेति सः}
{त्वदन्यः शम्भुना कोऽन्यो विरुध्येद् दृढमानसः} %॥१५॥

\twolineshloka
{परब्रह्मानन्दसमानन्ददस्त्वदृतेऽत्र कः}
{महामोहस्य सैन्येषु त्वादृशः कोऽस्ति वीर्यवान्} %॥१६॥

\twolineshloka
{अनिरुद्धपतिः कृष्णात्मजो कृष्णात्मजो यश्च सुरप्रभुः}
{मलयाचलसम्भूतचन्दनागरुवासितः} %॥१७॥

\twolineshloka
{दक्षिणादिङ्मातरिश्वा सहायस्ते जगज्जये}
{शरत्सुधांशुसन्मित्र जगत्सर्जनकारण} %॥१८॥

\twolineshloka
{नाथ त्वदस्त्रं परमममोघमतिदूरगम्}
{मर्मच्छिदामकरुणं रहितं प्रतिकारतः} %॥१९॥

\twolineshloka
{सुकुमारं श्रुतमपि निःसीमक्षोभकारणम्}
{स्वतुल्यस्य पदार्थस्य दर्शनादपि साधकम्} %॥२०॥

\twolineshloka
{प्रवृत्तिर्मुख्यालङ्कारः सहायेन जगज्जये}
{सर्वे श्रेष्ठास्त्वया देवा उपहास्याः कृता विभो} %॥२१॥

\twolineshloka
{ब्रह्मा कन्यालम्पटोऽभूद् वृन्दासक्तो हरिः स्मृतः}
{परदारकलङ्केन अस्पृष्टव्यः शिवो यतः} %॥२२॥

\twolineshloka
{स्वशक्त्यामेव निरतो बहुकालं व्यवायवान्}
{दुष्कर्मनिरतश्चेन्द्रो गौतमस्य वधूं प्रति} %॥२३॥

\twolineshloka
{द्विजराजो गुरोर्भार्यां बलादेवापहारवान्}
{विश्वामित्रस्तपोभ्रष्टः केनाकारि च भूयसा} %॥२४॥

\twolineshloka
{उक्ताः प्राधान्यतस्त्वेते किं बहूक्तेन मानद}
{विरलाः सन्ति लोकेऽस्मिन्ब्राह्मणा वशवर्तिनः} %॥२५॥

\twolineshloka
{तस्मात्प्रसीद भगवन्कृतया पूजयानया}
{पूजितः श्रावणे शुक्लत्रयोदश्यां मनोभवः} %॥२६॥

\twolineshloka
{प्रवृत्तिलम्पटस्यातिवीर्यं पुष्टिं ददात्यलम्}
{निवृत्तिमार्गनियतः स्वविकारं हरत्यपि} %॥२७॥

\twolineshloka
{सकामस्य स्त्रियो रम्याः पीनोत्तुङ्गपयोधराः}
{शरत्पूर्णसुधारश्मिवदनाः कमलेक्षणाः} %॥२८॥

\twolineshloka
{लम्बातिनीलकुरलस्निग्धकेश्यः सुनासिकाः}
{रम्भोरूर्वा गुप्तगुल्फा गतिनिर्जितकुञ्जराः} %॥२९॥

\twolineshloka
{कामागारा जिताश्वत्थपलाशा अतिशोभनाः}
{बृहच्छ्रोण्यः कम्बुकण्ठ्यो बृहज्जघनशोभिताः} %॥३०॥

\twolineshloka
{बिम्बोष्ठ्यः सिंहकट्यश्च नानालङ्कारभूषिताः}
{मनोरमा ददात्येष सन्तुष्टः श्रावणेऽर्चया} %॥३१॥

\twolineshloka
{शुक्लपक्षे त्रयोदश्यां ददाति च सुतान्बहून्}
{चिरायुषो गुणाढ्यांश्च सुखरूपान्सुसन्ततीन्} %॥३२॥

\twolineshloka
{कर्तव्यं यत् त्रयोदश्यामेतत्ते कथितं शुभम्}
{अतः परं चतुर्दश्यां कर्तव्यं शृणु मानद} %॥३३॥

\twolineshloka
{अष्टम्यां कथितं देव्याः पवित्रारोपणं तव}
{तत्र चेन्न कृतं तर्हि चतुर्दश्यां तु कारयेत्} %॥३४॥

\twolineshloka
{पवित्रं तु त्रिनेत्रस्य चतुर्दश्यां समर्पयेत्}
{पवित्रसाधनं सर्वं देवीविष्णुपवित्रवत्} %॥३५॥

\twolineshloka
{ऊहः परं प्रकर्तव्यः प्रार्थनादिषु नामसु}
{शैवागमे मया प्रोक्तं जाबालादिषु यत्परम्} %॥३६॥

\twolineshloka
{विकल्पात्कश्चिदस्तीह विशेषस्तं वदामि ते}
{एकादशाथ वा सूत्रैस्त्रिंशता चाष्टयुक्तया} %॥३७॥

\twolineshloka
{पञ्चाशता वा कर्तव्यं तुल्यग्रन्थ्यन्तरालकम्}
{द्वादशाङ्गुलमानानि तथा चाष्टाङ्गुलानि वा} %॥३८॥

\twolineshloka
{लिङ्गविस्तारमानानि चतुरङ्गुलिकानि वा}
{अर्पयेच्छिवतुष्ट्यर्थं विधिः पूर्वोक्त एव हि} %॥३९॥

\twolineshloka
{फलादि पूर्वमेवोक्तमन्ते कैलासमाप्नुयात्}
{एतत्ते कथितं वत्स किमन्यच्छ्रोतुमिच्छसि} %॥४०॥


{॥इति श्रीस्कन्दपुराणे ईश्वरसनत्कुमारसंवादे श्रावणमासमाहात्म्ये त्रयोदशीचतुर्दशीकर्तव्यकथनं नाम विंशोऽध्यायः॥२०॥}




\sect{त्रयोविंशोऽध्यायः}

\uvacha{ईश्वर उवाच}
\twolineshloka
{कृष्णाष्टम्यां नभोमासि वृषे चन्द्रे निशीथके}
{देवक्यजीजनत्कृष्णं योगेऽस्मिन्वसुदेवतः} %॥१॥

\twolineshloka
{सिंहराशिगते सूर्ये कर्तव्यः सुमहोत्सवः}
{सप्तम्यां लघुभुक्कुर्याद्दन्तधावनपूर्वकम्} %॥२॥

\twolineshloka
{उपवासस्य नियमं स्वपेद्रात्रौ जितेन्द्रियः}
{केवलेनोपवासेन कृष्णजन्मदिनं नयेत्} %॥३॥

\twolineshloka
{सप्तजन्मकृतात्पापान्मुच्यते नात्र संशयः}
{उपावृत्तस्तु पापेभ्यो यस्तु वासो गुणैः सह} %॥४॥

\twolineshloka
{उपवासः स विज्ञेयः सर्वभोगविवर्जितः}
{ततोऽष्टम्यां तिलैः स्नात्वा नद्यादौ विमले जले} %॥५॥

\twolineshloka
{सुदेशे शोभनं कुर्याद्देवक्याः सूतिकागृहम्}
{नानावर्णैः सुवासोभिः शोभितं कलशैः फलैः} %॥६॥

\twolineshloka
{पुष्पैर्दीपावलीभिश्च चन्दनागरुधूपितम्}
{हरिवंशस्य चरितं गोकुलं तत्र लेखयेत्} %॥७॥

\twolineshloka
{युक्तं वादित्रनिनदैर्नृत्यगीतादिमङ्गलैः}
{षष्ट्या देव्याधिष्ठितां च तन्मध्ये प्रतिमां हरेः} %॥८॥

\twolineshloka
{काञ्चनीं राजतीं ताम्रीं पैत्तलीं मृन्मयीं तु वा}
{वार्क्षी मणिमयीं वापि वर्णकैर्लिखितां यथा} %॥९॥

\twolineshloka
{सर्वलक्षणसम्पन्नां पर्यङ्के चाष्टशल्यके}
{प्रसूतां देवकीं तत्र स्थापयेन्मञ्चकोपरि} %॥१०॥

\twolineshloka
{सुप्तं बालं तत्र हरिं पर्यङ्के स्तनपायिनम्}
{यशोदां तत्र चैकस्मिन्प्रदेशे सूतिकागृहे} %॥११॥

\twolineshloka
{प्रसूतां कन्यकां चैव कृष्णपार्श्वे तु संलिखेत्}
{कृताञ्जलिपुटान्देवान्यक्षविद्याधरामरान्} %॥१२॥

\twolineshloka
{वसुदेवं च तत्रैव खड्गचर्मधरं स्थितम्}
{कश्यपो वसुदेवोऽयमदितिश्चैव देवकी} %॥१३॥

\twolineshloka
{शेषो बलो यशोदापि अदित्यंशाद् बभूव ह}
{नन्दः प्रजापतिर्दक्षो गर्गश्चापि चतुर्मुखः} %॥१४॥

\twolineshloka
{गोप्यश्चाप्सरसः सर्वा गोपाश्चापि दिवौकसः}
{कालनेमिश्च कंसोऽयं नियुक्तास्तेन चासुराः} %॥१५॥

\twolineshloka
{गोधेनुकुञ्जराश्वाश्च दानवाः शस्त्रपाणयः}
{लेखनीयाश्च तत्रैव कालियो यमुनाहृदे} %॥१६॥

\twolineshloka
{इत्येवमादि यत्किञ्चिच्चरितं हरिणा कृतम्}
{लेखयित्वा प्रयत्नेन पूजयेद्भक्तितत्परः} %॥१७॥

\onelineshloka
{उपचारैः षोडशभिर्देवकी चेति मन्त्रतः} %॥१८॥

\fourlineindentedshloka
{गायद्भिः किन्नराद्यैः सततपरिणुता वेणुवीणानिनादैर्}
{भृङ्गार्यादर्शदूर्वादधिकलशकरैः किन्नरैः सेव्यमाना}
{पर्यङ्के स्वासनस्था मुदिततरमुखी पुत्रिणी सम्यगास्ते}
{सा देवी दिव्यमाता विजयसुतसुता देवकी कान्तयुक्ता} %॥१९॥

\twolineshloka
{प्रणवादिनमोन्तैश्च पृथङ्नामानुकीर्तनैः}
{कुर्यात्पूजां विधिज्ञस्तु सर्वपापापनुत्तये} %॥२०॥

\twolineshloka
{देवक्या वसुदेवस्य वासुदेवस्य चैव हि}
{बलदेवस्य नन्दस्य यशोदायाः पृथक् पृथक्} %॥२१॥


\threelineshloka
{चन्द्रोदये शशाङ्काय अर्घ्यं दद्याद्धरिं स्मरन्}
{क्षीरोदार्णवसम्भूत अत्रिगोत्रसमुद्भव}
{नमस्ते रोहिणीकान्त अर्घ्यं नः प्रतिगृह्यताम्} %२२

\twolineshloka
{देवक्या वसुदेवं च नन्दं चैव यशोदया}
{रोहिण्या च सुधारश्मिं बलं च हरिणा सह} %॥२३॥

\twolineshloka
{सम्पूज्य विधिवद्देही किं नाप्नोति सुदुर्लभम्}
{एकादशीकोटिसङ्ख्यातुल्या कृष्णाष्टमी तथा} %॥२४॥

\twolineshloka
{एवं सम्पूज्य तद्रात्रौ प्रभाते नवमीदिने}
{यथा हरेस्तथा कार्यो भगवत्या महोत्सवः} %॥२५॥

\twolineshloka
{ब्राह्मणान्भोजयेद्भक्त्या दद्याद्वै गोधनादिकम्}
{यद्यदिष्टतमं तत्र कृष्णो मे प्रीयतामिति} %॥२६॥

\twolineshloka
{नमस्ते वासुदेवाय गोब्राह्मणहिताय च}
{शान्तिरस्तु शिवं चास्तु इत्युक्त्वा तं विसर्जयेत्} %॥२७॥

\twolineshloka
{ततो बन्धुजनैः सार्धं स्वयं भुञ्जीत वाग्यतः}
{एवं यः कुरुते देव्याः कृष्णस्य च महोत्सवम्} %॥२८॥

\twolineshloka
{प्रतिवर्षं विधानेन यथोक्तं लभते फलम्}
{पुत्रसन्तानमारोग्यं सौभाग्यमतुलं भवेत्} %॥२९॥

\twolineshloka
{इह धर्ममतिर्भूत्वा अन्ते वैकुण्ठमाप्नुयात्}
{उद्यापनमथो वक्ष्ये पुण्येऽह्नि विधिपूर्वकम्} %॥३०॥

\twolineshloka
{पूर्वेद्युरेकभक्ताशी स्वपेद्विष्णुं स्मरन्हृदि}
{प्रातः सन्ध्यादि सम्पाद्य ब्राह्मणैः स्वस्ति वाचयेत्} %॥३१॥

\twolineshloka
{आचार्यं वरयित्वा तु ऋत्विजश्चैव पूजयेत्}
{पलेन वा तदर्धेन तदर्धार्धेन वा पुनः} %॥३२॥

\twolineshloka
{प्रतिमां कारयेत्पश्चाद्वित्तशाठ्यविवर्जितः}
{मण्डपे मण्डले देवान्ब्रह्माद्यान्स्थापयेद् बुधः} %॥३३॥

\twolineshloka
{तत्र संस्थापयेत्कुम्भं ताम्रं मृण्मयमेव वा}
{तस्योपरि न्यसेत्पात्रं राजतं वैणवं तु वा} %॥३४॥

\twolineshloka
{वाससाच्छाद्य गोविन्दं तत्र सम्पूजयेद् बुधः}
{उपचारैः षोडशभिर्मन्त्रैर्वैदिकतान्त्रिकैः} %॥३५॥

\twolineshloka
{ततोऽर्घ्यं हरये दद्याद्देवकीसहिताय च}
{शङ्खे कृत्वा जलं शुद्धं सपुष्पफलचन्दनम्} %॥३६॥

\twolineshloka
{जानुभ्यामवनीं गत्वा नारिकेलफलान्वितम्}
{जातः कंसवधार्थाय भूभारोत्तारणाय च} %॥३७॥

\twolineshloka
{कौरवाणां विनाशाय दैत्यानां निधनाय च}
{गृहाणार्घ्यं मया दत्तं देवक्या सहितो हरे} %॥३८॥

\onelineshloka
{चन्द्रायार्घ्यं ततो दद्यात्पूर्वोक्तविधिना सुधीः} %॥३९॥

\twolineshloka
{नमस्तुभ्यं देवकीतनय प्रभो}
{वसुदेवात्मजानन्त त्राहि मां भवसागरात्} %॥४०॥

\twolineshloka
{इत्थं सम्प्रार्थ्य देवेशं रात्रौ जागरणं चरेत्}
{प्रत्यूषे विमले स्नात्वा पूजयित्वा जनार्दनम्} %॥४१॥

\twolineshloka
{मन्त्रेणेदं विष्णुरिति जुहुयाद्वै घृताहुतीः}
{होमशेषं समाप्याथ पूर्णाहुतिपुरःसरम्} %॥४३॥

\twolineshloka
{आचार्यं पूजयेत्पश्चाद्भूषणाच्छादनादिभिः}
{गामेकां कपिलां दद्याद् व्रतसम्पूर्णहेतवे} %॥४४॥

\twolineshloka
{पयस्विनीं सुशीलां च सवत्सां सगुणां तथा}
{स्वर्णशृङ्गीं रौप्यखुरां कांस्यदोहनिकायुताम्} %॥४५॥

\twolineshloka
{मुक्तापुच्छां ताम्रपृष्ठीं स्वर्णघण्टासमन्विताम्}
{वस्त्राच्छन्नां दक्षिणाढ्यामेवं सम्पूर्णतामियात्} %॥४६॥

\twolineshloka
{कपिलाया अभावेऽपि गौरन्यापि प्रदीयते}
{ततः प्रदद्यादृत्विग्भ्यो दक्षिणां च यथार्हतः} %॥४७॥

\twolineshloka
{ब्राह्मणान्भोजयेत्पश्चादष्टौ तेभ्यश्च दक्षिणाम्}
{कलशान्जलसम्पूर्णान्दद्याच्चैव समाहितः} %॥४८॥

\twolineshloka
{प्राप्यानुज्ञां तथा तेभ्यो भुञ्जीत सह बन्धुभिः}
{एवं कृते ब्रह्मपुत्र व्रतोद्यापनकर्मणि} %॥४९॥

\twolineshloka
{पायसेन तिलाज्यैश्च मूलमन्त्रेण भक्तितः}
{अष्टोत्तरशतं हुत्वा ततः पुरुषसूक्ततः} %॥४२॥


\threelineshloka
{निष्पापस्तत्क्षणादेव जायते विबुधोत्तमः}
{पुत्रपौत्रसमायुक्तो धनधान्यसमन्वितः}
{भुक्त्वा भोगांश्चिरं कालमन्ते वैकुण्ठमाप्नुयात्} %५०


{॥इति श्रीस्कन्दपुराणे ईश्वरसनत्कुमारसंवादे श्रावणमासमाहात्म्ये कृष्णजन्माष्टमीव्रतकथनं नाम त्रयोविंशोऽध्यायः॥२३॥}




\sect{चतुर्विंशोऽध्यायः}

\uvacha{ईश्वर उवाच}
\twolineshloka
{पुरा कल्पे ब्रह्मपुत्र दैत्यभारप्रपीडिता}
{ब्रह्माणं शरणं प्राप पृथ्वी दीनातिविह्वला} %॥१॥

\twolineshloka
{वृत्तान्तं तन्मुखाच्छ्रुत्वा ब्रह्मा देवगणैः सह}
{क्षीरार्णवे हरिं गत्वा तुष्टाव स्तुतिभिर्बहु} %॥२॥

\twolineshloka
{प्रादुरासीत्ततो दिक्षु श्रुत्वा सर्वं विधेर्मुखात्}
{मा भैष्ट देवा देवक्या जठरे वसुदेवतः} %॥३॥

\twolineshloka
{अवतीर्णो भविष्यामि हरिष्ये भूमिवेदनाम्}
{भवन्तु यादवा देवा इत्युक्त्वान्तर्दधे विभुः} %॥४॥

\twolineshloka
{देवक्या जठरे जातो वसुदेवेन गोकुले}
{स्थापितः कंसभीतेन ववृधे तत्र कंसहा} %॥५॥

\twolineshloka
{आगत्य मथुरां पश्चात्कंसं सगणमाहनत्}
{ततः सर्वे पौरजनाः प्रार्थयामासुरादरात्} %॥६॥

\twolineshloka
{कृष्ण कृष्ण महायोगिन्भक्तानामभयप्रद}
{प्रणतान्पाहि नो देव शरणागतवत्सल} %॥७॥

\twolineshloka
{किञ्चिद्विज्ञापये देव एतन्नो वक्तुमर्हसि}
{तव जन्मदिने कृत्यं न ज्ञातं केनचित् क्वचित्} %॥८॥

\twolineshloka
{ज्ञात्वा च तद्दिने सर्वे कुर्मो वर्धापनोत्सवम्}
{तेषां दृष्ट्वा च तां भक्तिं स्वस्मिञ्छ्रद्धां च सौहृदम्} %॥९॥

\twolineshloka
{कृत्यं जन्मदिने तेभ्यः कथयामास केशवः}
{श्रुत्वा तेऽपि तथा चक्रुर्विधानात्तेन तद् व्रतम्} %॥१०॥

\twolineshloka
{वरांश्च बहुधा प्रादाद् भगवान्व्रतकारिणे}
{अत्रैवोदाहरन्तीममितिहासं पुरातनम्} %॥११॥

\twolineshloka
{अङ्गदेशोद्भवो राजा मितजिन्नाम नामतः}
{तस्य पुत्रो महासेनः सत्यजित्सत्पथे स्थितः} %॥१२॥

\twolineshloka
{पालयामास सर्वज्ञो विधिवद्रञ्जयन्प्रजाः}
{तस्यैवं वर्तमानस्य कदाचिद्दैवयोगतः} %॥१३॥

\twolineshloka
{पाखण्डैः सह संवासो बभूव बहुवासरम्}
{तत्संसर्गात्स नृपतिरधर्मे निरतोऽभवत्} %॥१४॥

\twolineshloka
{वेदशास्त्रपुराणानि निनिन्द बहुशो नृपः}
{वर्णाश्रमगते धर्मे विद्वेषं परमं गतः} %॥१५॥

\twolineshloka
{एवं बहुतिथे काले प्रयाते मुनिसत्तम}
{कालेन निधनं प्राप्तो यमदूतवशं गतः} %॥१६॥

\twolineshloka
{बध्वा पाशैर्नीयमानो यमदूतैर्यमान्तिकम्}
{पीडितस्ताड्यमानोऽसौ दुष्टसङ्गतियोगतः} %॥१७॥

\twolineshloka
{नरके पातितः प्राप यातना बहुवत्सरम्}
{भुक्त्वा पापस्य शेषेण पैशाचीं योनिमास्थितः} %॥१८॥

\twolineshloka
{क्षुधातृष्णासमाक्रान्तो भ्रमन्स मरुधन्वसु}
{कस्यचित्त्वथ वैश्यस्य देहमाविश्य संस्थितः} %॥१९॥

\twolineshloka
{सह तेनैव संयातो मथुरां पुण्यदां पुरीम्}
{समीपे रक्षकैस्तस्य तस्माद् गेहाद् बहिष्कृतः} %॥२०॥

\twolineshloka
{बभ्राम विपिने सोऽथ ऋषीणामाश्रमेषु च}
{कदाचिद्दैवयोगेन हरेर्जन्माष्टमीदिने} %॥२१॥

\twolineshloka
{क्रियमाणां महापूजां व्रतिभिर्मुनिभिर्द्विजैः}
{रात्रौ जागरणं चैव नामसङ्कीर्तनादिभिः} %॥२२॥

\twolineshloka
{ददर्श सर्वं विधिवच्छुश्राव च हरेः कथाम्}
{निष्पापस्तत्क्षणादेव शुद्धो निर्मलमानसः} %॥२३॥

\twolineshloka
{प्रेतदेहं समुत्सृज्य विष्णुलोके विमानगः}
{यमदूतैः परित्यक्तो दिव्यभोगसमन्वितः} %॥२४॥

\twolineshloka
{विष्णुसान्निध्यमापन्नो व्रतस्यास्य प्रभावतः}
{नित्यमेतद् व्रतं चैव पुराणे सार्वलौकिकम्} %॥२५॥

\twolineshloka
{कथ्यते विधिवत्सम्यङ् मुनिभिस्तत्त्वदर्शिभिः}
{सर्वकामिकमेवैतत्कृत्वा कामानवाप्नुयात्} %॥२६॥

\twolineshloka
{एवं यः कुरुते कृष्णजन्माष्टम्यां व्रतं शुभम्}
{भुक्त्वेह विविधान्भोगाञ्छुभान्कामानवाप्नुयात्} %॥२७॥

\twolineshloka
{तत्र देवविमानेन वर्षलक्षं विधेः सुत}
{भोगान्नानाविधान्भुक्त्वा पुण्यशेषादिहागतः} %॥२८॥

\twolineshloka
{सर्वकामसमृद्धस्तु सर्वाशुभविवर्जितः}
{कुले नृपतिवर्याणां जायते मदनोपमः} %॥२९॥

\twolineshloka
{यस्मिन्सदैव विषये लिखितं स्यात्परार्पितम्}
{कृष्णजन्मोपकरणं सर्वशोभासमन्वितम्} %॥३०॥

\twolineshloka
{पूज्यते विश्वसृद् तत्र व्रतैरुत्सवसंयुतैः}
{परचक्रभयं तत्र न कदाचिद्भविष्यति} %॥३१॥

\twolineshloka
{पर्जन्यः कामवर्षी स्यादीतिभ्यो न भयं क्वचित्}
{गृहे वा पूजयेद्यस्तु चरितं देवकीजनः} %॥३२॥


\threelineshloka
{तत्र सर्वसमृद्धं स्यान्नोपसर्गाद्भयं भवेत्}
{संसर्गेणापि यो भक्त्या व्रतं पश्येदनाकुलः}
{सोऽपि पापविनिर्मुक्तः प्रयाति हरिमन्दिरम्} %३३


{॥इति श्रीस्कन्दपुराणे ईश्वरसनत्कुमारसंवादे श्रावणमासमाहात्म्ये जन्माष्टमीव्रतकथनं नाम चतुर्विंशोऽध्यायः॥२४॥}




\sect{पञ्चविंशोऽध्यायः}

\uvacha{ईश्वर उवाच}
\twolineshloka
{वक्ष्ये मुनिश्रेष्ठ पिठोरीव्रतमुत्तमम्}
{अमायां श्रावणे मासि सर्वसम्पत्प्रदायकम्} %॥१॥

\twolineshloka
{सर्वाधिष्ठानमेतद्यद्गृहं पीठं ततो मतम्}
{आरस्तत्रसमूहः स्याद्वस्तुमात्रस्य पूजने} %॥२॥

\twolineshloka
{पिठोरमितिसंज्ञा च व्रतस्यातो मुनीश्वर}
{तत्प्रकारं च वक्ष्येऽहं सावधानमनाः शृणु} %॥३॥

\twolineshloka
{कुड्ये विलिप्य ताम्रेण कृष्णेनाथ सितेन वा}
{धातुना तत्र ताम्रे तु पीतेन विलिखेत्सुधीः} %॥४॥

\twolineshloka
{शुक्लेन वाथ कृष्णेन पूर्ववच्चैव संलिखेत्}
{सितपीतेन रक्तेन कृष्णेन हरितेन वा॥५} %॥५॥

\twolineshloka
{मध्ये शिवं शिवायुक्तं लिङ्गं वा मूर्तिमेव वा}
{विस्तीर्णं कुड्यमालिख्य सर्वसंसारमालिखेत्} %॥६॥

\twolineshloka
{चतुःशालासमायुक्तं पाकागारं सुरालयम्}
{शय्यागृहं सप्तकोशांस्तथान्तः स्त्रीनिकेतनम्} %॥७॥

\twolineshloka
{प्रासादाट्टालिकाशोभं शालवृक्षसमुद्भवम्}
{इष्टकाभिश्च पाषाणैश्चूर्णनद्धैः सुशोभनम्} %॥८॥

\twolineshloka
{द्वाराणि च विचित्राणि वलभीचेष्टिकास्तथा}
{अजा गावो महिष्यश्च अश्वा उष्ट्रा मतङ्गजाः} %॥९॥

\twolineshloka
{गन्त्रीरथप्रभृतयः शकटानां प्रभेदकाः}
{स्त्रियो बालाश्च वृद्धाश्च तरुण्यः पुरुषास्तथा} %॥१०॥

\onelineshloka
{पालक्यान्दोलिका चैव मञ्चका बहुरूपकाः} %॥११॥

\twolineshloka
{हैमानि रौप्याणि च ताम्रकाणि सैसानि लौहानि च मृन्मयानि}
{रङ्गप्रसूतानि च पैत्तलानि पात्राणि नानाविधकारकाणि} %॥१२॥

\twolineshloka
{यावन्तः कशिपुभेदा उपबर्हणजातयः}
{मार्जाराः सारिकाश्चैव शुभा अन्येऽपि पक्षिणः} %॥१३॥

\twolineshloka
{पुरुषाणामलङ्काराः स्त्रीणां चैवाप्यनेकशः}
{यानि चास्तरणानीह तथा प्रावरणानि च} %॥१४॥


\threelineshloka
{यज्ञपात्राणि यावन्ति स्तम्भदण्डौ च मन्थने}
{रज्जुत्रयं च तद्धेतु दुग्धं च नवनीतकम्}
{दधि तक्रं तथा मस्तु आज्यं तैलं तिलांस्तथा} %१५

\twolineshloka
{गोधूमशालितुवरीयवयावनालवार्तानलञ्च चणका मसुराः कुलित्थाः}
{मुद्गप्रियङ्गतिलकोद्रवकातसीतिश्यामाकमाषचवला इति धान्यवर्गाः} %॥१६॥

\twolineshloka
{दृषदं चोपलं चुल्लीं तथा सम्मार्जनीमपि}
{पुरुषाणां च वस्त्राणि नारीणां चैव सर्वशः} %॥१७॥

\twolineshloka
{वेणुजन्यं च शूर्पादि तथा तृणभवानि च}
{उलूखलं च मुसलं यन्त्रं दलयुगान्वितम्} %॥१८॥

\twolineshloka
{व्यजनं चामरं छत्रमुपानत्पादुकाद्वयम्}
{दास्यो दास्या भृत्यपोष्याः पशुभक्ष्यं तृणादिकम्} %॥१९॥

\twolineshloka
{धनुर्बाणशतघ्न्यश्च खड्गाः कुन्ताश्च शक्तयः}
{चर्मपाशाङ्कुशगदास्त्रिशूलं भिन्दिपालकाः} %॥२०॥

\twolineshloka
{तोमरो मुद्गरश्चैव परशुः पट्टिशस्तथा}
{भुशुण्डी परिघश्चैव चक्रयन्त्रादिकं च यत्} %॥२१॥

\twolineshloka
{जलयन्त्रं मषीपात्रं लेखनी पुस्तकादिकम्}
{फलजातं सर्वमपि छुरिका कर्तरी तथा} %॥२२॥

\twolineshloka
{नानाविधानि पुष्पाणि बिल्वश्च तुलसी तथा}
{दीपिकाश्चैव दीपाश्च तथा तत्साधनानि च} %॥२३॥

\twolineshloka
{शाकं नानाविधं भक्ष्यं पक्वान्नानां च या भिदः}
{लेख्यं तच्चैव सकलमनूक्तमपि चैव हि} %॥२४॥

\twolineshloka
{कियल्लेख्यं जनेनात्र वक्तव्यं वा मया कियत्}
{एकैकस्य पदार्थस्य भेदाः शतसहस्रशः} %॥२५॥

\twolineshloka
{उपचारैः षोडशभिः सर्वेषां पूजनं भवेत्}
{नानाविधश्च गन्धः स्यात्पुष्पधूपोऽपि चन्दनम्} %॥२६॥

\twolineshloka
{ब्राह्मणान्भोजयेद् बालान्सुवासिन्यश्च पुष्कलान्}
{प्रार्थयेच्च शिवं साम्बं व्रतं सम्पूर्णमस्त्विति} %॥२७॥

\twolineshloka
{शिव साम्ब दयासिन्धो गिरीश शशिशेखर}
{व्रतेनानेन सन्तुष्टः प्रयच्छास्मान्मनोरथान्} %॥२८॥

\twolineshloka
{एवं कृत्वा पञ्चवर्षं तत उद्यापनं चरेत्}
{आज्येन बिल्वपत्रैश्च होमः स्याच्छिवमन्त्रतः} %॥२९॥

\twolineshloka
{ग्रहहोमः पुरा कार्यः पूर्वेद्युरधिवासनम्}
{अष्टोत्तरसहस्रं वा शतमष्टोत्तरं तु वा} %॥३०॥

\twolineshloka
{होमसङ्ख्या भवेद्वत्स आचार्यं पूजयेत्ततः}
{भूयसीं दक्षिणां दद्यात्स्वयं भोजनमाचरेत्} %॥३१॥


\threelineshloka
{इष्टबन्धुजनैः सार्धं कुटुम्बसहितो बुधः}
{एवं कृते विधाने तु सर्वान्कामानवाप्नुयात्}
{यद्यदिष्टतमं लोके तत्सर्वं लभते नरः} %३२

\twolineshloka
{एतत्ते कथितं वत्स पिठोरीव्रतमुत्तमम्}
{व्रतेनानेन सदृशं सर्वकामसमृद्धिदम्} %॥३३॥

\twolineshloka
{शिवप्रीतिकरं चैव न भूतं न भविष्यति}
{भित्तौ यद्यल्लिखेद्वस्तु तत्तदाप्नोति निश्चितम्} %॥३४॥


{॥इति श्रीस्कन्दपुराणे ईश्वरसनत्कुमारसंवादे श्रावणमासमाहात्म्ये अमावास्यायां पिठोरीव्रतकथनं नाम पञ्चविंशोऽध्यायः॥२५॥}




\sect{षड्विंशोऽध्यायः}

\uvacha{ईश्वर उवाच}
\twolineshloka
{यत्कर्तव्यं नभोमासि अमावास्यादिने भवेत्}
{प्रसङ्गतश्च यच्चान्यत्स्मृतं तदपि ते ब्रुवे} %॥१॥

\twolineshloka
{पुरा नानाविधैर्दैत्यैर्महाबलपराक्रमैः}
{जगद्विध्वंसकैर्दुष्टैर्देवतोच्छेदकारिभिः } %॥२॥

\twolineshloka
{सङ्ग्रामा बहवो जाता आरुह्य वृषभं शुभम्}
{महासत्त्वो महावीर्यो न कदाचिज्जहौ च माम्} %॥३॥

\twolineshloka
{अन्धकासुरयुद्धे तु तेन छिन्नतनुः कृतः}
{भिन्नत्वग्रुधिरस्रावी प्राणमात्रावशेषितः} %॥४॥

\twolineshloka
{तथापि धैर्यमालम्ब्य यावद्धन्मि च तं खलम्}
{उवाह तावन्मां नन्दी तस्य तज्ज्ञातवाहनम्} %॥५॥

\twolineshloka
{हत्वा तमन्धकं दैत्यं तुष्टोऽहं नन्दिनं तदा}
{कर्मणा ते प्रसन्नोऽस्मि वरं वरय सुव्रत} %॥६॥

\twolineshloka
{व्रणास्ते प्रशमं यान्तु निरोगो बलवान्भव}
{पूर्वस्मादपि ते वीर्यं रूपं चापि विवर्धताम्} %॥७॥

\onelineshloka
{यं यं वरं याचसे त्वं तं तं दास्याम्यसंशयम्} %॥८॥


\uvacha{नन्दिकेश्वर उवाच}
\twolineshloka
{ममास्ति तथापि याचनीयं न देवदेव महेश्वर}
{ममोपरि प्रसन्नोऽसि किं वैभवमतः परम्} %॥९॥

\twolineshloka
{तथापि भगवन् याचे लोकोपकृतये शिव}
{अद्यामा श्रावणस्यास्ति यस्यां तुष्टो भवान्मम} %॥१०॥

\twolineshloka
{एतस्यां वृषभाः पूज्या गोभिर्युक्ताः सुमृन्मयाः}
{अद्यैवामादिने जन्म कामधेनूपमं भवेत्} %॥११॥

\twolineshloka
{अतोऽप्यस्यां वरं देहि भवत्वेषेच्छितप्रदा}
{प्रत्यक्षं वृषभा गावः पूजनीयाश्च भक्तितः} %॥१२॥

\twolineshloka
{धातुभिर्गैरिकाद्यैश्च भूषणीयाः प्रयत्नतः}
{शृङ्गेषु स्वर्णरौप्यादिपट्टिकाबन्धशोभनम्} %॥१३॥

\twolineshloka
{कौशेयगुच्छान्महतः शृङ्गयोरपि बन्धयेत्}
{पृष्ठं नानाविधैर्वर्णैश्चित्रितेन सुवाससा} %॥१४॥

\twolineshloka
{आच्छादयेद् गले घण्टां बध्नीयाद्रम्यशब्दिताम्}
{दिनाष्टांशे बहिर्नीत्वा सायं ग्रामं प्रवेशयेत्} %॥१५॥

\twolineshloka
{पिण्यातकं च नैवेद्यं अन्नं नानाविधं च यत्}
{अर्पयेत्तस्य भवतु गोधनं वृद्धिगं सदा} %॥१६॥

\twolineshloka
{गावो यत्र गृहे न स्युः श्मशानसदृशं च तत्}
{पञ्चामृतं पञ्चगव्यं न भवेद् गोरसं विना} %॥१७॥

\twolineshloka
{सम्मार्जनं पूततमं गोमयेन विना न हि}
{पिपीलिकादिजन्तूनामुपसर्गाश्च तत्र हि} %॥१८॥

\twolineshloka
{प्रोक्षणं यत्र गोमूत्रान्न भवेच्च सुरोत्तम}
{भोजनस्य महादेव को रसो गोरसं विना} %॥१९॥

\twolineshloka
{एतेऽन्येऽपि वरा देयाः प्रसन्नोऽसि यदि प्रभो}
{इति नन्दिवचः श्रुत्वा तुष्टोऽहमधिकं तदा} %॥२०॥

\twolineshloka
{सर्वमस्तु वृषश्रेष्ठ यथा ते याचितं तथा}
{अन्यच्च शृणु भो नन्दिन्नामास्य तु दिनस्य यत्} %॥२१॥

\twolineshloka
{न वाह्यते यो वृषभः केनचित्कर्मणि क्वचित्}
{तृणमश्नन्पिबन्नीरं तूष्णीं यो वर्धते वृषः} %॥२२॥

\twolineshloka
{महावीरश्च बलवान्पोल इत्युच्यते हि सः}
{तन्नाम्नेदं दिनं नन्दिन्पोला इति भविष्यति} %॥२३॥


\threelineshloka
{तत्रोत्सवो महान् कार्य इष्टबन्धुजनैः सह}
{इति दत्ता मया वत्स वराः श्रेष्ठा हि तद्दिने}
{तेन श्रेष्ठं दिनं चैतत्पोलासंज्ञं मतं जनैः} %२४

\twolineshloka
{अत्रोत्सवो महान् कार्यो वृषाणां सर्वकामदः}
{अतः परं प्रवक्ष्यामि अस्यामेव कुशग्रहम्} %॥२५॥

\twolineshloka
{नभोमासस्य दर्शे तु शुचिर्दर्भान्समाहरेत्}
{अयातयामास्ते दर्भा विनियोज्याः पुनः पुनः} %॥२६॥

\twolineshloka
{कुशाः काशा यवा दूर्वा उशीराश्च सकूदकाः}
{गोधूमा व्रीहयो मौञ्ज्या दश दर्भाः सबल्वजाः} %॥२७॥

\twolineshloka
{विरिञ्चिना सहोत्पन्न परमेष्ठीनिसर्गज}
{नुद पापानि सर्वाणि दर्भ स्वस्तिकरो भव} %॥२८॥

\twolineshloka
{एवं सन्मन्त्रमुच्चार्य ततः पूर्वोत्तरामुखः}
{हुंफट्कारेण मन्त्रेण सकृच्छित्वा समुद्धरेत्} %॥२९॥

\twolineshloka
{अच्छिन्नाग्रा अशुष्काग्राः पैत्र्ये तु हरिताः स्मृताः}
{अमूला देवकार्येषु प्रयोज्याश्च जपादिषु} %॥३०॥

\twolineshloka
{सप्तपत्राः कुशाः शस्ता दैवे पित्र्ये च कर्मणि}
{अनन्तर्गर्भिणौ साग्रौ प्रादेशौ च पवित्रके} %॥३१॥

\twolineshloka
{चतुर्भिर्दर्भपिञ्जलैः पवित्रं ब्राह्मणस्य तु}
{एकैकं न्यूनमुद्दिष्टं वर्णे वर्णे यथाक्रमम्} %॥३२॥

\twolineshloka
{सर्वेषां वा भवेद् द्वाभ्यां पवित्रं ग्रन्थिशोभितम्}
{इदं तु धारणार्थं स्यात्पवित्रं कथितं तव} %॥३३॥

\twolineshloka
{दर्भद्वयं तु सर्वेषां भवेदुत्पवनाय च}
{पञ्चाशता भवेद् ब्रह्मा तदर्धेन तु विष्टरः} %॥३४॥

\twolineshloka
{निष्कासनीयं नो हस्तादाचमे तु पवित्रकम्}
{विकिरेऽग्नौ कृते चैव कृते पाद्ये तु सन्त्यजेत्} %॥३५॥

\twolineshloka
{नास्ति दर्भसमं पुण्यं पवित्रं पापनाशनम्}
{दर्भाधीनानि कर्माणि दैवपित्र्याणि सर्वशः} %॥३६॥

\twolineshloka
{तादृग्विधानां दर्भाणाममायां ग्रहणं भवेत्}
{अयातयामता चैव किं वर्ण्यामा नभस्यतः} %॥३७॥

\twolineshloka
{इत्येतत्कथितं कृत्यममायां श्रावणे तु यत्}
{अन्यच्च श्रावणे कृत्यं तच्चापि कथयामि ते} %॥३८॥


{॥इति श्रीस्कन्दपुराणे ईश्वरसनत्कुमारसंवादे श्रावणमासमाहात्म्ये अमायां वृषभपूजनं कुशग्रहणं नाम षड्विंशोऽध्यायः॥२६॥}




\sect{सप्तविंशोऽध्यायः}

\uvacha{ईश्वर उवाच}
\twolineshloka
{अथातः श्रावणे कर्कसिंहसङ्क्रान्तिसम्भवः}
{प्राप्यते तत्र यत्कृत्यं तच्चापि कथयामि ते} %॥१॥

\twolineshloka
{सिंहकर्कटयोर्मध्ये सर्वा नद्यो रजस्वलाः}
{तासु स्नानं न कुर्वीत वर्जयित्वा समुद्रगाः} %॥२॥

\twolineshloka
{अगस्त्योदयपर्यन्तं केचिदूचुर्महर्षयः}
{यावन्नोदेति भगवान्दक्षिणाशाविभूषणः} %॥३॥

\twolineshloka
{तावद्रजोवहा नद्य अल्पतोयाः प्रकीर्तिताः}
{याः शोषमुपगच्छन्ति ग्रीष्मे तु सरितो भुवि} %॥४॥

\twolineshloka
{तासु प्रावृषि न स्नायादपूर्णे दशवासरे}
{धनुःसहस्राण्यष्टौ च गतिर्यासां स्वतो नहि} %॥५॥

\twolineshloka
{न ता नदीशब्दवाच्या गर्तास्ते परिकीर्तिताः}
{प्रारम्भे कर्कसङ्क्रान्तेर्महानद्यो रजस्वलाः} %॥६॥

\twolineshloka
{त्रिदिनं तु चतुर्थेऽह्नि शुद्धाः स्युर्योषितो यथा}
{महानदीः प्रवक्ष्यामि शृणुष्वावहितो मुने} %॥७॥

\twolineshloka
{गोदावरी भीमरथी तुङ्गभद्रा च वेणिका}
{तापी पयोष्णी विन्ध्यस्य दक्षिणे षट् प्रकीर्तिताः} %॥८॥

\twolineshloka
{भागीरथी नर्मदा च यमुना च सरस्वती}
{विशोका च वितस्ता च विन्ध्यस्योत्तरतोऽपि षट्} %॥९॥

\twolineshloka
{द्वादशैता महानद्यो देवर्षिक्षेत्रसम्भवाः}
{महानद्यो देविका च कावेरी वञ्जरा तथा} %॥१०॥

\twolineshloka
{कृष्णा रजस्वला एताः कर्कटादावहर्मुने}
{कर्कटादौ रजोदुष्टा गौतमी वासरत्रयम्} %॥११॥

\twolineshloka
{चन्द्रभागा सती सिन्धुः शरयूर्नर्मदा तथा}
{गङ्गा च यमुना चैव प्लक्षजाला सरस्वती} %॥१२॥

\twolineshloka
{रजसा नाभिभूयन्ते ये चान्ये नदसंज्ञिताः}
{शोणः सिन्धुर्हिरण्याख्यः कोकिलाहितघर्घराः} %॥१३॥

\twolineshloka
{शतद्रुश्च नदाः सप्त पावनाः परिकीर्तिताः}
{गङ्गा धर्मद्रवा पुण्या यमुना च सरस्वती} %॥१४॥

\twolineshloka
{अन्तर्गता रजोदोषाः सर्वावस्थासु चामलाः}
{अपामयं रजोदोषो न भवेत्तीरवासिनाम्} %॥१५॥

\twolineshloka
{जलं रजोदुष्टमपि गङ्गातोयेन पावनम्}
{अजा गावो महिष्यश्च योषितश्च प्रसूतिकाः} %॥१६॥

\twolineshloka
{भूमेर्नवोदकं चैव दशरात्रेण शुध्यति}
{अभावे कूपवापीनामन्यासां च पयोऽमृतम्} %॥१७॥

\twolineshloka
{रजोदुष्टेऽपि वयसि ग्रामभोगो न दुष्यति}
{अन्येन चोद्धृते नीरे रजोदोषो न विद्यते} %॥१८॥

\twolineshloka
{उपाकर्मणि चोत्सर्गे प्रातः स्नाने विपत्सु च}
{चन्द्रसूर्यग्रहे चैव रजोदोषो न विद्यते} %॥१९॥

\twolineshloka
{अतः परं प्रवक्ष्यामि सिंहे गोप्रसवो यदि}
{भानौ सिंहगते चैव यस्य गौः सम्प्रसूयते} %॥२०॥

\twolineshloka
{मरणं तस्य निर्दिष्टं षड्भिर्मासैर्न संशयः}
{तत्र शान्तिं प्रवक्ष्यामि येन सम्पाद्यते सुखम्} %॥२१॥

\twolineshloka
{प्रसूतां तत्क्षणादेव तां गां विप्राय दापयेत्}
{ततो होमं प्रकुर्वीत घृताक्तै राजसर्षपैः} %॥२२॥

\twolineshloka
{आहुतीनां घृताक्तानां तिलानां जुहुयात्ततः}
{सहस्रेण व्याहृतिभिरष्टसङ्ख्याधिकेन च} %॥२३॥

\twolineshloka
{सोपवासः प्रयत्नेन दद्याद् विप्राय दक्षिणाम्}
{सिंहराशौ गते सूर्ये गोष्ठसूतिर्यदा भवेत्} %॥२४॥

\twolineshloka
{तदानिष्टं भवेत्किञ्चित्तच्छान्त्यै शान्तिकं चरेत्}
{अस्य वामेति सूक्तेन तद्विष्णोरिति मन्त्रतः} %॥२५॥

\twolineshloka
{जुहुयाच्च तिलाज्येन शतमष्टोत्तराधिकम्}
{मृत्युञ्जयविधानेन जुहुयाच्च तथायुतम्} %॥२६॥

\twolineshloka
{श्रीसूक्तेन ततः स्नायाच्छान्तिसूक्तेन वा पुनः}
{एवं कृतविधानेन न भयं जायते क्वचित्} %॥२७॥

\twolineshloka
{एवमेव नभोमासि सूयेत वडवा दिने}
{अत्रापि शान्तिकं कार्यं तदा दोषो विनश्यति} %॥२८॥

\twolineshloka
{कर्के सिंहे नभोदानमथ वक्ष्ये शुभप्रदम्}
{घृतधेनुप्रदानं च कर्कटस्थे दिवाकरे} %॥२९॥

\twolineshloka
{ससुवर्णं छत्रदानं शस्तं सिंहे निगद्यते}
{श्रावणे वस्त्रदानस्य कीर्तितं सुमहत्फलम्} %॥३०॥

\twolineshloka
{घृतं च घृतकुम्भाश्च घृतधेनुः फलानि च}
{श्रावणे श्रीधरप्रीत्यै दातव्यानि विपश्चिते} %॥३१॥

\twolineshloka
{अन्यान्यपि च दानानि मत्तोषाय कृतानि च}
{अक्षय्यफलदानि स्युरन्यमासेभ्य एव हि} %॥३२॥

\twolineshloka
{द्वादशस्वपि मासेषु नास्ति चैतादृशः प्रियः}
{आगच्छति नभोमासि प्रतीक्षां च करोम्यहम्} %॥३३॥

\twolineshloka
{करिष्यते व्रतं योऽत्र स मे प्रियतरो भवेत्}
{ब्राह्मणानां विधू राजा सूर्यः प्रत्यक्षदैवतम्} %॥३४॥

\twolineshloka
{ममाक्षिणी तयोरत्र सङ्क्रान्ती भवतो यतः}
{कर्कसंज्ञा सिंहसंज्ञा माहात्म्यं किमतः परम्} %॥३५॥

\twolineshloka
{प्रातः स्नानं मासमात्रमत्र यः कुरुते नरः}
{द्वादशस्वपि मासेषु प्रातः स्नानफलं लभेत्} %॥३६॥

\twolineshloka
{न करोति नभोमासि प्रातः स्नानं यदा नरः}
{द्वादशस्वपि मासेषु कृतं निष्फलतामियात्} %॥३७॥

\twolineshloka
{महादेव दयासिन्धो श्रावणे मास्युषस्यहम्}
{प्रातः स्नानं करिष्यामि निर्विघ्नं कुरु मे प्रभो} %॥३८॥

\twolineshloka
{स्नात्वा शिवं समभ्यर्च्य नभोमाहात्म्यसत्कथाम्}
{शृणुयात्प्रत्यहं भक्त्या एवं मासं नयेत्सुधीः} %॥३९॥

\twolineshloka
{अन्यत्र मासः कृष्णादिरत्र शुक्लादिरिष्यते}
{नभोमासकथायास्तु माहात्म्यं केन वर्ण्यते} %॥४०॥

\twolineshloka
{सप्तधापि च या वन्ध्या सा पुत्रं लभते शुभम्}
{विद्यार्थी लभते विद्यां बलार्थी लभते बलम्} %॥४१॥

\twolineshloka
{रोगी चारोग्यमाप्नोति बद्धो मुच्येत बन्धनात्}
{धनं धनार्थी लभते धर्मे चैव रतिर्भवेत्} %॥४२॥

\twolineshloka
{भार्यार्थी लभते भार्यां किं बहूक्तेन मानद}
{यद्यत्कामयते तत्तत्प्राप्नोत्यत्र न संशयः} %॥४३॥

\twolineshloka
{अन्ते मम पुरं प्राप्य मोदते मम सन्निधौ}
{पूजयेद् वाचकं सम्यग्वासोऽलङ्करणादिभिः} %॥४४॥

\twolineshloka
{वाचकस्तोषितो येन तेनाहं तोषितः शिवः}
{श्रुत्वा श्रावणमाहात्म्यं वाचकं यो न पूजयेत्} %॥४५॥

\twolineshloka
{छिनत्ति रविजस्तस्य कर्णौ स बधिरो भवेत्}
{तस्माच्छक्त्या प्रकुर्वीत वाचकस्य सुपूजनम्} %॥४६॥

\twolineshloka
{इदं श्रावणमाहात्म्यं यः पठेच्छृणुयादपि}
{श्रावयेद्वापि सद्भक्त्या तस्य पुण्यमनन्तकम्} %॥४७॥


{॥इति श्रीस्कन्दपुराणे ईश्वरसनत्कुमारसंवादे श्रावणमासमाहात्म्ये नदीरजोदोषसिंहगोप्रसव-सिंहकर्कटश्रावणस्तुतिवाचकपूजाकथनं नाम सप्तविंशोऽध्यायः॥२७॥}




\sect{अष्टाविंशोऽध्यायः}

\uvacha{ईश्वर उवाच}
\twolineshloka
{अतः परं प्रवक्ष्यामि अगस्त्यार्घ्यविधिं परम्}
{येन चीर्णेन वैधात्र सर्वान्कामानवाप्नुयात्} %॥१॥

\twolineshloka
{कालस्य नियमो ज्ञेय अगस्त्यस्योदयात्पुरा}
{समरात्राद्भवेद्यावदुदयः सप्तरात्रकम्} %॥२॥

\twolineshloka
{दद्यादर्घ्यं प्रत्यहं च तद्विधिं ते वदाम्यहम्}
{प्रातः शुक्लतिलैः स्नात्वा शुक्लमाल्याम्बरो गृही} %॥३॥

\twolineshloka
{स्थापयेदव्रणं कुम्भं सुवर्णादिविनिर्मितम्}
{पञ्चरत्नसमायुक्तं घृतपात्रेण संयुतम्} %॥४॥

\twolineshloka
{कुम्भोद्भवस्य प्रतिमां तत्र पात्रे निधापयेत्}
{अङ्गुष्ठमात्रं पुरुषं सौवर्णं च चतुर्भुजम्} %॥६॥

\twolineshloka
{नानाभक्ष्यफलैर्युक्तं माल्यवस्त्रविभूषितम्}
{ताम्रेण पूर्णपात्रेण उपरिस्थेन भूषितम्} %॥५॥

\twolineshloka
{पीनात्यायतदोर्दण्डं दक्षिणाभिमुखं मुनिम्}
{सुशोभनं प्रशान्तं च जटामण्डलधारिणम्} %॥७॥

\twolineshloka
{कमण्डलुकरं शिष्यैर्बहुभिः परिवारितम्}
{तथा दर्भाक्षतधरं लोपामुद्रासमन्वितम्} %॥८॥

\twolineshloka
{आवाहयेत्पूजयेच्च गन्धपुष्पादिभिस्तथा}
{उपचारैः षोडशभिर्नैवेद्यैर्बहुविस्तरैः} %॥९॥

\twolineshloka
{दध्योदनबलिं दद्याद्भक्तियुक्तेन चेतसा}
{ततश्चार्घः प्रदातव्यस्तं चैव विधिवच्छृणु} %॥१०॥

\twolineshloka
{सौवर्णे वाथ रौप्ये वा ताम्रे वेणुमयेऽथ वा}
{पात्रे नारङ्गखर्जूरीनारिकेलफलानि च} %॥११॥

\twolineshloka
{कूष्माण्डकारवल्लीनि कदलीदाडिमानि च}
{वृन्ताकबीजपूराणि अक्षोटाः पिस्तकास्तथा} %॥१२॥

\twolineshloka
{नीलोत्पलानि पद्मानि कुशदूर्वाङ्कुरास्तथा}
{अन्यान्यपि च साध्यानि फलानि कुसुमानि च} %॥१३॥

\twolineshloka
{नानाप्रकारभक्ष्याणि सप्तधान्यानि चैव हि}
{सप्ताङ्कुराः पल्लवाश्च पञ्च वस्त्राणि चैव हि} %॥१४॥

\twolineshloka
{एतान्पदार्थान्संस्थाप्य सम्यक्प्रपूजयेत्}
{जानुभ्यामवनिं गत्वा तत्पात्रं नम्रमूर्धनि} %॥१५॥

\twolineshloka
{धृत्वावाचिमुखो भूत्वा ध्यायेत्कुम्भोद्भवं मुनिम्}
{दद्यादर्घ्यं प्रयत्नेन श्रद्धाभक्तिपुरःसरम्} %॥१६॥

\twolineshloka
{काशपुष्पप्रतीकाश वह्निमारुतसम्भव}
{मित्रावरुणयोः पुत्र कुम्भयोने नमोऽस्तु ते} %॥१७॥

\twolineshloka
{विन्ध्यवृद्धिक्षयकर मेघतोयविषापह}
{रत्नवल्लभ देवर्षे लङ्कावास नमोऽस्तु ते} %॥१८॥

\twolineshloka
{आतापी भक्षितो येन वातापी च महाबलः}
{लोपामुद्रापतिः श्रीमान्योऽसौ तस्मै नमो नमः} %॥१९॥

\twolineshloka
{येनोदितेन पापानि विलयं यान्ति चाधयः}
{व्याधयस्त्रिविधास्तापास्तस्मै नित्यं नमो नमः} %॥२०॥

\twolineshloka
{यादः पूर्णः सरिन्नाथो येन वै शोषितः पुरा}
{सपुत्राय सशिष्याय सपत्नीकाय वै नमः} %॥२१॥

\twolineshloka
{अगस्त्यस्यायमर्घो वै द्विजातिर्वेदमन्त्रतः}
{शूद्रः पौराणमन्त्रेण दत्त्वार्घ्यं प्रणमेत्सुधीः} %॥२२॥

\twolineshloka
{राजपुत्रि महाभागे ऋषिपनि वरानने}
{लोपामुद्रे नमस्तुभ्यमर्घ्यं मे प्रतिगृह्यताम्} %॥२३॥

\twolineshloka
{ततो होमं प्रकुर्वीत अर्घ्यमन्त्रेण मन्त्रवित्}
{आज्येनाष्टसहस्रं वा शतमष्टोत्तरं तु वा} %॥२४॥

\twolineshloka
{कृत्वैवं च ततोऽगस्त्यं प्रणिपत्य विसर्जयेत्}
{अचिन्त्यचरितागस्त्य यथागस्त्यः प्रपूजितः} %॥२५॥

\twolineshloka
{ऐहिकामुष्मिकीं कृत्वा कार्यसिद्धिं व्रजस्व भोः}
{विसर्जयित्वागस्त्यं तं विप्राय प्रतिपादयेत्} %॥२६॥

\twolineshloka
{वेदवेदाङ्गविदुषे दरिद्राय कुटुम्बिने}
{अगस्त्यो द्विजरूपेण प्रतिगृह्णातु सत्कृतः} %॥२७॥

\twolineshloka
{अगस्त्यः प्रतिगृह्णाति ह्यगस्त्यो वै ददाति च}
{उभयोस्तारकोऽगस्त्यो ह्यगस्त्याय नमो नमः} %॥२८॥

\twolineshloka
{मन्त्रद्वयेन दद्यात्तु ब्राह्मणस्तु जपेन्मनुम्}
{वैदिकं पूर्वविहितं पौराणं शूद्र एव तु} %॥२९॥

\twolineshloka
{श्वेतां धेनुं ततो दद्याद्धेमशृङ्गीं पयस्विनीम्}
{सहवत्सां रौप्यखुरां ताम्रपृष्ठीं सुशोभनाम्} %॥३०॥

\twolineshloka
{कांस्यदोहनिकायुक्तां घण्टावस्त्रसमन्विताम्}
{एवं सप्तदिनं दत्त्वा अर्घ्यं प्रागुदयान्मुनेः} %॥३१॥

\twolineshloka
{सप्तमे दिवसे धेनुं प्रदद्याच्च सदक्षिणाम्}
{एवं कृत्वा सप्तवर्षमकामश्चेन्न जन्मभाक्} %॥३२॥

\twolineshloka
{सकामश्चक्रवर्तित्वं रूपारोग्यसमन्वितः}
{ब्राह्मणः स्याच्चतुर्वेदसर्वशास्त्रविशारदः} %॥३३॥

\twolineshloka
{क्षत्रियः पृथिवीं सर्वां प्राप्नोत्यर्णवमेखलाम्}
{वैश्यश्चेद्धान्यनिष्पत्तिं गोधनं चापि विन्दति} %॥३४॥

\twolineshloka
{शूद्राणां धनमारोग्यं सत्यं चैवाधिकं भवेत्}
{स्त्रीणां पुत्राः प्रजायन्ते सौभाग्यं गृहमृद्धिमत्} %॥३५॥

\twolineshloka
{विधवानां महापुण्यं वर्धते विधिनन्दन}
{कन्या भर्तारमाप्नोति व्याधेर्मुच्येत दुःखितः} %॥३६॥

\twolineshloka
{येषु देशेष्वगस्त्यस्य पूजनं क्रियते नरैः}
{तेषु देशेषु पर्जन्यः कामवर्षी प्रजायते} %॥३७॥

\twolineshloka
{ईतयः प्रशमं यान्ति नश्यन्ति व्याधयस्तथा}
{पठन्ति ये त्वगस्त्यस्य ह्यर्घ्यं शृण्वन्ति केचन} %॥३८॥

\twolineshloka
{ते सर्वे पापनिर्मुक्ताश्चिरं स्थित्वा महीतले}
{हंसयुक्तविमानेन स्वर्गे यान्ति नरोत्तमाः} %॥३९॥

\onelineshloka
{यावज्जीवं करिष्यन्ति निष्कामं मुक्तिभागिनः} %॥४०॥


{॥इति श्रीस्कन्दपुराणे ईश्वरसनत्कुमारसंवादे श्रावणमासमाहात्म्ये अगस्त्यार्घ्यविधिर्नामाष्टाविंशोऽध्यायः॥२८॥}




\sect{एकोनत्रिंशोऽध्यायः}

\uvacha{ईश्वर उवाच}
\twolineshloka
{सनत्कुमार वक्ष्यामि उक्तानां व्रतकर्मणाम्}
{काले कदा तु किं कार्यं तच्छृणुष्व महामुने} %॥१॥

\twolineshloka
{का तिथिः श्रावणे मासि किङ्कालव्यापिनी भवेत्}
{ग्राह्या प्रधानं किं तत्र पूजाजागरणादिकम्} %॥२॥

\twolineshloka
{तत्तत्कथनकाले तु केषाञ्चित्काल ईरितः}
{नक्तव्रतस्य कालस्तु उक्तस्तद्व्रतकर्मणि} %॥३॥

\twolineshloka
{प्रधानं रात्रिभुक्तिस्तु भोजनाभावयुग्दिवा}
{उद्यापनं तु सर्वेषां तत्तद्व्रततिथौ भवेत्} %॥४॥

\twolineshloka
{असम्भवे तु पञ्चाङ्गशुद्धे स्यादधिवासनम्}
{द्वितीयदिवसे कुर्याद्धोमादिविधिमादृतः} %॥५॥

\twolineshloka
{धारणा पारणा चैव ह्रासवृद्धी न कारणम्}
{नभः शुक्लप्रतिपदि सङ्कल्प्योपोषणं चरेत्} %॥६॥

\twolineshloka
{द्वितीयदिवसे भुक्तिस्ततोऽन्यस्मिन्नुपोषणम्}
{एवं क्रमेण कुर्वीत हविष्याशी तु पारणे} %॥७॥

\twolineshloka
{एकादशीपारणाहे उपवासत्रयं तथा}
{रविवारव्रतार्चायाः कालः स्यात्प्रातरेव हि} %॥८॥

\twolineshloka
{सोमवारे प्रधानः स्यात्सायङ्कालः प्रकीर्तितः}
{भौमे बुधे गुरौ मुख्यः प्रातःकालश्च पूजने} %॥९॥

\twolineshloka
{शुक्रवारे पूजनं स्यात्कल्पे रात्रौ च जागरः}
{नृसिंहपूजने मन्दे सायङ्कालश्च पूजनम्} %॥१०॥

\twolineshloka
{शनिव्रते शनैर्दाने मध्याह्ने मुख्य इष्यते}
{हनूमतोऽपि मध्याह्नः प्रातरश्वत्थपूजनम्} %॥११॥

\twolineshloka
{रोटकाख्ये व्रते प्रतिपत्सोमसंयुता}
{त्रिमुहूर्तोत्तरा सा स्यादन्यथा पूर्वयोगिनी} %॥१२॥

\twolineshloka
{वत्स औदुम्बरी द्वितीया तु सायाह्नव्यापिनी मता}
{तृतीयासंयुता ग्राह्या द्वयोश्चेत्पूर्ववेधिता} %॥१३॥

\twolineshloka
{तृतीया स्वर्णगौर्याख्या सा चतुर्थीयुता भवेत्}
{चतुर्थी गणनाथस्य मातृविद्धा प्रशस्यते} %॥१४॥

\twolineshloka
{नागानां पूजने शस्ता षष्ठीयुक्ता च पञ्चमी}
{सूपौदनव्रते षष्ठी सायाने सप्तमीयुता} %॥१५॥

\twolineshloka
{मध्याह्नव्यापिनी ग्राह्या सप्तमी शीतलाव्रते}
{पवित्रारोपणेऽष्टम्यां देव्या रात्रियुता तिथिः} %॥१६॥

\twolineshloka
{कुमारीनवमी नक्तव्यापिनी तु प्रशस्यते}
{आशासंज्ञा तु दशमी सा नक्तव्यापिनी भवेत्} %॥१७॥

\twolineshloka
{त्याज्या विद्धैकादशी तु तत्र वेधं मुने शृणु}
{अरुणोदयवेधस्तु दशम्या वैष्णवान्प्रति} %॥१८॥

\twolineshloka
{आदित्योदयवेधस्तु स्मार्तानां निन्द्य एव सः}
{अरुणोदयकालस्तु यामार्धं चरमं निशः} %॥१९॥

\twolineshloka
{एवं रीत्या यस्य भवेद् द्वादशी सा पवित्रके}
{त्रयोदशी त्वनङ्गस्य व्रते स्याद् रात्रियोगिनी} %॥२०॥

\twolineshloka
{द्वितीययामे तत्रापि प्रशस्ततरा भवेत्}
{पवित्रारोपणे शम्भो रात्रिगा स्याच्चतुर्दशी} %॥२१॥

\twolineshloka
{अतिप्रशस्ता तत्रापि निशीथव्यापिनी तु या}
{उपाकर्मणि चोत्सर्गे पूर्णिमा श्रवणं च भम्} %॥२२॥

\twolineshloka
{त्रिमुहूर्तं द्वितीयेऽह्नि तदा ग्राह्यं परं दिनम्}
{नोचेदनुष्ठितः पूर्वं तैत्तिराणां च बह्वृचाम्} %॥२३॥

\twolineshloka
{तैत्तिराणां च यजुषां मुहूर्तत्रयगामपि}
{उत्तरस्मिन्पूर्वमेव दिनं स्यात्कर्मणि द्वयोः} %॥२४॥

\twolineshloka
{मुहूर्तानन्तरं पूर्वदिने चेत्सङ्गतिर्भवेत्}
{पूर्णिमाश्रवण च मुहूर्तद्वितयात्पुरा} %॥२५॥

\twolineshloka
{उत्तरस्मिन्समाप्तं चेत्तदा पूर्वदिनं भवेत्}
{हस्तभे त्वपराह्णे स्याद् ग्राह्यं तत्सामवेदिभिः} %॥२६॥

\twolineshloka
{दिनद्वये तदा स्याच्चेत्पूर्वमेव दिनं भवेत्}
{उपाकर्मप्रयोगान्ते कालो दीपस्य संसदः} %॥२७॥

\twolineshloka
{श्रवणाकर्मणि प्रोक्तः कालः सर्वबलौ तथा}
{पर्वणोऽह्नि भवेद् रात्रौ स्वस्वगृह्यानुसारतः} %॥२८॥

\twolineshloka
{स्यादस्तमययोगिनी}
{हयग्रीवोत्सवे पूर्णा मध्याह्नव्यापिनी भवेत्} %॥२९॥

\twolineshloka
{पौर्णिमात्र प्रशस्ता वा अपराह्णव्यापिनी स्याद्रक्षाबन्धनकर्मणि}
{चन्द्रोदयव्यापिनी च स्यात्सङ्कष्टचतुर्थिका} %॥३०॥

\twolineshloka
{उभयत्र यदा सा स्यान्न स्याद्वा पूर्वगा भवेत्}
{चतुर्थी च तृतीयायां महापुण्यफलप्रदा} %॥३१॥

\twolineshloka
{कर्तव्या व्रतिभिर्वत्स गणनाथसुतोषिणी}
{गणेशगौरीबहुलाव्यतिरिक्ताः प्रकीर्तिताः} %॥३२॥

\twolineshloka
{चतुर्थ्यः पञ्चमीविद्धा देवतान्तरपूजने}
{निशीथव्यापिनी ग्राह्या कृष्णजन्माष्टमी तिथिः} %॥३३॥

\twolineshloka
{षट्प्रकारा तु सर्वत्र निर्णये तिथिरिष्यते}
{पूर्णव्याप्तिर्द्वयोरनोरव्याप्तिरपि केवला} %॥३४॥

\twolineshloka
{अंशतश्च समा व्याप्तिरंशतो विषमा तथा}
{सम्पूर्णव्याप्तिरेकत्र अंशतश्च परेऽहनि} %॥३५॥

\twolineshloka
{अंशतो व्याप्तिरेकत्र अव्याप्तिरपरत्र च}
{पक्षत्रये तु सन्देहो यथा नास्ति तथा शृणु} %॥३६॥

\twolineshloka
{अंशतो विषमव्याप्तावधिका व्याप्तिरुत्तमा}
{एकत्र पूर्णा चान्यत्र सापूर्णा चोच्यते तिथिः} %॥३७॥

\twolineshloka
{अव्याप्तिरंशतो व्याप्तिस्तत्रांशव्याप्तिरुत्तमा}
{अंशव्याप्तिर्यदा पूर्णा अंशतश्च समा यदा} %॥३८॥

\twolineshloka
{संशयस्तत्र भवति तस्य स्यान्निर्णयो भिदा}
{क्वचिद्भवेद्युग्मवाक्याद्वारनक्षत्रयोगतः } %॥३९॥

\twolineshloka
{प्रधानद्वययोगेन पारणायोगतः क्वचित्}
{जन्माष्टम्यां तु सन्देहे त्रिपक्षे तु परा भवेत्} %॥४०॥

\twolineshloka
{अष्टम्यन्ते पारणं स्याद्यदि यामत्रयात्परा}
{समाप्येत तदूर्ध्वं चेदष्टम्युषसि पारणा} %॥४१॥

\twolineshloka
{व्रते पिठोरसंज्ञेऽमा मध्यानव्यापिनी शुभा}
{वृषभाणां पूजने तु अमा सायन्तनी भवेत्} %॥४२॥

\twolineshloka
{दर्भाणां सञ्चये चैव सङ्गवः काल ईरितः}
{त्रिंशत्पुण्याः पूर्वनाड्यः कर्कसङ्क्रमणे रवेः} %॥४३॥

\twolineshloka
{पुण्याः षोडश नाड्यस्तु सिंहे पूर्वाः परा अपि}
{केचिदिच्छन्ति मुनयः पूर्वा एव तु षोडश} %॥४४॥

\twolineshloka
{अगस्त्यार्घ्यस्य कालस्तु व्रत एव प्रकीर्तितः}
{अयं ते कथितो वत्स कर्मणां कालनिर्णयः} %॥४५॥

\twolineshloka
{य इदं शृणुतेऽध्यायं यश्चापि परिकीर्तयेत्}
{नभोमासि कृतानां स व्रतानां लभते फलम्} %॥४६॥


{॥इति श्रीस्कन्दपुराणे ईश्वरसनत्कुमारसंवादे श्रावणमासमाहात्म्ये व्रतनिर्णयकालनिर्णयकथनं नाम एकोनत्रिंशोऽध्यायः॥२९॥}




\sect{त्रिंशोऽध्यायः}

\uvacha{ईश्वर उवाच}
\twolineshloka
{कियत्कियत्ते कथितं माहात्म्यं श्रावणस्य हि}
{सर्वं वर्णयितुं शक्यं नालं वर्षशतैरपि} %॥१॥

\twolineshloka
{प्रियेयं मम कल्याणी हुत्वा दक्षाध्वरे तनुम्}
{हिमाचलसुता जाता तेनेयं योजिता मया} %॥२॥

\twolineshloka
{सेवने श्रावणे मासि तेन मे प्रियकृन्नभाः}
{नातिशीतो नाति चोष्णः श्रावणे मासि भूपतिः} %॥३॥

\twolineshloka
{उद्धूलयित्वा स्वतनुं सर्वां श्रौतेन भस्मना}
{श्वेतेनाथ जलार्द्रेण त्रिपुण्ड्रान्द्वादशांश्चरेत्} %॥४॥

\twolineshloka
{भाले वक्षसि नाभौ च बाह्वोः कूर्परयोस्तथा}
{मणिबन्धद्वये चैव कण्ठे मूर्धनि पृष्ठके} %॥५॥

\twolineshloka
{मानस्तोकेति मन्त्रेण सद्योजातादिमन्त्रतः}
{षडक्षरेण मन्त्रेण भस्मना शोभयेत्तनुम्} %॥६॥

\twolineshloka
{धारयेच्चैव रुद्राक्षानष्टाधिकशतं तनौ}
{द्वात्रिंशद्धारयेत्कण्ठे मूर्ध्नि द्वाविंशतिस्तथा} %॥७॥

\twolineshloka
{कर्णद्वये द्वादशैव चतुर्विंशत्करद्वये}
{अष्टाष्ट भुजयोर्भाले एकमेकं शिखाग्रगम्} %॥८॥

\twolineshloka
{एवं कृत्वा तु मामर्च्य जपेत्पञ्चाक्षरं मनुम्}
{श्रावणे मासि विप्रेन्द्र सोऽहमेव न संशयः} %॥९॥

\twolineshloka
{ज्ञात्वेमं मत्प्रियं मासं पूजयेत्केशवं च माम्}
{कृष्णाष्टमी च तत्रापि मम प्रियतरा तिथिः} %॥१०॥

\twolineshloka
{देवक्या जठरात्तस्मिन्दिने प्रादुरभूद्धरिः}
{एतत्ते कथितं लेशात्किमन्यच्छ्रोतुमिच्छसि} %॥११॥


\uvacha{सनत्कुमार उवाच}
\twolineshloka
{यद्यत्कृत्यं पार्वतीश नभोमासि त्वयेरितम्}
{आनन्दाब्धौ निमग्नत्वाद् बहुत्वाच्चावधारणा} %॥१२॥

\twolineshloka
{न स्थिता क्रमशो नाथ ब्रूहि सर्वं यथा तथा}
{श्रुत्वा चाव्यवधानेन धारयिष्यामि भक्तितः} %॥१३॥


\uvacha{ईश्वर उवाच}
\twolineshloka
{शृणुष्वावहितो भूत्वा अनुक्रमणिकां शुभाम्}
{आदौ प्रश्नः शौनकस्य ततः सूतस्य चोत्तरम्} %॥१४॥

\twolineshloka
{श्रोतुर्गुणास्तव प्रश्ना निरुक्तिः श्रावणस्य च}
{तस्य स्तुतिः पुनः प्रश्नस्तव विस्तरशो मुने} %॥१५॥

\twolineshloka
{मम स्तुतिस्त्वत्कृता च नामनिर्वचनादिना}
{भूयो ममोत्तरं तत्र उद्देश्यः क्रमतोऽखिलः} %॥१६॥

\twolineshloka
{विशेषतस्तव प्रश्नास्ततो नक्तव्रते विधिः}
{रुद्राभिषेककथनं लक्षपूजाविधिस्ततः} %॥१७॥

\twolineshloka
{दीपदानं परित्यागः कस्यचित्प्रियवस्तुनः}
{फलं रुद्राभिषेकेण तथा पञ्चामृतेन च} %॥१८॥

\twolineshloka
{फलं भूशयनस्यापि तथा मौनव्रतस्य च}
{धारणा पारणा चैव ततो मासोपवासने} %॥१९॥

\twolineshloka
{सोमाख्याने ततो लक्षरुद्रवर्तिविधिः स्मृतः}
{कोटिलिङ्गविधानं च व्रतं चानौदनाभिधम्} %॥२०॥

\twolineshloka
{हविष्याशनमप्यत्र पत्रावल्यां च भोजनम्}
{शाकत्यागो भूशयनं प्रातः स्नानं दमः शमः} %॥२१॥

\twolineshloka
{स्फटिकादिषु लिङ्गेषु अर्चा जपफलं ततः}
{प्रदक्षिणा नमस्कारान्वेदपारायणं तथा} %॥२२॥

\twolineshloka
{विधिः पुरुषसूक्तस्य ग्रहयज्ञविधिस्ततः}
{रविचन्द्रकुजानां च क्रमशो व्रतविस्तरः} %॥२३॥

\twolineshloka
{बुधगुर्वोव्रतं पश्चाच्छुक्रे जीवन्तिकाव्रतम्}
{शनौ नृसिंहस्य शनेरनिलाश्वत्थयोस्तथा} %॥२४॥

\twolineshloka
{रोटकव्रतमाहात्म्यं तत औदुम्बरव्रतम्}
{स्वर्णगौरीव्रतं पश्चाद् दूर्वागणपतिव्रतम्} %॥२५॥

\twolineshloka
{नागव्रतं च पञ्चम्यां षष्ठ्यां सूपौदनव्रतम्}
{शीतलाख्यं व्रतं देव्याः पवित्रारोपणं ततः} %॥२६॥

\twolineshloka
{दुर्गाकुमारीपूजा च आशाव्रतमतः परम्}
{उभयैकादशी पश्चात्पवित्रारोपणं हरेः} %॥२७॥

\twolineshloka
{अनङ्गस्य त्रयोदश्यां ततः शम्भोः पवित्रकम्}
{उपाकर्मोत्सर्जने च श्रवणाकर्म चैव हि} %॥२८॥

\twolineshloka
{ततः सर्पबलिर्वाजिग्रीवजन्ममहोत्सवः}
{सभादीपस्तथा रक्षाबन्धः सङ्कटनाशनम्} %॥२९॥

\twolineshloka
{व्रतं ततः कृष्णजन्माष्टमीव्रतकथानकम्}
{व्रतं पिठोरसंज्ञं तु पोलासंज्ञं वृषव्रतम्} %॥३०॥

\twolineshloka
{दर्भाणां सङ्ग्रहश्चैव नदीनां सरजस्कता}
{सिंहे गोप्रसवे शान्तिः कर्कसिंहनभेषु च} %॥३१॥

\twolineshloka
{दानानि स्नानमाहात्म्यं माहात्म्यश्रवणं तथा}
{ततो वाचकपूजा च अगस्त्यार्घ्यं ततः परम्} %॥३२॥

\twolineshloka
{कर्मणां च व्रतानां च कालनिर्णय ईरितः}
{एतन्मासि कृतानां स व्रतानां फलभाग्भवेत्} %॥३३॥

\onelineshloka
{सनत्कुमार हृदये धारयस्व क्रमं शुभम्} %॥३४॥

\twolineshloka
{य इमं शृणुतेऽध्यायं माहात्म्यं श्रावणस्य यत्}
{तत्फलं समवाप्नोति व्रतानां चैव यत्फलम्} %॥३५॥

\twolineshloka
{किं बहूक्तेन विप्रर्षे श्रावणे विहितं तु यत्}
{तस्य चैकस्य कर्तापि मम प्रियतरो भवेत्} %॥३६॥


\uvacha{सूत उवाच}
\twolineshloka
{सनत्कुमारः पीत्वेदं शिववाक्यामृतं परम्}
{श्रुतिद्वारा चाप मोदं कृतकृत्यो बभूव ह} %॥३७॥

\twolineshloka
{नभोमासं स्तुवञ्छम्भुं स्मरन्स हृदये शिवम्}
{शङ्करेणाभ्यनुज्ञातो ययौ देवर्षिसत्तमः} %॥३८॥

\twolineshloka
{इदं रहस्यं परमं नाख्येयं यस्य कस्यचित्}
{भवतो योग्यतां दृष्ट्वा मयैतत्कथितं प्रभो} %॥३९॥


{॥इति श्रीस्कन्दपुराणे ईश्वरसनत्कुमारसंवादे श्रावणमासमाहात्म्ये अनुक्रमणिकाकथनं नाम त्रिंशोऽध्यायः॥३०॥}

{॥इति श्रावणमासमाहात्म्यं सम्पूर्णम्॥}
